% ############################################### PREAMBLE AND CONTROL PANEL #

% ########################################################### DOCUMENT CLASS #
\documentclass[  %   draft
             % , demo  % black boxes instead of figures
              , 8pt, twoside, openright]{extbook}

% ########################################################## AND CONTROL PANEL #

% true - hmarginration 1:1, urlcolor - darkgreen, version - for searching
\newif\ifscreen 	   \screentrue
% true - no apendix, no author of grammer, no backmatter ISBN, 
\newif\ifPDF               \PDFtrue
% true - rotate pages of morph verb tables
\newif\ifPDFrotate         \PDFrotatetrue
% true - adds white searchable value !
\newif\ifPDFsearchvalue    \PDFsearchvaluetrue
\newif\ifmakecovers        \makecoverstrue
\newif\ifmakefrontmatter   \makefrontmattertrue
\newif\ifinputintroduction \inputintroductiontrue
\newif\ifinputletters      \inputletterstrue
\newif\ifinputphon         \inputphontrue
\newif\ifinputmorpho       \inputmorphotrue
\newif\ifmakebackmatter    \makebackmattertrue
\newif\ifshowcrop           \showcroptrue
% corrections
\newif\ifC                  \Ctrue

\newif\ifPRINTHOUSE 	   \PRINTHOUSEfalse
\newif\ifPDFSEARCH  	   \PDFSEARCHfalse
\newif\ifPDFPRINT 	   \PDFPRINTfalse
\newif\ifPDFtest	   \PDFtesttrue

%versions
%PRINTHOUSE
\ifPRINTHOUSE
\makecoversfalse
\screenfalse
\PDFfalse
\PDFrotatefalse
\PDFsearchvaluefalse
\showcroptrue
\Ctrue
\fi

%PDFSEARCH
\ifPDFSEARCH
\screentrue
\PDFtrue
\PDFrotatetrue
\PDFsearchvaluetrue
\showcropfalse
\Ctrue
\fi

%PDFPRINT
\ifPDFPRINT
\screenfalse
\PDFtrue
\PDFrotatefalse
\PDFsearchvaluefalse
\showcropfalse
\Ctrue
\fi

%PDFtest
\ifPDFtest
\makecoversfalse
\makefrontmatterfalse
\inputintroductionfalse
\inputletterstrue
\inputphonfalse
\inputmorphofalse
\makebackmatterfalse
\screentrue
\PDFtrue
\PDFrotatefalse
\PDFsearchvaluefalse
\showcropfalse
\Cfalse
\fi

% 1. PDF (for printing)  
% a) searchable value, b) no apendix, c) version= PDF d) no ISBN and no ISBN bar
% e) automatic rotation of landscape pages, f) no author of Icelandic Grammer
% 2. PDF (for viewing offline and searching) + screen (horizontal ratio 1:1) 
% 3. PRINTHOUSE
% a) no searchable value, b) apendix, c) version ISBN d) ISBN bar
% e) no rotation of landscape pages, f)urllinkcolor black
% f) A4 paper dimentions, C5 layout, shows cropping marks

% ########################################################## PAGE GEOMETRIES #

\usepackage%[showframe]
           {geometry}
           

\ifshowcrop
% Basic geometry of document
\geometry
  { headsep    =   \baselineskip
  , textwidth  = 42\baselineskip
  , textheight = 60\baselineskip
  , hmarginratio = \ifscreen 1:1\else 2:3\fi
  , vmarginratio = 2:3
  , bindingoffset = 0cm
  , onecolumn
  , a4paper 
  , layout=c5paper,
  , layouthoffset=\dimexpr(\paperwidth-\csname Gm@layoutwidth\endcsname)/2\relax
   , layoutvoffset=\dimexpr(\paperheight-\csname Gm@layoutheight\endcsname)/2\relax
  , showcrop
  }
\else
\geometry
  { headsep    =   \baselineskip
  , textwidth  = 42\baselineskip
  , textheight = 60\baselineskip
  , hmarginratio = \ifscreen 1:1\else 2:3\fi
  , vmarginratio = 2:3
  , bindingoffset = 0cm
  , onecolumn,
  , c5paper
  }
\fi

\ifshowcrop
% Geometry of the dictionary section
\newcommand\dictionarygeometry{\newgeometry
  { textwidth  = 42\baselineskip
  , textheight = 60\baselineskip
  , hmarginratio = \ifscreen 1:1\else 2:3\fi
  , vmarginratio = 3:2
  , bindingoffset = 0cm
  , twocolumn
  , a4paper 
  , layout=c5paper
  , layouthoffset=\dimexpr(\paperwidth-\csname Gm@layoutwidth\endcsname)/2\relax
   , layoutvoffset=\dimexpr(\paperheight-\csname Gm@layoutheight\endcsname)/2\relax
  , showcrop
  }}
   % Geometry of the cover pages
\newcommand\covergeometry{\newgeometry
  { textwidth  = 42\baselineskip
  , textheight = 60\baselineskip
  , hmarginratio = 1:1
  , vmarginratio = 2:3
  , bindingoffset = 0cm
  , onecolumn
  , layout=c5paper
 , layouthoffset=\dimexpr(\paperwidth-\csname Gm@layoutwidth\endcsname)/2\relax
   , layoutvoffset=\dimexpr(\paperheight-\csname Gm@layoutheight\endcsname)/2\relax
  , showcrop
  }} 
\else
% Geometry of the dictionary section
\newcommand\dictionarygeometry{\newgeometry
  { textwidth  = 42\baselineskip
  , textheight = 60\baselineskip
  , hmarginratio = \ifscreen 1:1\else 2:3\fi
  , vmarginratio = 3:2
  , bindingoffset = 0cm
  , twocolumn,
  , c5paper
  }}
   % Geometry of the cover pages
\newcommand\covergeometry{\newgeometry
  { textwidth  = 42\baselineskip
  , textheight = 60\baselineskip
  , hmarginratio = 1:1
  , vmarginratio = 2:3
  , bindingoffset = 0cm
  , onecolumn
  , c5paper
  }} 
\fi

% ################################################################# LANGUAGES #

\usepackage[icelandic, latin, czech]{babel}
\usepackage{csquotes}


% ################################################################## ENCODING #

\usepackage[utf8]{inputenc}

% Smashing uppercase letters with diacritics
% ------------------ THIS IS EVIL AND SHOULD NEVER BE DONE -------------------
% Don't smash anything if we are on the first run to produce the index
\IfFileExists{figures.idx}{
  % Czech uppercase letters with diacritics:
  \DeclareUnicodeCharacter{00C1}{\strut\smash{\' A}} % Á
  \DeclareUnicodeCharacter{010C}{\strut\smash{\v C}} % Č
  \DeclareUnicodeCharacter{010E}{\strut\smash{\v D}} % Ď
  \DeclareUnicodeCharacter{00C9}{\strut\smash{\' E}} % É
  \DeclareUnicodeCharacter{011A}{\strut\smash{\v E}} % Ě
  \DeclareUnicodeCharacter{00CD}{\strut\smash{\' I}} % Í
  \DeclareUnicodeCharacter{0147}{\strut\smash{\v N}} % Ň
  \DeclareUnicodeCharacter{00D3}{\strut\smash{\' O}} % Ó
  \DeclareUnicodeCharacter{0158}{\strut\smash{\v R}} % Ř
  \DeclareUnicodeCharacter{0160}{\strut\smash{\v S}} % Š
  \DeclareUnicodeCharacter{0164}{\strut\smash{\v T}} % Ť
  \DeclareUnicodeCharacter{00DA}{\strut\smash{\' U}} % Ú
  \DeclareUnicodeCharacter{016E}{\strut\smash{\r U}} % Ů
  \DeclareUnicodeCharacter{00DD}{\strut\smash{\' Y}} % Ý
  \DeclareUnicodeCharacter{017D}{\strut\smash{\v Z}} % Ž
  % Icelandic uppercase letters with diacritics: ÁÉÍÓÚÝ (ÐÞÆ) and
  \DeclareUnicodeCharacter{00D6}{\strut\smash{\" O}} % Ö
}{}
% ------------------ THAT WAS EVIL AND SHOULD NEVER BE DONE ------------------

% Non-breakable space
\DeclareUnicodeCharacter{00A0}{~}

% Non-breakable hyphen
\newcommand*{\nobreakhyphen}{\mbox{-}}
\DeclareUnicodeCharacter{2011}{\nobreakhyphen}


% ############################################################ FONTS, SYMBOLS #

\usepackage[T1]{fontenc}

\usepackage{tgpagella}
\usepackage[scaled=0.88]{helvet}      % relative scale of the two fonts

\def\phvfamily{\fontfamily{phv}\selectfont}

%\usepackage{stmaryrd}
% Let's keep this minimal instead of using the package:
\DeclareSymbolFont{stmry}{U}{stmry}{m}{n}
%\SetSymbolFont{stmry}{bold}{U}{stmry}{b}{n}
\DeclareMathSymbol\shortrightarrow\mathrel{stmry}{"01}
\DeclareMathSymbol\shortuparrow\mathrel{stmry}{"02}

% Methinks these are not necessary
%\usepackage{dingbat}
%\usepackage{manfnt}
%\usepackage{latexsym}

\usepackage{pifont}




% #################################################################### LAYOUT #

% \normalsize should be {8pt}{9.6pt}
\def\HUGE {\fontsize{23.887872pt}{3\baselineskip}\selectfont}
\def\Huge {\fontsize{19.906560pt}{3\baselineskip}\selectfont}
\def\huge {\fontsize{16.588800pt}{3\baselineskip}\selectfont}
\def\LARGE{\fontsize{13.824000pt}{2\baselineskip}\selectfont}
\def\Large{\fontsize{11.520000pt}{2\baselineskip}\selectfont}
\def\large{\fontsize{ 9.600000pt}{2\baselineskip}\selectfont}

% Two columns layout ruler
\setlength\columnsep    {2\baselineskip}
\setlength\columnseprule{0.4pt}

% Necessary for baseline alignment
\topskip=\baselineskip
\raggedbottom
\setlength\parskip{0pt} % it's better to avoid glue

% Temporarily suppress warnings
\hbadness=5000
\vbadness=5000

\setlength\emergencystretch{17pt}

% Allow smaller emergencystretch in several cases
\newenvironment{xtolerant}[2]{%
  \par
  \ifx\empty#1\empty\else\tolerance=#1\relax\fi
  \ifx\empty#2\empty\else\emergencystretch=#2\relax\fi
}{%
  \par
}

% ######################################################### ASSORTED PACKAGES #

\usepackage{tikz}
\usetikzlibrary{calc}

\usepackage{tipa}                      % IPA phonetics

\usepackage{fourier-orns}              % ornaments
\usepackage{amsmath}                   % non-breakable dash
\usepackage{hyphenat}                  % no hyphen in abbreviations

\usepackage{hanging}
\usepackage{fix2col}
\usepackage{fixltx2e}                  % subscript
\usepackage{pagecolor}
\usepackage{afterpage}                 % Used in coloring trick for covers
\usepackage{pdfpages}

\usepackage{multirow}
\usepackage{multicol}
\usepackage{rotating}

\usepackage{etoolbox} % http://ctan.org/pkg/etoolbox
%\usepackage{newunicodechar} 		% new definition of middle dot
\usepackage[ISBN=978-80-260-2385-2,SC0]{ean13isbn}

\usepackage{pgffor}

\usepackage{newunicodechar}

\newunicodechar{·}{\chejnikcdot}
\makeatletter
\newcommand{\chejnikcdot}{%
  \ifnum\lastpenalty=10042
    \kern-0.06667em\relax % dots should always be kerned
  \fi
  \textperiodcentered
  \@ifnextchar|{\kern\chejnikkern\relax}{\penalty10042 }%
}

\newcommand{\closeupdotbar}{%
  \renewcommand\chejnikkern{-0.06667em}%
}
\newcommand\chejnikkern{0pt} % default

% Middle dot + | little bit closer, nicer
%\newunicodechar{·}{\advancedcdot}
%\makeatletter
%\newcommand*\advancedcdot{\textperiodcentered
%  \@ifnextchar|{\hspace{-0.03667em}}{\checkifcdotnext}}
%\newcommand*\checkifcdotnext[2]
%  {\def\tmpa{·}\def\tmpb{#1#2}\ifx\tmpa\tmpb\hspace{-0.06667em}\fi#1#2}
%\makeatother

% ############################################################ BASELINE GRID #

\usepackage{eso-pic}
\usepackage{tikzpagenodes}

% Command to draw a baseline grid
\newcommand\drawbaselinegrid{%
  \begin{tikzpicture}[overlay,remember picture]
    \draw [red!30!white, ultra thin, dashed]
      (0,0) grid [ ystep  = \baselineskip, xstep = \textwidth
                 , shift  = (current page text area.north west)
                 , yshift = -\dp\strutbox
                 ] ++(\textwidth,-\textheight);
    \draw [red!30!white, thin]
      (0,0) grid [ step = \baselineskip, xstep = \textwidth
                 , shift=(current page text area.north)
                 ] ++(0.5\textwidth,-\textheight)
      (0,0) grid [ step = \baselineskip, xstep = \textwidth
                 , shift=(current page text area.north)
                 ] ++(-0.5\textwidth,-\textheight);
    \draw[black!10!white]
      (current page text area.north west)
        rectangle (current page text area.south east);
  \end{tikzpicture}}

% Draw a baseline grid on every page
% \AddToShipoutPicture{\drawbaselinegrid}


% ################################################################## METADATA #

\title{%
  \textbf{Islandsko-český studijní slovník}%
  \thanks{Tato kniha byla vytvořena v \LaTeX{}u pod Ubuntu 14.04.%
    Poděkování patří všem autorům, kteří publikují pod svobodnými licencemi.}}
\author{Aleš Chejn, Ján Zaťko, Jón Gíslason}
\date{říjen 2011}


% ############################################################## BIBLIOGRAPHY #

\usepackage[style=authoryear]{biblatex}
\addbibresource{slovnik.bib}


% ##################################################################### INDEX #

\usepackage{imakeidx}

% Index of authors of photographs
\makeindex
  [ name = figures
  , title = {Seznam autorů fotografií a ilustrací}
  , intoc
  , columnseprule
  , columnsep = \columnsep
  , noautomatic                       % We handle this in make.sh
  ]

% icelandic rules start
\usepackage{filecontents}
\def\mygroup#1{{\bfseries#1}}
\begin{filecontents*}{style.xdy}
;; style.xdy
(markup-letter-group :open-head "~n\mygroup{" :close-head "}")
\end{filecontents*}
% icelandic rules end


% ############################################################### PAGE STYLES #

\usepackage{fancyhdr}

% -------------------------------------- BASIC PAGE STYLE
\fancypagestyle{plain}{
  \fancyhf\relax                       % Clear header/footer
  \renewcommand \headrulewidth {0.0pt} % No header rule
  \renewcommand \footrulewidth {0.0pt} % No footer rule
  \fancyfoot[C]{\thepage}              % Page number in footer, centred
}

% -------------------------------------- INDEX PAGE STYLE  
%\fancypagestyle{indexstyle}{
%  \fancyhf\relax                       % Clear header/footer
%  \renewcommand \headrulewidth {0.4pt} % Thin header rule
%  \renewcommand \footrulewidth {0.0pt} % No footer rule
%  \fancyhead[C]{\thepage}              % Page number in footer, centred
%}

% -------------------------------------- DICTIONARY PAGE STYLE
\fancypagestyle{myheadings}{
  \fancyhf\relax                       % Clear header/footer
  \renewcommand \headrulewidth {0.4pt} % Thin header rule
  \fancyhead[CO,CE]{\thepage}          % Page number in header, centred
  % NOTE: the page numbers will be printed when the dictionary is ready
  \fancyhead[LE,LO]{\phvfamily\bfseries\rightmark}
  \fancyhead[RE,RO]{\phvfamily\bfseries\leftmark}
}

\pagestyle{plain}  

% ##################################################################### LISTS #

\usepackage{expdlist}
\usepackage{enumitem}

\setlist{nosep, topsep=\baselineskip}
\setlist[itemize]{leftmargin=*,label=\scalebox{.6}{\textbullet}}


% #################################################################### FLOATS #

\usepackage{caption}
\captionsetup{labelformat=empty}

% Alter some LaTeX defaults for better treatment of figures:
% See p.105 of "TeX Unbound" for suggested values. 
% See pp. 199-200 of Lamport's "LaTeX" book for details.
% ------------------------------------- Parameters for ALL pages:
\renewcommand \topfraction    {0.9}   %   max fraction of floats at top
\renewcommand \bottomfraction {0.8}   %   max fraction of floats at bottom
% ------------------------------------- Parameters for TEXT (non float) pages:
\setcounter{topnumber}    {1}
\setcounter{bottomnumber} {1}
\setcounter{totalnumber}  {1}         %   2 may work better
\setcounter{dbltopnumber} {2}         %   for 2-column pages
\renewcommand \dbltopfraction   {0.9} %   fit big float above 2-col. text
\renewcommand \textfraction     {0.07}%   allow minimal text w. figs
% ------------------------------------- Parameters for FLOAT (non text) pages:
\renewcommand \floatpagefraction{0.7} %   require fuller float pages
% N.B.: floatpagefraction MUST be less than topfraction !!
\renewcommand{\dblfloatpagefraction}{0.7} % require fuller float pages
% remember to use [htp] or [htpb] for placement

\makeatletter
\setlength{\@fptop}{0pt}
\setlength{\@fpbot}{0pt plus 1fil}
\makeatother

% This kills spacing around floats
\setlength \intextsep        {0pt}
\setlength \floatsep         {0pt}
\setlength \textfloatsep     {0pt}
\setlength \dbltextfloatsep  {0pt}
\setlength \dblfloatsep      {0pt}

% This kills spacing around captions
\setlength \abovecaptionskip {0pt}
\setlength \belowcaptionskip {0pt}


% #################################################################### TABLES #

\usepackage{booktabs}
\usepackage{tabularx}
\usepackage{longtable}
\usepackage{ltxtable}
\usepackage{lscape}

\ifPDFrotate
\else
\makeatletter			       
\renewcommand\PLS@Rotate[1]{}          % this command prevents pdflscape 
\makeatother                           % to rotate pages
\fi

%\LetLtxMacro\oldtabularx\tabularx
%\LetLtxMacro\oldendtabularx\endtabularx
%\renewenvironment{tabularx}[1]
%  {\oldtabularx{\columnwidth}{#2}}
%  {\oldendtabularx}

\newcommand\LTXfw[1]{\LTXtable{\columnwidth}{#1}}

%%% \LetLtxMacro\oldtabularx\tabularx
%%% \LetLtxMacro\oldendtabularx\endtabularx
%%% \newenvironment{longtabularx}[2]
%%%   {\oldtabularx{#1}{#2}}
%%%   {\oldendtabularx}

% ######################################################### czech babel vs cmidrule
% ###################################### http://tex.stackexchange.com/a/112001/8380
\begingroup
  \makeatletter
  \catcode`\-=\active
  \AtBeginDocument{
  \def\@@@cmidrule[#1-#2]#3#4{\global\@cmidla#1\relax
    \global\advance\@cmidla\m@ne
    \ifnum\@cmidla>0\global\let\@gtempa\@cmidrulea\else
    \global\let\@gtempa\@cmidruleb\fi
    \global\@cmidlb#2\relax
    \global\advance\@cmidlb-\@cmidla
    \global\@thisrulewidth=#3
    \@setrulekerning{#4}
    \ifnum\@lastruleclass=\z@\vskip \aboverulesep\fi
    \ifnum0=`{\fi}\@gtempa
    \noalign{\ifnum0=`}\fi\futurenonspacelet\@tempa\@xcmidrule}
  }
\endgroup

% macro that supress cmidrule at page brake
% http://tex.stackexchange.com/questions/110548/using-longtable-with-cmidrule-page-breaking-issues
\makeatletter
\def\LT@output{%
  \ifnum\outputpenalty <-\@Mi
    \ifnum\outputpenalty > -\LT@end@pen
      \LT@err{floats and marginpars not allowed in a longtable}\@ehc
    \else
      \setbox\z@\vbox{\unvbox\@cclv}%
      \ifdim \ht\LT@lastfoot>\ht\LT@foot
        \dimen@\pagegoal
        \advance\dimen@-\ht\LT@lastfoot
        \ifdim\dimen@<\ht\z@
          \setbox\@cclv\vbox{\unvbox\z@
\copy\LT@foot\vss}%
          \@makecol
          \@outputpage
          \setbox\z@\vbox{\box\LT@head}%
        \fi
      \fi
      \global\@colroom\@colht
      \global\vsize\@colht
      \vbox
        {\unvbox\z@
\box\ifvoid\LT@lastfoot\LT@foot\else\LT@lastfoot\fi}%
    \fi
  \else
    \setbox\@cclv\vbox{\unvbox\@cclv

\setbox0\lastbox
\ifdim\ht0=0.29999pt \unskip % hope that was a cline we threw away
\else
\nointerlineskip
\box0     % put it back, whatever it was
\fi

\copy\LT@foot\vss}%
    \@makecol
    \@outputpage
      \global\vsize\@colroom
    \copy\LT@head\nobreak
  \fi}

\makeatother

% These values ensure that for top, mid and bottom rules
% space above + rule width + space below = baselineskip
\setlength \abovetopsep    {0.43\baselineskip}
\setlength \heavyrulewidth {0.10\baselineskip}
\setlength \belowbottomsep {0.43\baselineskip}
\setlength \aboverulesep   {0.47\baselineskip}
\setlength \lightrulewidth {0.06\baselineskip}
\setlength \belowrulesep   {0.47\baselineskip}

%\usepackage{letltxmacro}

%%% \LetLtxMacro\oldtabular\tabular
%%% \LetLtxMacro\oldendtabular\endtabular
%%% \renewenvironment{tabular}[2][t]
%%%   {\oldtabular[#1]{#2}}
%%%   {\oldendtabular}

%%% \LetLtxMacro\oldtable\table
%%% \LetLtxMacro\oldendtable\endtable
%%% \renewenvironment{table}[1][h]
%%%   {\leavevmode\oldtable[#1]\vspace{\dp\strutbox}}
%%%   {\vspace{-\dp\strutbox}\oldendtable}
%%%   % This aligns a [t] aligned tab-envo to baseline grid


%%% % declention and conjugation tables
%%% \usepackage{floatrow}
%%% \DeclareFloatFont{footnotesize}{\footnotesize}
%%% % "scriptsize" is defined by floatrow, "tiny" not
%%% \floatsetup[table]{font=footnotesize}

% ################################################################### FIGURES #

\usepackage{adjustbox} % loads also graphicx

\graphicspath{%
  {/home/chejnik/Dokumenty/hvalur.org/images/biolib/full/}%
  {/home/chejnik/Dokumenty/hvalur.org/images/uploaded_files/}}

\newlength\dicfigheight
\newlength\dicfigskip

\newcommand\dicFigure[3]{ % {filename}{caption}{index entry}
  \setlength\fboxsep{0pt}\setlength\fboxrule{0.5pt}
  \def\dicfigbox{\fbox{%
    \adjincludegraphics [ max height = 0.8\columnwidth
                        , max width  = 0.8\columnwidth ] {#1}}}
  \setlength\abovecaptionskip {\dp\strutbox}
  \setlength\belowcaptionskip {\ht\strutbox} 
  \setlength\dicfigheight {\heightof\dicfigbox+\fboxrule}
  \setlength\dicfigskip
    {\baselineskip*\numexpr1+\dicfigheight/\baselineskip\relax-\dicfigheight}
  \begin{figure}[ht]
    \centering\vspace*{\dicfigskip}\dicfigbox\caption{#2}\index[figures]{#3}
  \end{figure}
}


% #################################################################### COLORS #

\usepackage{color}

% http://www.colorschemer.com/online.html
% main color - grammar color
% #66023C
\definecolor {darkgreen}          {rgb} {0.40, 0.01, 0.24}

\definecolor {royalazure}         {rgb} {0.00, 0.22, 0.66}
\definecolor {brown}              {rgb} {0.40, 0.01, 0.24}

% COLORS FOR THUMB INDEXES
\definecolor {babyblueeyes}       {rgb} {0.63, 0.79, 0.95}
\definecolor {unitednationsblue}  {rgb} {0.36, 0.57, 0.90}
\definecolor {blue(ryb)}          {rgb} {0.01, 0.28, 1.00}
\definecolor {darkblue}           {rgb} {0.00, 0.00, 0.55}
\definecolor {screamingreen}      {rgb} {0.46, 1.00, 0.44}
\definecolor {limegreen}          {rgb} {0.20, 0.80, 0.20}
\definecolor {islamicgreen}       {rgb} {0.00, 0.56, 0.00}
\definecolor {upforestgreen}      {rgb} {0.00, 0.27, 0.13}
\definecolor {icterine}           {rgb} {0.99, 0.97, 0.37}
\definecolor {orange(colorwheel)} {rgb} {1.00, 0.50, 0.00}
\definecolor {orange-red}         {rgb} {1.00, 0.27, 0.00}
\definecolor {oucrimsonred}       {rgb} {0.60, 0.00, 0.00}
\definecolor {cottoncandy}        {rgb} {1.00, 0.74, 0.85}
\definecolor {orchid}             {rgb} {0.85, 0.44, 0.84}
\definecolor {vividcerise}        {rgb} {0.85, 0.11, 0.51}
\definecolor {patriarch}          {rgb} {0.50, 0.00, 0.50}

\definecolor {title}              {RGB} {  16,   13,   32}
\definecolor {golden}             {RGB} { 241,  184,   45}
\definecolor {freqstar}           {RGB} {  91,    3,   99}
\definecolor {newdev}             {RGB} {   3,   11,   99}

\definecolor {color1}             {RGB} { 182,   86,    0}
\definecolor {color2}             {RGB} { 143,    9,    6}
\definecolor {color3}             {RGB} {   3,   23,  118}
\definecolor {color4}             {RGB} {   0,   82,  168}
\definecolor {color5}             {RGB} {   0,   85,  142}
\definecolor {color6}             {RGB} {   0,  115,  162}
\definecolor {color7}             {RGB} {  34,  146,  186}
\definecolor {color8}             {RGB} {  40,  159,  153}
\definecolor {color9}             {RGB} {   0,  125,  111}
\definecolor {color10}            {RGB} {   4,  107,   60}
\definecolor {color11}            {RGB} {  71,  134,   81}
\definecolor {color12}            {RGB} { 109,  134,   42}
\definecolor {color13}            {RGB} { 205,  194,   18}
\definecolor {color14}            {RGB} { 204,  162,   24}

%\definecolor {color1}             {RGB} { 239,  222,  205}
%\definecolor {color2}             {RGB} { 251,  206,  177}
%\definecolor {color3}             {RGB} { 251,  185,  142}
%\definecolor {color4}             {RGB} { 248,  166,  112}
%\definecolor {color5}             {RGB} { 216,  231,  246}
%\definecolor {color6}             {RGB} { 186,  215,  244}
%\definecolor {color7}             {RGB} { 152,  197,  243}
%\definecolor {color8}             {RGB} { 118,  181,  245}
%\definecolor {color9}             {RGB} { 181,  249,  185}
%\definecolor {color10}            {RGB} { 160,  248,  167}
%\definecolor {color11}            {RGB} { 128,  247,  136}
%\definecolor {color12}            {RGB} {  87,  248,   98}
%\definecolor {color13}            {RGB} { 250,  241,  138}
%\definecolor {color14}            {RGB} { 250,  240,  111}
%\definecolor {color15}            {RGB} { 244,  239,   85}
%\definecolor {color16}            {RGB} { 245,  239,   57}

\definecolor {lightgrey}          {RGB} { 105,  105,  105}
\definecolor {darkgreen_real}     {rgb} {0.00, 0.50, 0.00}


% ############################################################# THUMB INDEXES #

% Thumb indexes' colors
%\newcommand\BoxColor{%
%  \ifcase\theletternum darkgreen!30\or babyblueeyes\or unitednationsblue\or blue(ryb)\or screamingreen\or limegreen\or islamicgreen\or upforestgreen\or icterine\or orange(colorwheel)\or orange-red%
%  \or oucrimsonred\or cottoncandy\or orchid\or vividcerise\or patriarch\or babyblueeyes\or unitednationsblue\or blue(ryb)\or screamingreen\or limegreen\or islamicgreen\or upforestgreen\or icterine\or orange(colorwheel)\or orange-red%
%  \or oucrimsonred\or cottoncandy\or orchid\or vividcerise\or patriarch\else darkgreen!30\fi}
  
% Thumb indexes' colors
%\newcommand\BoxColor{%
%  \ifcase\theletternum darkgreen!30\or color1\or color2\or color3\or color4\or color5\or color6\or color7\or color8\or color9\or color10%
%  \or color11\or color12\or color13\or color14\or color1\or color2\or color3\or color4\or color5\or color6\or color7\or color8\or color9%
%  \or color10\or color11\or color12\or color13\or color14\or color1\or color2\or color3\or color4\or color5\or color6\else color7\or color8\or color9\or color10\or color11\or color12\fi}

%\def\BoxColor#1{darkgreen!\the\numexpr103-#1-#1-#1\relax!blue}

\newcommand\BoxColor[1]{%
\ifcase#1 darkgreen!30\or color13\or color14\or color1\or color2\or color3\or color4\or color5\or color6\or color7\or color8\or color9\or color10%
\or color11\or color12\or color13\or color14\or color1\or color2\or color3\or color4\or color5\or color6\or color7\or color8\or color9%
\or color10\or color11\or color12\or color13\or color14\or color1\or color2\or color3\or color4\or color5\or color6\else color7\or color8\or color9\or color10\or color11\or color12\fi}

\ifshowcrop
%%%%%%%%%%%%%%%%%%%%%%%%%%%%%%%%%%%%%%%%%%%%%%%%%%%%%%%%%%%%%%%%%%%%%%%%%%%%%%%%

\newcounter{letternum}
\newcounter{lettersum}                 \setcounter{lettersum}{34}
% these margins separate top thumb from top of page and last thumb from
% bottom of page. BEWARE, it's from top of the (cut out) page, not the folio!
\newlength{\thumbtopmargin}            \setlength{\thumbtopmargin}{1cm}
\newlength{\thumbbottommargin}         \setlength{\thumbbottommargin}{1cm}
\newlength{\thumbwidth}                \setlength{\thumbwidth}{2.20cm}
\newlength{\thumbheight}
\pgfmathsetlength{\thumbheight}%
  {(\csname Gm@layoutheight\endcsname-\thumbtopmargin-\thumbbottommargin)/\value{lettersum}}
% inner xsep of letter from box margin
\newlength{\thumbsep}                  \setlength{\thumbsep}{.30cm}
% overlap over cutting line
\newlength{\thumboverlap}              \setlength{\thumboverlap}{3cm}

\tikzset{
  thumb/.style={
    text=white,
    outer sep=0pt,
    inner xsep=\thumbsep,
    minimum height=\thumbheight,
    text width=\thumbwidth-2\thumbsep,
    font=\sffamily\bfseries,},
  thumb east/.style={thumb,
    xshift=\dimexpr(+\thumboverlap-\paperwidth+\csname Gm@layoutwidth\endcsname)/2\relax,
    align=left,anchor=north east,},
  thumb west/.style={thumb,
    xshift=\dimexpr(-\thumboverlap+\paperwidth-\csname Gm@layoutwidth\endcsname)/2\relax,
    align=right,anchor=north west,},
}

\def\thumbnew{}
\def\thumbold{}
\usepackage{everypage}
\AddEverypageHook{\if\relax\thumbnew\relax\xdef\thumbnew{\thumbold}\fi}

\def\ethumbs#1,#2\relax%
  {\if\relax#1\relax\else\eventhumb{#1}\fi%
   \if\relax#2\relax\else\ethumbs#2\relax\fi%
   \gdef\thumbnew{}%
   \gdef\thumbold{#1,}}

\def\othumbs#1,#2\relax%
  {\if\relax#1\relax\else\oddthumb{#1}\fi%
   \if\relax#2\relax\else\othumbs#2\relax\fi%
   \gdef\thumbold{#1,}%
   \gdef\thumbnew{}}

\newcommand{\drawthumb}[2]{%
  \begin{tikzpicture}[remember picture, overlay]
    \node [thumb #2, fill=\BoxColor{#1}]
       at ($(current page.north #2)-%
            (0,\csname Gm@layoutvoffset\endcsname+\thumbtopmargin-\thumbheight+%
               #1*\thumbheight)$) {\csname Let#1\endcsname};
   \end{tikzpicture}}



%%%%%%%%%%%%%%%%%%%%%%%%%%%%%%%%%%%%%%%%%%%%%%%%%%%%%%%%%%%%%%%%%%%%%%%%%%%%%%%%
\else
% THUMB INDEXES
% new counter to hold the current number of the letter to determine the vertical position
\newcounter{letternum}
% newcounter for the sum of all letters to get the right height of a box
\newcounter{lettersum}
\setcounter{lettersum}{34}
% some margin settings
\newlength{\thumbtopmargin}
\setlength{\thumbtopmargin}{\ifshowcrop 3cm\else 1.5cm\fi} %1cm
\newlength{\thumbbottommargin}
\setlength{\thumbbottommargin}{\ifshowcrop 6cm\else 2cm\fi} %2.5cm
% calculate the box height by dividing the page height
\newlength{\thumbheight}
\pgfmathsetlength{\thumbheight}{%
(\paperheight-\thumbtopmargin-\thumbbottommargin)%
/%
\value{lettersum}
}
% box width
\newlength{\thumbwidth}
\setlength{\thumbwidth}{\ifshowcrop 2cm\else 0.5cm\fi} %0.5cm
% style the boxes
\tikzset{
thumb/.style={
   text=white,
   minimum height=\thumbheight,
   text width=\thumbwidth,
   outer sep=0pt,
   font=\sffamily\bfseries,
 }
 }

\def\thumbnew{}
\def\thumbold{}
\usepackage{everypage}
\AddEverypageHook{\if\relax\thumbnew\relax\xdef\thumbnew{\thumbold}\fi}

\def\ethumbs#1,#2\relax{\if\relax#1\relax\else\eventhumb{#1}\fi%
                        \if\relax#2\relax\else\ethumbs#2\relax\fi%
                        \gdef\thumbnew{}%
                        \gdef\thumbold{#1,}%
}

\def\othumbs#1,#2\relax{\if\relax#1\relax\else\oddthumb{#1}\fi%
                        \if\relax#2\relax\else\othumbs#2\relax\fi%
                        \gdef\thumbold{#1,}%
                        \gdef\thumbnew{}%
}

\newcommand{\drawthumb}[2]{%
  % see pgfmanual.pdf for more information about this part
  \begin{tikzpicture}[remember picture, overlay]
    \node [thumb, fill=\BoxColor{#1}, text centered, anchor=north #2]
       at ($(current page.north #2)-%
%            (0,\thumbtopmargin+\value{letternum}*\thumbheight)%
            (0,\thumbtopmargin+#1*\thumbheight)$) {\csname Let#1\endcsname};
   \end{tikzpicture}}
\fi

\newcommand{\oddthumb} [1]{\drawthumb{#1}{west}}
\newcommand{\eventhumb}[1]{\drawthumb{#1}{east}}

\newcommand{\lettergroup}[1]%
  {\refstepcounter{letternum}%
   \expandafter\gdef\csname Let\theletternum\endcsname{#1}%
   \xdef\thumbnew{\theletternum,\thumbnew}%
   \fancyhead[LO]{\phvfamily\bfseries\rightmark%
     \expandafter\ethumbs\thumbnew\relax\relax\relax}%
   \fancyhead[RE]{\phvfamily\bfseries\leftmark%
     \expandafter\othumbs\thumbnew\relax\relax\relax}}
     
% ############################################################## LOCALIZATION #

% renames the index name
%\addto\captionsczech{%
%  \renewcommand\indexname{{Seznam autorů fotografií a ilustrací}}}

% renames the contents name
\addto\captionsczech{%
  \renewcommand\contentsname{Obsah}}

\renewcommand \contentsname {Obsah}
\renewcommand \chaptername  {Kapitola}


% ################################################################ SECTIONING #
\usepackage[explicit]{titlesec}

\setlength{\beforetitleunit}{\baselineskip}
\setlength{\aftertitleunit} {\baselineskip}

\titlespacing*{\chapter}   {0em}{*0}{*3} % Remembere there's a topskip too
\titlespacing*{\section}   {0em}{*1}{*2}
\titlespacing*{\subsection}{0em}{*1}{*2}

\newcommand\ghostdrop[2]{%
  \raisebox{-#1\baselineskip}[0pt][0pt]{#2}}

\titleformat {\chapter} {} {} {0pt} % {#1}
  {\ghostdrop{1.0}{\parbox[t]{\columnwidth}
     {\LARGE\bfseries\centering\MakeUppercase{#1}}}}

\def\longchapterskip{\vspace{2\baselineskip}}
\def\longsectionskip{\vspace{2\baselineskip}}

\titleformat {\section} {\bfseries} {} {0pt}
  {\ghostdrop{1.0}{\parbox[t]{\columnwidth}
     {\Large\bfseries\centering#1}}}

\titleformat {\subsection} {\bfseries} {} {0pt}
  {\ghostdrop{1.0}{\hspace{1em}\parbox[t]{\columnwidth-1em}
     {\large\bfseries#1}}}


\setlength\multicolsep{0pt}
%\BeforeBeginEnvironment{multicols}{\vspace{-\dp\strutbox}}


% ######################################################## TABLE OF CONTENTS #
\usepackage{titletoc}

\titlecontents {chapter} [0pt] {\addvspace\baselineskip\bfseries}
  {} {} {\titlerule*[2\baselineskip]{}\contentspage}

\titlecontents {section} [\baselineskip] {} {} {}
  {\titlerule*[1\baselineskip]{.}\contentspage}


% #################################################### ASSORTED CUSTOM MACROS #

\def\textIS#1{\foreignlanguage{icelandic}{#1}}
\def\textCS#1{\foreignlanguage{czech}{#1}}
\def\textLA#1{\foreignlanguage{latin}{#1}}

\usepackage{accsupp}

\makeatletter                                                                                                          
  \def\ACCSUPP@bdc{\special{pdf:code \ACCSUPP@span BDC}}                                                               
  \def\ACCSUPP@emc{\special{pdf:code EMC}}                                                                             
\makeatother

\usepackage{pgffor}

\newcommand\blspace[1][1]{\vspace{#1\baselineskip}}

\usepackage{xparse}


\newcommand \dicsymFrequent  {\raisebox{-.2ex}{\color{darkgreen}\ding{167}}}
\newcommand \dicsymSee       {$\shortrightarrow$}
\newcommand \dicsymCompare   {$\shortuparrow$}
\newcommand \dicsymIdiom     {$\diamondsuit$}
\newcommand \dicsymExampleIS {$\triangleright$}
\newcommand \dicsymExampleCS {\guilsinglright}
\newcommand \dicsymProverb   {{\color{newdev}\ding{96}}}
\newcommand \dicsymPhoto
  {\includegraphics[keepaspectratio,height=0.65em]{photo_icon}}

%\NewDocumentCommand\dicTerm{mo}
 % {{\phvfamily\bfseries\textIS{#1}}\IfValueT{#2}{, \phvfamily\bfseries\textIS{#2}}}

\NewDocumentCommand\dicEntry{o}
  {\par\hangpara{0.25\baselineskip}{1}%
   \IfValueT{#1}{\markboth{#1}{#1}}%
    \ifPDFsearchvalue
    \settowidth{\eyja}{#1}%
  \makebox[\eyja]{\color{white}#1}\hspace{-\eyja}%
  \fi
   \ignorespaces}
   
   \NewDocumentCommand\dicEntryGuide{o}
  {\par\hangpara{0.29\baselineskip}{1}%
   \IfValueT{#1}{\markboth{#1}{#1}}%
    \ifPDFsearchvalue
    \settowidth{\eyja}{#1}%
  \makebox[\eyja]{\color{white}#1}\hspace{-\eyja}%
  \fi
   \ignorespaces}
   
\newcommand\dicSubEntry
  {\\\hspace*{-\hangindent}}

\NewDocumentCommand\textAlt{mm}
  {\BeginAccSupp{method=pdfstringdef,ActualText={#2}}#1\EndAccSupp{}}

\newcounter{terms}

\NewDocumentCommand\dicStyleIntTerm{m}{\textIS{{\phvfamily\bfseries#1}}}
\NewDocumentCommand\dicStyleExtTerm{m}{\textIS{\textit{#1}}}

   \NewDocumentCommand\dicTerm{mo}
  {{\closeupdotbar \phvfamily\bfseries\textIS{#1}}\IfValueT{#2}{, \phvfamily\bfseries\textIS{#2}}}

  \newcommand\dicTermList[1]
  {\foreach \t [count=\n] in {#1}{\ifnum\n=1\else, \fi\dicTerm{\t}}}
  
\def\dicIPA#1{\textipa{[#1]}}

\NewDocumentCommand\dicPos{smo}
  {{\color{darkgreen}\textbf{\small#2}%
   \IfValueT{#3}{\nobreak\textsubscript{#3}}}}

%\NewDocumentCommand\dicFlx{smo}
 % {{\color{darkgreen}\footnotesize#1\IfValueT{#2}{\textsubscript{#2}}}}

  \NewDocumentCommand\dicFlx{smo}
  {{\color{darkgreen}\footnotesize#2%
   \IfValueT{#3}{\nobreak\textsubscript{#3}}}}
  
\newcommand\dicExampleCS[1]{\dicsymExampleCS~\textit{#1}}
\newcommand\dicExampleIS[1]{\dicsymExampleIS~\textIS{\textit{#1}}}

\NewDocumentCommand\dicLink{sm}
  {\IfBooleanF#1{\dicsymSee\nobreak\hspace{.08667em plus .03333em}}\nobreak\hspace{.03667em}\dicTerm{#2}}
\NewDocumentCommand\dicSynonym{sm}
  {(\IfBooleanF#1{\dicsymSee\nobreak\hspace{.08667em plus .03333em}}\nobreak\hspace{.03667em}\textIS{\textit{#2}})}
\NewDocumentCommand\dicAntonym{sm}
  {(\IfBooleanF#1{\dicsymCompare\nobreak\hspace{.08667em plus .03333em}}\nobreak\hspace{.03667em}\textIS{\textit{#2}})}
\NewDocumentCommand\dicCompare{sm}
  {(\IfBooleanF#1{\dicsymCompare\nobreak\hspace{.08667em plus .03333em}}\nobreak\hspace{.03667em}\textIS{\textit{#2}})}


\newcommand\dicPhraseIS[1]{\textIS{\textbf{#1}}}

\NewDocumentCommand\dicIdiom{mo}
  {\dicSubEntry{\color{newdev}\dicStyleIntTerm{#1} %
     \IfValueTF{#2}{+ \dicStyleIntTerm{#2} \markboth{#1 + #2}{#1 + #2}}{\markboth{#1}{#1}}}\dicsymIdiom}
     
\newcommand\dicProverb
  {\dicSubEntry\dicsymProverb}



\newcommand\dicLangCat[1]{{\footnotesize#1}}
\newcommand\dicFieldCat[1]{{\footnotesize#1}}


\newcommand\dicDirectTranslationCS[1]{#1}
\newcommand\dicIndirectTranslationCS[1]{{\footnotesize#1}}

\usepackage{placeins}

\newcommand*{\dicLetter}[2]{%
  \FloatBarrier\newpage\phantomsection\pdfbookmark[1]{#1}{#2}%
  \noindent\parbox[b][9\baselineskip][c]{\columnwidth}
    {\centering\HUGE\strut\smash{\MakeUppercase{#1}~\MakeLowercase{#1}}}%
      %\par\xdef\tpd{\the\prevdepth}
      %\par\prevdepth\tpd%
  \expandafter\lettergroup{\MakeLowercase{#1}}}

  
%tabulky
\newcommand{\specialcell}[2][l]{%
\begin{tabular}[#1]{@{}l@{}}#2\end{tabular}}


%\newcommand{\devisionguide}[1]{\hspace*{0.3cm}{{{{\foreignlanguage{icelandic}{\color{newdev}{\phvfamily{\textbf{#1}}}}}}}}}
%\newcommand{\devisionguideindent}[1]{\hspace*{-0.1cm}{{{{\foreignlanguage{icelandic}{\color{newdev}{\phvfamily{\textbf{#1}}}}}}}}}

%\newcommand\n{$n$\nobreakdash-\hspace{0pt}} %nonbreakable dash in grammar endings
%\def \nobreakseq {\nobreak \hskip 0pt \hbox} %nolinebreak in case like (-s, -)

% rule line
\newcommand{\HRule}{\rule{\linewidth}{0.1mm}}

% width of pdf found box
\newlength\eyja

%\newcommand\strutsmash[1]{\strut\smash{#1}}


% Dictionary symbols


% Command synonyms
%\let\freqstar\dicsymFrequent
%\let\photoicon\dicsymPhoto

\newcommand{\vstretch}[1]{\vspace*{\stretch{#1}}}

\newcommand*{\addthin}{\hskip0.16667em\relax}
\newcommand*{\addthinS}{\hskip0.06667em\relax}
\newcommand*{\addthinM}{\hskip0.10007em\relax}
\newcommand*{\addthinSS}{\hskip0.00007em\relax}
% #################################################################### COVERS #

\def\dictnameCZ{islandsko-český studijní slovník}
\def\dictnameIS{íslensk-tékknesk stúdentaorðabók}

\newcommand*{\frontcoverimages}{%
  ds_image_blaberjalyng_0_2.jpg,
  ds_image_lundi_0_1.jpg,
  ds_image_hestur_0_1.jpg,
  ds_image_hrafnafifa_0_2.jpg,
  ds_image_baldursbra_0_1.jpg,
  20948.jpg,
  ds_image_hvalur_0_1.jpg,
  ds_image_blodberg_0_1.jpg}

\newcommand*{\backcoverimages}{%
  ds_image_spoi_0_1.jpg,
  ds_image_mosalyng_0_1.jpg,
  ds_image_blaklukkulyng_0_1.jpg,
  20787.jpg,
  ds_image_ufsi_0_2.jpg,
  %ds_image_gullbra_0_1.jpg,
  ds_image_vatnsberi_0_1.jpg,
  ds_image_islandsfifill_0_2.jpg,
  21137.jpg}

\newcommand\makecoverwith[1]{
 \pagecolor{title}\afterpage\nopagecolor
 \begin{center}
 \vspace*{0.3cm}%{\fill}
 {\Huge\scshape\color{white}%
   \dictnameCZ\\\dictnameIS}
 \vspace*{1em}                        % I think this is necessary
 \begin{figure}[ht]\centering%
   \setlength\fboxsep{0pt}\setlength\fboxrule{0.5pt}%
   \foreach \file in #1{
     \fbox{\includegraphics[width=5.5cm]{\file}}}
 \end{figure}
 \vspace*{\fill}
 \end{center}
}


% ###################################################### HYPERREF & PDF INFO #

% unicode neccessary so that national characters in hypersetup appear ok
\usepackage[pdftex, unicode, hyperfootnotes=false]{hyperref}

% hyperlinks in black
\makeatletter
\let\Hy@linktoc\Hy@linktoc@none
\makeatother

\hypersetup
  { pdftitle={Islandsko-český studijní slovník}
  , pdfauthor={Aleš Chejn, Jón Gíslason, Ján Zaťko}
  , pdfsubject={Islandsko-český studijní slovník}
  , pdfkeywords={Islandsko-český studijní slovník%
               , islandština, čeština, slovník}
  , bookmarks=true
  , colorlinks=true
  , citecolor=black
  , urlcolor=\ifscreen darkgreen\else black\fi
  , linkcolor=black}


% ################################################################# MICROTYPE #

\usepackage[final]{microtype}

% ############################################################### HYPHENATION #

% this allows to include brackets inside hyphenation rules f.e (pro)-dis-ku-to-vat
\lccode`\(`\(
\lccode`\)`\)
\lccode`\[`\[
\lccode`\]`\]

\begin{hyphenrules}{icelandic}
\hyphenation{
  frum-semja
  burm-neskur
  ó-breytan-le-gur
  pla-tó-nskt
  skatta-fram-tal
  skrúf-gang-ur
  strand-rauð-viður
  prent-sver-ta
  grun-nur
  uppi-skrop-pa
  ger-vi-mennska
  sö-lu-stað-ur
  sý-na-gó-ga
  mark-mið
  trau-st-ur
  ís-kyggi-leg-ur
  brjó-sta-höld
  kné-krjú-pa
  nel-li-ka
  kaffi-kvörn
  margs-konar
  skrúf-gang-ur
  uppi-skrop-pa
  veð-rað-ur
  vest-ur-átt
  mis-mun-andi
  Bar-ba-do-si
  tung-u-mál
  au-ga-brún
  há-sker-pu-sjón-varp
  rúnn-stykki
  ger-vi-tungl
  yfir-borðs-menn-ska
  gæ-fu-mun-ur
  hlí-fi-skil-di
  á-vís-un
  ó-um-flýj-an-leg-ur
  san-mar-ín-sk-ur
  gæ-ða-stað-all
}
\end{hyphenrules}
\begin{hyphenrules}{czech}
\hyphenation{
  teo-rie
  mi-liardy
  (z)-od-po-věd-nost
  tan-tié-ma
  alian-ce
  pe-dia-tr(ička)
  pa-cienty
  zá-chra-nář(ka)
  dia-man-to-vá
  poe-zie
  ma-te-riál-ní
  egoi-s-ta
  egoist-ka
  nai-vi-ta
  ne-upřím-nost
  pě-ti-úhel-ník
  še-s-ti-úhel-ník
  za-šmo-dr-chat
  ar-cheo-lo-gic-ký
  ne-omlu-vi-tel-ný
  (vy)-pro-du-ko-vat
  dia-gno-s-ti-ko-vat
  teo-lož-ka
  ateist-ka
  oceá-nu
  im-pe-ria-li-s-ta
  (vy)-zdvih-nout
  ne-úměr-ný
  kon-ti-nuál-ní
  neu-t-rum
  ne-opráv-ně-ně
  vig-vam
  pro-zaic-ký
  geo-lož-ka
  seiz-mic-ká
  teo-rém
  olean-dr
  poe-ta
  (vy)-fo-to-gra-fo-vat
  (vy)-ven-ti-lo-vat
  pneu-mo-nie
  Lu-cem-bur-čan-(ka)
  ma-te-riály
  am-bi-cióz-ní
  Myan-ma(r)
  Myan-ma-řan-(ka)
  mol-dav-šti-na
  si-tua-cí
  ne-ukáz-ně-ný
  he-te-ro-trof-ní
  ne-omlu-vi-tel-ná
  ne-upřím-ný
  dvou-bliz-ný
  ne-úmy-sl-ný
  ne-ustá-lý-mi
  ne-ohle-du-pl-ný
  ne-osob-ní
  ne-oso-le-ný
  ne-uvě-ři-tel-ný
  ne-ohro-že-ný
  ne-oby-čej-ný
  ne-útul-ný
  ne-únav-ný
  ne-ože-nit
  po-ukáz-ka
  za-úto-čit
  aero-dy-na-mic-ký
  hru-bián-ství
  (pro)-di-s-ku-to-vat
  (vy)-ča-lou-nit
  (na)-oč-ko-vá-ní
  Bhú-tán-ka
  (za)-šně-ro-vat
  soud-ní
  re-vo-lu-cio-nář(ka)
  á-zer-báj-d-žán
  čtyř-hran-ný
  ze-m-dle-lý
  klou-zač-ka
  Vest-mann-ské
  tří-stran-ný
  (za)-bu-rá-ce-ní
  os-tud-ný
  pro-du-cent-(ka)
  chvá-s-tal-(ka)
  apar-theid
  asis-ten-ce
  asis-tent-(ka)
  (u)-či-nit
  (u)-za-mknout
  past-vi-na
  (za)-zá-lo-ho-vat
  Af-ri-čan-(ka)
  (za)-an-ga-žo-vá-ní
  se-rióz-ní
  opo-nent-(ka)
  pro-tiv-ník
  zvlášt-ní
  (po)-div-ný
  ho-mo-se-xuál-ní
  ko-alič-ní
  si-tua-ce
  si-tua-cí
  egoi-s-mus
  spe-ciál-ní
  spe-cia-li-zo-va-ný
  spe-cia-li-zo-va-né
  sobě-stač-ný
  sa-mo-zvo-le-ný
  sa-mo-čin-ně
 chvá-s-ta-vost
 klient-(ka)
 ku-rióz-ní
 (na)-šrou-bo-vat
 (za)-šrou-bo-vat
 krea-ti-vi-ta
 pneu-ma-ti-ka
 toa-le-ta
 so-cia-list-ka
 (vy)-dez-in-fi-ko-va-ný
 vše-uměl
 vše-uměl-(kyně)
 (pře)-škrt-nout
 (vy)-svlék-nout
 vy-učo-va-cí
 vy-učo-va-cích
 (vy)-chla-s-tat
 (vy)-te-to-vat
 pa-cient-(ka)
 ne-opod-s-tat-ně-ný
 de-se-ti-úhel-ník
 pe-rio-di-kum
 (za)-mum-lá-ní
 ne-smě-lá
 ne-dů-vě-ři-vá
 tee-na-ger-(ka)
 (z)-de-ge-ne-ro-vá-ní
 mo-ci-chti-vý
 bez-úhon-nost
 zne-uctít
 vo-dě-odol-ný
 we-bo-vý
 web-de-sign
 in-že-nýr-(ka)
 zá-pa-do-is-land-ský
 klient-(ka)
 za-uja-tý
 vy-uží-va-né
 vy-uží-vat
 (za)-re-gi-s-t-ro-vat
 (z)-kom-pri-mo-va-ný
 na-cio-na-li-s-mus
 ma-k-ro-eko-no-mic-ká
 bez-bo-lest-ný
 viol-ka
 troj-úhel-ník
 (za)-bu-rá-ce-ní
 ne-ústup-nost
 (vy)-cvi-če-ný
 zprá-va
 ge-ria-tr-(ička)
 osud-ný
 ná-pra-vy
 vstup-ní
 obec-ně
 re-pub-li-ka
 (po)-div-ný
 pa-slun-ce
 vliv-ný
 ba-ham-ský
 pe-dia-tr-(ička)
 ha-sič-(ka)
 vnuč-ka
 bro-s-kvoň
 lo-di-vod-(ka)
 mi-s-tr-(yně)
 in-že-nýr-(ka)
 dr-ži-tel-(ka)
 to-pič-(ka)
 pro-mi-nent-ní
 takt-ní
 frag-men-tár-ní
 po-zo-ru-hod-ný
 úspě-š-ný
 trans-fu-ze
 bur-žoa-zie
 pře-stro-je-ní
 oh-ni-vzdor-ný
 dis-tan-ční
 de-frau-da-ce
 fi-s-kál-ní
 zběh-lý
 kvap-ný
 střed-ní
 sa-mo-hlá-s-ka
 pře-vo-di-tel-ný
 zban-k-ro-to-va-ný
 zá-clo-na
 zná-mý
 zá-klad-ní
 na-uka
 ro-š-ťa-č-ka
 lí-bán-ky
 sou-du
 zkou-š-ka
 mlý-nek
 prá-šiv-ka
 kul-ma
 la-hod-ný
 zá-meč-ni-ce
 ne-vý-znam-ný
 ovzdu-ší
 mno-ho-úhel-ní-ko-vý
 žra-vý
 zcest-ný
 esen-ce
 (při)-nu-ce-ný
 skr-b-lic-ký
 skle-s-lý
 bez-bož-ný
 ne-opráv-ně-ný
 ošu-mě-lý
 pru-der-ní
 chvá-s-ta-vý
 sou-drž-ný
 van-dal-(ka)
 sa-bo-tér-(ka)
 (po)-div-ně
 mo-k-řad
 plyt-ký
 na-spo-du
 izo-la-ce
 le-toun
 za-han-be-ný
 ří-ze-ní
 kom-plex-ní
 pr-kén-ko
 ba-vič-(ka)
 plyt-ký
 ne-mo-tor-ně
 šéf-(ka)
 přes-čas
 zji-s-ti-tel-ný
 ta-ko-vý
 kňou-ral-(ka)
 se-me-s-tr
 po-pra-s-ka-ný
 han-ba
 vzpří-me-ný
 chys-ta-ný
 lik-vi-da-ce
 or-ga-ni-za-ce
 ri-s-kant-ní
 (za)-pís-kat
 (za)-hvíz-dat
 vrou-cí
  }
\end{hyphenrules}


% ############################################################# CZECH TYPOGRAPHY #
% https://math.feld.cvut.cz/tkadlec/ftp/text/textypo.pdf



% ################################################################# DOCUMENT #

\begin{document}
\pagenumbering{gobble}

% -------------------------------------------------------------- Front cover -

\ifmakecovers
\covergeometry
  \makecoverwith\frontcoverimages
\restoregeometry  
\fi


% ============================================================= FRONT MATTER =

\ifmakefrontmatter

% ------------------------------------------------------------------ Patitul -

\cleardoublepage

\begin{center}
  %\vspace*{2cm}{\fontsize{20}{20}\selectfont\scshape\dictnameCZ\\\dictnameIS}
  \vspace*{\stretch 1}
  {\Huge\scshape\dictnameCZ\\\dictnameIS}
  \vspace{\stretch 4}
\end{center}

% -------------------------------------------------------------------- Titul -

\cleardoublepage

\begin{center}
  %\vspace*{2cm}{\fontsize{24}{24}\selectfont\scshape\dictnameCZ\\\dictnameIS}
  \vspace*{\stretch 1}
  {\HUGE\scshape\dictnameCZ\\\dictnameIS}
  \vspace{\stretch 2}
\end{center}

% ----------------------------------------------------------------- Impresum -

\clearpage

\noindent\begin{minipage}[t][\textheight][t]{\textwidth}
Autoři \textit{Islandsko-českého studijního slovníku}:\\
  \textbf{Aleš Chejn, Ján Zaťko, Jón Gíslason}\\[\baselineskip]
\ifPDF
\else
Autor \textit{Islandské gramatiky}:\\
  \textbf{Vojtěch Kupča}\\[\baselineskip]
\fi
Poděkování:\\
  \textbf{Dorotě Nierychlewské-Chejn, Jiřině Hanzlové, Jiřině Chejnové,
          Simoně Petrásové,\\ Petru Mikešovi, Vojtěchu Kupčovi, 
          Renatě Peškové Emilsson, Amiru Mulahumici, \\ Kristýně Antonové, 
          Martině Kašparové, Lucii Brabcové, Samoriele, Adamu Kožouškovi}\\[\baselineskip]
Finanční příspěvek na výtisk exemplářů určených knihovnám\\ v České republice,
Slovenské republice a na Islandu poskytli:\\
  \textbf{Petr Mikeš, Dorota Nierychlewska-Chejn, Jiřina Chejnová, Otta Chejn}
  \ifPDF
  \\[39\baselineskip]
  \else
  \\[36\baselineskip]
  \fi
\textbf{Hvalur.org - 2016}\\
Licence: \textbf{GNU Free Documentation License}\\
Sazba: \textbf{Aleš Chejn} v prostředí  \textbf\LaTeX\\
  Zvláštní poděkování patří \textbf{Paolu al-Imk\=an\=\i\ Brasolin} a uživatelům  \textbf{\TeX\ - \LaTeX\ Stack Exchange}\\[\baselineskip]
\ifPDF
Verze: \textbf{PDF verze} \textit{Islandsko-českého studijního slovníku} vhodná pro \textbf{\ifscreen offline prohlížení a vyhledávání\else vlastní tisk\fi}\\ volně dostupná na \url{www.hvalur.org}\\[\baselineskip]
\else
\textbf\ISBN\\[\baselineskip]
\fi
\ifPDF
\else
Tisk: \textbf{Tiskárna Melmen, Pardubice}\\
\fi
Heslová část vysázena patkovým fontem \TeX\ Gyre Pagella,
hesla bezpatkovým fontem {\phvfamily Helvetica}.
\end{minipage}

% -------------------------------------------------------- Table of contents -

\cleardoublepage 

\pdfbookmark\contentsname{toc}    
\tableofcontents
\thispagestyle{empty}
\cleardoublepage 

\fi

\pagenumbering{arabic}
\setcounter{page}{3}
% ============================================================= INTRODUCTION =

\ifinputintroduction

% ------------------------------------- The study Icelandic-Czech dictionary -

\chapter{O Islandsko-českém studijním slovníku}
Cílem vytvoření\textit{ Islandsko-českého studijního slovníku} je zpřístupnění islandského jazyka všem zájemcům o~tento jazyk. 

\textit{Islandsko-český studijní slovník} je určen především českým a~slovenským studentům islandštiny, začátečníkům i~pokročilým, kteří zde naleznou ucelené fonetické, morfologické a~syntaktické informace, jakož
i~slovní zásobu potřebnou k~aktivnímu i~pasivnímu užití s~rozsahem vhodným ke studiu islandštiny na vysoké škole.
Slovník je rovněž určen všem zájemcům o~islandštinu a~Island, jimž slovník nabízí encyklopedické informace, úvod do fonetiky a~gramatiky a~velké množství fotografií a~ilustrací. 
Třebaže se jedná o~slovník primárně zaměřený na potřeby českých a~slovenských uživatelů, obsahuje také řadu informací, které usnadní práci islandskému uživateli (např. použití více jak 10 000 synonym a~antonym, překlady příkladů ap.).

Slovní zásoba původně vychází z~\textit{Concise Icelandic-English Dictionary} (\cite {ic_en}), značně však byla rozšířena díky přístupu k~webovému projektu \textit{ISLEX} (\cite {int1}). Velké množství hesel bylo
přidáno na podnět uživatelů online verze \textit{Islandsko-českého studijního slovníku}.

Slovník obsahuje 30~575 hesel s~46~725 významy nebo slovními spojeními. Z~celkového počtu 30~575 je 28~097 svébytných hesel obsahujících alespoň jeden význam nebo slovní spojení. 
Zbylá hesla (v~počtu 2~478) jsou varianty, dublety nebo nepravidelné tvary.

Slovník je primárně online elektronický slovník, volně přístupný na internetových stránkách \url{www.hvalur.org}  (\cite {int14}). 
Tištěná verze slovníku je publikována na vlastní náklady v~počtu 30 kusů, přičemž většina exemplářů je určena knihovnám a~institucím v~České republice, Slovenské republice a~na Islandu
a~několik exemplářů je určeno autorům nebo zájemcům z~řad veřejnosti. Dvě PDF verze (vyhledávání\,/\addthin prohlížení offline a~vlastní tisk) jsou dostupné zdarma na internetových stránkách slovníku (viz výše).

\textit{Islandsko-český studijní slovník} je otevřený svobodný projekt, jehož se v~průběhu tvorby zúčastnilo velké množství spoluautorů (viz Historie tvorby IČSS), kterým tímto jménem autorů velmi děkuji.
Děkuji také všem uživatelům online verze slovníku za trpělivost, dobrou vůli a~veškeré připomínky a~návrhy, jmenovitě především Petru Mikešovi, Adamu Kožouškovi a~Samoriele za neutuchající víru v~náš projekt.
Poděkování patří také našim rodinám a~blízkým, díky kterým jsme se mohli soustavně věnovat tvorbě slovníku.
Srdečné poděkování patří mé ženě Dorotce za podporu a~nadšení potřebné k~tvorbě slovníku.

\blspace[5]

{\centering Aleš Chejn, Brwinów, 24. prosince 2015\par}


% ------------ The history of the formation Icelandic-Czech dictionary study -

\chapter{Historie tvorby Islandsko-českého studijního slovníku}
\longchapterskip
Projekt \textit{Islandsko-českého studijního slovníku} (IČSS) vznikl v~roce 2001 na nepřímý podnět doc. PhDr. Heleny Kadečkové, CSc. z~Univerzity Karlovy,
kdy byl vypracován seznam slov. 
Tento seznam obsahoval islandská slova a~jejich překlady, později byly přidány základní morfologické informace o~slovním druhu a~deklinační\,/\addthin konjugační koncovky podle vzoru z~\textit{Concise Icelandic-English Dictionary} (\cite {ic_en}). 
Seznam v~této době obsahoval okolo 1 200 slov. Byl zveřejněn pomocí statických webových stránek a~později byl umístěn na webu Univerzity Karlovy. 

V~roce 2006 byl slovník rozšířen o~nových 1 500 heslových slov a~stávající hesla byla rozšířena o~nové významy, zvláště významy frázových sloves. 
Ke konci roku 2006  byl slovník převeden z programu Excel do databáze MySQL a~byly zahájeny práce na vytvoření aplikace \textit{Dictionary System} (DS) (\cite {int19}), která by podporovala tvorbu IČSS. 
Obsah slovníku byl zveřejněn pod svobodnou licencí na webových stránkách \url{www.hvalur.org} (\cite {int14}), kde se nachází dodnes. 

V~roce 2007 byla vytvořena první verze aplikace DS. V~této době se k~projektu přidala Renata Pešková Emilsson, která se nyní věnuje podpoře výuky češtiny na Islandu. 
Ta slovník obohatila rozsáhlým seznamem slov z~praktického života na Islandu. Dále se slovník rozšiřoval o~heslová slova podle \textit{Concise Icelandic-English Dictionary} (\cite {ic_en}). 
V průběhu následujících třech let bylo doplněno přibližně 15~000 slov. V této době byla hesla rozšířena o~synonyma, příklady, frekvenci, oborovou a~stylovou charakteristiku. K projektu se na Islandu připojila na krátkou dobu Kristýna Antonová, studentka islandštiny na Islandské univerzitě, která 
nás upozornila na nesrovnalosti v~uživatelském rozhraní aplikace DS.

Během léta 2009 byl vytvořen skript na generování skloňování a~časování heslových slov. Výsledky generování byly porovnány s~\textit{Beygingarlýsing íslensks nútímamáls} (\cite {DalvikVM}) a~během následujících let konzultovány
s~Jónem Gíslasonem. Kristín Bjarnadóttir, vedoucí projektu \textit{Beygingarlýsing íslensks nútímamáls}, zodpověděla řadu našich dotazů týkajících se deklinací a~konjugací.

V~roce 2009 začala spolupráce s~\textit{Biolib.cz} (\cite {int5}), jmenovitě s~Ondřejem Zichou. Obdrželi jsme  databázi s~náhledy obrázků, autory a~licencemi a~latinskými názvy botanických a~zoologických druhů.

V~zimním semestru 2009 jsme požádali s~Amirem Mulahumicem, bosenským studentem islandštiny na Islandské univerzitě, Studentský fond Islandské univerzity o~grant na tvorbu zvukové databáze islandských slov. 
Studentský fond návrh přijal. Peníze z~grantu jsme určili pro islandského rodilého mluvčího, Jóna Gíslasona, lektora islandského jazyka na Islandské univerzitě v~Reykjavíku, který nahrál více jak 22~000 islandských slov do elektronické formy. 
S~Dorotou Nierychlewskou-Chejn jsme připravili dva seznamy heslových slov (jeden pro nahrávajícího a~druhý pro zpracování nahrávek). Na podporu nahrávek jsme vytvořili webové stránky nazvané \textit{Databáze s~výslovností islandských slov} (\cite {int13}) a~tyto stránky jsme přeložili do 6~jazyků (Jón Gíslason do islandštiny, Ján Zaťko do slovenštiny, francouzštiny a~angličtiny, Amir Mulamuhic do srbochorvatštiny a~Dorota Nierychlewska-Chejn do polštiny).

V~roce 2010 jsme přidali do databáze pravidla islandské výslovnosti, která jsou uvedena v~knize \textit{Handbók um íslenskan framburð} (\cite {is2}). Z~těchto pravidel jsme automaticky vygenerovali fonetický zápis výslovnosti všech hesel ve slovníku. 
Jón Gíslason ověřil správnost fonetického zápisu u~všech heslových slov.

Na konci roku 2011 jsme spojili webové stránky slovníku se stránkami aplikace a~sjednotili jsme a~vylepšili vzhled webových stránek. Cennými připomínkami se na vytvoření aplikace podílela
RNDr. Ingrid Nagyová, PhD., vysokoškolská pedagožka Ostravské Univerzity. Ján Zaťko přeložil nové webové stránky, Průvodce po slovníku do slovenštiny, francouzštiny a~angličtiny a~celý manuál aplikace do angličtiny. 
Dorota Nierychlewska-Chejn přeložila stránky do polštiny, Jón Gíslason do islandštiny a~Marcos Helena do portugalštiny.

Od prosince 2011 do března 2014 jsme realizovali první čtení, ve kterém jsme doplnili chybějící informace, přidali nová hesla a~opravili pravopisné chyby. 
Obsah slovníku jsme porovnávali především s~\textit{Íslensk orðabók} (\mbox{\cite {is}}) a~\textit{Islandsko-českým slovníkem} (\cite {is7}). S~Vojtěchem Kupčou jsme konzultovali mnohé nesnáze při překladech.

V~letech 2013 až 2015 byl vytvořen morfologický klíč, který rozšiřuje klasický systém zobrazení morfologické informace v~islandských slovnících o~přesné zařazení daného hesla do určité deklinační\,/\addthin konjugační skupiny, jejíž
skloňování\,/\addthin časování je ukázáno na příkladu. V~roce 2014 jsme obdrželi roční grant od Ostravské univerzity v~Ostravě na tvorbu morfologického klíče. S~realizací grantu nám velmi pomohla doc. PhDr. Diana Svobodová, Ph.D. a~Ing. Zuzana Havrlantová.

Od března 2014 do srpna 2015 jsme realizovali druhé finální čtení, ve kterém jsme porovnávali informace s~webovým slovníkem \textit{ISLEX} (\cite {int1}). 
Na požádání jsme obdrželi od Þórdís Úlfarsdóttir (vedoucí projektu \textit{ISLEX}) seznam víceslovných příslovečných frází. Takto poskytnutý seznam nám usnadnil opravy slovníku. 
S~Þórdís Úlfarsdóttir jsme konzultovali různé gramatické jevy a~záležitosti týkající se deklinací\,/\addthin konjugací.
Martina Kašparová, studentka islandštiny na Islandské univerzitě, prošla celý slovník, přičemž se soustředila mimo jiné na stylistické a~pravopisné korektury překladů. 
Lucie Brabcová udělala jazykové a~typografické korektury celého slovníku.

Vojtěch Kupča vydal v~roce 2014 \textit{Islandskou gramatiku} (\cite {is77}) a~souhlasil s~připojením této publikace ke~knižní verzi IČSS.
V~roce 2015 jsme s~přispěním Paola al-Imk\=an\=\i\ Brasolin vysázeli slovník v~typografickém prostředí {\LaTeX}. V~únoru 2016 byl slovník vytištěn v~Tiskárně MELMEN, s.r.o v~Pardubicích a svázán 
v~Knihařství Vojtíšková Hrubá v~Pardubicích. 

\blspace[5]

{\centering Všem spoluautorům děkujeme za spolupráci.\par}


% -------------------------------------------------  Guide to the dictionary -

\chapter{Průvodce po slovníku}
\begin{multicols}{2}
\def\tableA{\noindent\begin{tabularx}{\columnwidth}{>{\bfseries}lll}
  \toprule
  \textit{přípona} & \textit{přípona s~koncovkou} & \textit{příklad} \\
  \midrule
  -ing-   & -ingur   & rugl··ingur \\
          & -ing     & hreyf··ing \\
          & -ingi    & erf··ingi \\
          & -ingar   & lækn··ingar \\
  -ling-  & -lingur  & vett··lingur \\
  -ning-  & -ningur  & flut··ningur \\
          & -ning    & set··ning \\
          & -ningar  & kos··ningar \\
  -að-    & -aður    & mark··aður \\
  -nað-   & -naður   & skil··naður \\
  -ar-    & -ari     & dóm··ari \\
  -skap-  & -skapur  & kveð··skapur \\
  -dóm-   & -dómur   & lær··dómur \\
  -leik-  & -leikur  & sann··leikur \\
          & -leiki   & mögu··leiki \\
  -lát-   & -látur   & laus··látur \\
  -átt-   & -átta    & víð··átta \\
  -sam-   & -samur   & starf··samur \\
  -sem-   & -semi    & hamingju··semi \\
  -nesk-  & -neskur  & tékk··neskur \\
          & -neska   & tékk··neska \\
  -leg-   & -legur   & blíð··legur \\
          & -lega    & blíð··lega \\
  -ung-   & -ungur   & gler··ungur \\
          & -unga    & snopp··unga \\
          & -ungi    & sveit··ungi \\
          & -ung     & nýj··ung \\
  -ótt-   & -óttur   & blett··óttur \\
  -ræn-   & -rænn    & suð··rænn \\
          & -ræna    & nor··ræna \\
  -ern-   & -erni    & fað··erni \\
  -ind-   & -indi    & vís··indi \\
  -and-   & -andi    & gef··andi \\
  -ug-    & -ugur    & göf··ugur \\
          & -ugt     & stöð··ugt \\
  -ust-   & -usta    & þjón··usta \\
  -ald-   & -aldur   & far··aldur \\
  -uð-    & -uður    & söfn··uður \\
  -úð-    & -úð      & dul··úð \\
  -versk- & -verskur & ung··verskur \\
          & -verska  & ung··verska \\
  -ist-   & -isti    & hass··isti \\
  -ism-   & -ismi    & húman··ismi \\
  -ang-   & -angur   & far··angur \\
  \bottomrule
  \end{tabularx}
  \captionof{table}{Seznam přípon označených ve slovníku}
  \label{table1}}

\def\tableB{\noindent\begin{tabularx}{\columnwidth}{>{\bfseries}lX}
  \toprule
  \textbf{symbol} & \textbf{vysvětlivka} \\
  \midrule
  m                   & označuje slovní druh \\
  \textsubscript{1}   & označuje třídu \\
  (-a, -ar)           & označuje deklinační / konjugační koncovky hesla \\
  \textsubscript{8}   & označuje podtřídu\\
  \bottomrule
  \end{tabularx}
  \captionof{table}{Morfologická informace}
  \label{table2}}

\def\tableC{\noindent\begin{tabularx}{\columnwidth}
    {>{\bfseries}l>{\footnotesize\itshape}lllll}
  \toprule
  \multirow{3}{*}{{{\textbf{f{\textsubscript{1}}} \Large{\textbf{7}}}}}
    & nom & kisa & kisan    & kisur & kisurnar \\
    & acc & kisu & kisuna   & kisur & kisurnar \\
    & dat & kisu & kisunni  & kisum & kisunum \\
    & gen & kisu & kisunnar & kisa  & kisanna \\
  \bottomrule
  \end{tabularx}
  \captionof{table}{Skloňování hesla \textit{kisa}}
  \label{table3}}

\def\tableD{\noindent\begin{tabularx}{\columnwidth}{>{\bfseries}lllll}
  \toprule
  \textit{zkratka} & \textit{slovo} & \textit{pád} & \textit{číslo}
    & \specialcell{\textit{předmět}\\\textit{podmět}} \\
  \midrule
  e-að  & eitthvað   & 1. & sg & neživotný \\
  e-ð   & eitthvað   & 4. & sg & neživotný \\
  e-ir  & einhverjir & 1. & pl & životný \\
  e-ja  & einhverja  & 4. & pl & životný \\
  e-jum & einhverjum & 3. & pl & životný \\
  e-m   & einhverjum & 3. & sg & životný \\
  e-n   & einhvern   & 4. & sg & životný \\
  e-r   & einhver    & 1. & sg & životný \\
  e-rra & einhverra  & 2. & pl & životný \\
  e-rs  & einhvers   & 2. & sg & životný \\
  e-s   & einhvers  & 2. & sg & neživotný \\
  e-u   & einhverju  & 3. & sg & neživotný \\
  \bottomrule
  \end{tabularx}
  \captionof{table}{Seznam použitých zkratek neurčitých zájmen}
  \label{table4}}

%#############################################################################

\section{Heslo}

\subsection*{Řazení slov}

Heslová slova jsou zobrazena tučně a~jsou seřazena podle islandské abecedy:

\blspace
{\bfseries\centering
  a~á b c d ð e é f g h i~í j k~l m n\\
  o~ó p q r s~t u~ú v~w x y ý z~þ æ ö\par}
\blspace

K~snadnější orientaci v~heslové části je použit u~každého písmene tzv. palcový index, který se nachází na hřbetě slovníku. Každé písmeno je v~palcovém indexu barevně odlišeno.

\subsection*{Homonyma}

Homonyma (heslová slova stejného reprezentativního tvaru) jsou označena povýšenými arabskými číslicemi uváděnými za heslovým slovem.

\blspace
  \dicEntry \dicTerm{vor\textsuperscript{1}} \dicPos{n} \dicFlx{(-s, -)}
  \dicEntry \dicTerm{vor\textsuperscript{2}} \dicPos{pron} \dicFlx{pers gen pl}
\blspace

Homonyma jsou seřazena podle českého ustáleného řazení slovních druhů. Poté následují předpony a~zkratky.
V~případě dvou podstatných jmen je pořadí mužský rod, ženský rod, střední rod. V~případě dvou podstatných jmen stejného rodu má přednost slovo se slabým skloňováním před slovem se silným skloňováním. V~případě dvou sloves stejného reprezentativního tvaru má přednost sloveso se slabým časováním před slovesem se silným časováním.

\subsection*{Varianty a~dublety}

U~některých heslových slov jsou v~záhlaví hesla uvedeny a) varianty (obměny)

\blspace
  \dicEntry \dicTerm{ærsla··fenginn}, \dicTerm{ærslafullur} \dicPos{adj}
\blspace

nebo b) dublety (dvojtvary).

\blspace
  \dicEntry \dicTerm{{ævin··týri}, \dicTerm{ævintýr}} \dicPos{n}
\blspace

Varianty a~dublety jsou méně frekventované než samotné heslo. Lze je vyhledat ve slovníku jako hesla obsahující fonetické a~morfologické informace s~odkazem na původní heslové slovo.

\blspace
  \dicEntry \dicTerm{ærsla··|fullur} \dicIPA{{a}{i}{\textsubring{r}}{s}{\textsubring{d}}{l}{a}{f}{\textscy}{\textsubring{d}}{l}{\textscy}{\textsubring{r}}} \dicPos{adj} \dicFlx{(comp -fyllri, sup -fyllstur)} \dicLink{ærslafenginn}
  \dicEntry \dicTerm{ævin··týr} \dicIPA{{a}{i}{\textlengthmark}{v}{\textsci}{n}{t\textsuperscript{h}}{i}{\textsubring{r}}} \dicPos{n} \dicFlx{(-s, -)} \dicLink{ævintýri}
\blspace

\subsection*{Vybrané tvary jako hesla}

Ve slovníku jsou uvedeny vybrané tvary heslových slov, u~kterých by mohl mít uživatel potíže s~identifikací. Jedná se o~a) genitiv jednotného čísla ({\textbf{sg}}) a~nominativ množného čísla ({\textbf{pl}}) podstatných jmen, b) ženský a~střední rod přídavných jmen a~2. a~3. stupeň rodu mužského přídavných jmen, c) oznamovací způsob ({\textbf{ind}}) přítomného času ({\textbf{praes}}) a~minulého času ({\textbf{pf}}) 1. osoby ({\textbf{pers}}) sg a~pl a~spojovací způsob ({\textbf{con}}) minulého času 1. osoby sg a~jiné nepravidelné tvary.

\blspace
  \dicEntry \dicTerm{öttum} \dicPos{v} \dicFlx{ind pf pl 1 pers} \dicLink{etja}
\blspace

\subsection*{Dělení heslových slov}

K~dělení heslových slov je použita střední tečka \textbf{'·'}, která odděluje části složeného slova, přípony a~předpony. Platí následující pravidla:

\renewcommand{\labelenumi}{\alph{enumi})}
\begin{enumerate}
\item jediné dělení heslového slova je označeno dvěma středními tečkami \textbf{'··'}. Může se jednat o~dělení částí složeného slova
\blspace
\dicEntry \dicTerm{ráð··hús}
\blspace

nebo o~dělení mezi kořenem a~předponou / příponou.

\blspace
\dicEntry \dicTerm{þjón··usta}
\blspace
\item V~případě, že se v~hesle vyskytuje více dělení, jsou použity dvě střední tečky \textbf{'··'} pro významově poslední dělení a~jedna střední tečka \textbf{'·'} pro zbylé / zbylá dělení.
\blspace
\dicEntry \dicTerm{við·bótar··líf·eyris·sparn·aður}
\blspace
\item V~hesle s~předponou je označena jedna nebo i~více předpon (každá zvlášť).
\blspace
\dicEntry \dicTerm{ó··að·finnan·legur}
\blspace
\item V~hesle s~příponou je označena jedna nebo i~více přípon (každá zvlášť).
\blspace
\dicEntry \dicTerm{létt·úð··ugur}
\blspace
\item V~hesle s~příponou je označena pouze přípona / přípony poslední části složeného slova. Důvodem pro neoznačení přípon jiných částí složených slov je lepší čitelnost hesla. Příklad označení pouze přípony poslední části složeniny:
\blspace
\dicEntry \dicTerm{flutninga··miðl·ari}
\end{enumerate}

Následuje seznam přípon, které jsou ve slovníku označeny:

\tableA

\section{Frekvence slov}

Symbolem \dicsymFrequent\ je označeno 2~825 hesel patřících mezi velmi frekventovaná islandská slova. Tento seznam vychází z~korpusu \textit{Íslenskur Orðasjóður} (\cite{int7}) a~knihy \textit{Íslensk orðtíðnibók} (\cite{is2}).

\blspace
  \dicEntry \dicTerm{koma}\textsuperscript{2} \dicsymFrequent
\blspace

\section{Fonetický zápis}

V~hranatých závorkách v~záhlaví hesla je uveden fonetický zápis výslovnosti hesla v~mezinárodní fonetické abecedě (IPA, angl. International Phonetic Alphabet).

\blspace
  \dicEntry \dicTerm{öldu··rót} \dicIPA{{\oe}{l}{\textsubring{d}}{\textscy}{r}{ou}{\textsubring{d}}}
\blspace

Znak '{\textlengthmark}' označuje dlouhou samohlásku nebo souhlásku.

\blspace
  \dicEntry \dicTerm{ör}\textsuperscript{2} \dicIPA{{\oe}{\textlengthmark}{\textsubring{r}}}
\blspace

\section{Slovní druhy, předpony a~zkratky}

V~záhlaví hesla je uvedena morfologická informace popisující a) slovní druh hesla, b) předponu nebo c) zkratku. Pokud se heslo řadí do více slovních druhů, jsou jednotlivé zkratky pro slovní druhy odděleny lomítkem.

\blspace
  \dicEntry \dicTerm{af} \dicPos{prep / adv}
\blspace

Informace o~slovním druhu, předponě či zkratce je vysázena barevně a~tučně.

\subsection*{Podstatná jména}

Podstatná jména jsou uvedena v~1. pádě (\textbf{nom}) jednotného čísla (\textbf{sg}) nebo v~1. pádě množného čísla (\textbf{pl}), pokud se heslo v~jednotném čísle nevyskytuje. Podstatná jména jsou označena následovně: \textbf{m} (l. masculinum) - podstatné jméno rodu mužského,
\textbf{f} (l. femininum) - podstatné jméno rodu ženského, \textbf{n} (l. neutrum) - podstatné jméno rodu středního, \textbf{subs} (l. substantivum) - podstatné jméno bez jednoznačného mluvnického rodu.

\blspace
  \dicEntry \dicTerm{hestur} \dicPos{m}
  \dicEntry \dicTerm{Hanoí} \dicPos{subs}
\blspace

\subsection*{Přídavná jména}

Přídavná jména jsou uvedena v~1. pádě (\textbf{nom}) rodu mužského (\textbf{m}) jednotného čísla (\textbf{sg}) 1. stupně (\textbf{pos}) (l. positivus) a~jsou označena zkratkou \textbf{adj} (l. adjectivum).

\blspace
  \dicEntry \dicTerm{fallegur} \dicPos{adj}
\blspace

V~případě, že u~přídavného jména dochází v~rodě ženském (\textbf{f}) nebo středním (\textbf{n}) k~přehlásce anebo je tvar druhého stupně (\textbf{comp} (l. comparativus)) či třetího stupně (\textbf{sup} (l. superlativus)) nepravidelný, pak jsou tyto tvary uvedeny v~nominativu jako heslová slova s~odkazem na hlavní heslový tvar

\blspace
  \dicEntry \dicTerm{yngri} \dicPos{adj} \dicFlx{comp m} \dicLink{ungur}
\blspace

nebo jako svébytná hesla s~významy.

\blspace
  \dicEntry \dicTerm{fremri} \dicPos{adj} \dicFlx{comp (sup fremstur)}
\blspace

\subsection*{Zájmena}

Zájmena jsou označena zkratkou \textbf{pron} (l. pronomen) a~jsou tříděna na zájmena ukazovací \textbf{dem} (angl. demonstrative), přivlastňovací \textbf{poss} (angl. possessive), osobní \textbf{pers} (angl. personal), neurčitá \textbf{indef} (angl. indefinite), zvratná \textbf{refl} (angl. reflexive) a~vztažná \textbf{rel} (angl. relative).

\blspace
  \dicEntry \dicTerm{þessi} \dicPos{pron} \dicFlx{dem}
\blspace

U~zájmen přivlastňovacích je v~3. osobě označen rod.

\blspace
  \dicEntry \dicTerm{hans} \dicPos{pron} \dicFlx{poss m}
\blspace

U~zájmen osobních je označen pád a~číslo.

\blspace
  \dicEntry \dicTerm{mig} \dicPos{pron} \dicFlx{pers acc sg}
\blspace

\subsection*{Číslovky}

Číslovky jsou označeny zkratkou \textbf{num} (l. numerale). V~případě vybraných číslovek (1 až 4) jsou ve slovníku uvedeny jako heslová slova také tvary ženského i~středního rodu. V~takovém případě jsou označeny zkratkou \textbf{m} pro mužský rod, zkratkou \textbf{f} pro ženský rod a~zkratkou \textbf{n} pro střední rod. Zkratkou \textbf{ord} (angl. ordinal number) jsou označeny číslovky řadové.

\blspace
  \dicEntry \dicTerm{tut··tugu} \dicPos{num}
  \dicEntry \dicTerm{tveir} \dicPos{num} \dicFlx{pl m}
  \dicEntry \dicTerm{fyrsti} \dicPos{num} \dicFlx{ord}
\blspace

\subsection*{Slovesa}

Slovesa jsou uvedena v~infinitivu a~jsou označena zkratkou \textbf{v} (l. verbum).

\blspace
  \dicEntry \dicTerm{fara} \dicPos{v}
\blspace

Slovesa ve středním rodě jsou označena jako \textbf{refl} (angl. reflexive).

\blspace
  \dicEntry \dicTerm{nálgast} \dicPos{v} \dicFlx{refl}
\blspace

Slovesa, která jsou vždy neosobní, jsou označena \textbf{impers} (angl. impersonal).

\blspace
  \dicEntry \dicTerm{svima} \dicPos{v} \dicFlx{impers}
\blspace

U~sloves je dále uvedeno, s~jakým pádem se pojí. Zkratka \textbf{acc} označuje 4. pád, \textbf{dat} 3. pád, \textbf{gen} 2. pád a~\textbf{nom} 1. pád. V~případě, že se sloveso pojí s~různými pády, je použito lomítko.

\blspace
  \dicEntry \dicTerm{klóra} \dicPos{v} \dicFlx{acc / dat}
\blspace

V~případě, že se sloveso pojí s~více pády, je použito znaménko plus.

\blspace
  \dicEntry \dicTerm{gefa} \dicPos{v} \dicFlx{dat + acc}
\blspace

\subsection*{Příslovce}

Příslovce jsou označena zkratkou \textbf{adv} (l. adverbium). Pokud chybí informace o~stupni, jedná se o~1. stupeň stupňování (\textbf{pos}). 2. stupeň je vždy označen zkratkou \textbf{comp} a~3. stupeň zkratkou \textbf{sup}.

\blspace
  \dicEntry \dicTerm{þar} \dicPos{adv}
  \dicEntry \dicTerm{ofar} \dicPos{adv} \dicFlx{comp (pos uppi, sup efst)}
  \dicEntry \dicTerm{síðast} \dicPos{adv} \dicFlx{sup (comp síðar)}
\blspace

\subsection*{Předložky}

Předložky jsou označeny zkratkou \textbf{prep} (l. praepositio). Pokud se předložka pojí výhradně s~jedním pádem, jsou použity zkratky pro 4. pád (\textbf{acc}), 3. pád (\textbf{dat}), 2. pád (\textbf{gen}).

\blspace
  \dicEntry \dicTerm{af} \dicPos{prep / adv} \dicFlx{dat}
\blspace

V~případě, že se předložka pojí s~různými pády, je použito lomítko.

\blspace
  \dicEntry \dicTerm{fyrir} \dicPos{prep / adv} \dicFlx{acc / dat}
\blspace

\subsection*{Spojky}

Spojky jsou označeny zkratkou \textbf{conj} (l. conjunctio).

\subsection*{Částice}

Částice jsou označeny zkratkou \textbf{part} (l. particula).

\subsection*{Citoslovce}

Citoslovce jsou označeny zkratkou \textbf{inter} (l. interjectio).

\subsection*{Předpony}

Předpony jsou označeny zkratkou \textbf{predp}.

\subsection*{Zkratky}

Zkratky jsou označeny zkratkou \textbf{zkr}.

\section{Skloňování, časování a~stupňování}

Každé heslo obsahuje zápis morfologické informace, která přesně informuje, jak se heslo skloňuje / časuje nebo stupňuje.

V~slovníku jsou použity dva odlišné systémy zápisu morfologické informace, které se vzájemně doplňují:

\begin{enumerate}
\item tradiční systém - do závorky za slovním druhem se zapisují deklinační / konjugační koncovky. Počet koncovek závisí na slovním druhu (viz dále).
\item numerický systém (vychází ze systému Helgi Haraldssona (\cite {is6})) - číslice odkazující na morfologický klíč, kde jsou podrobně vypsány deklinační / konjugační třídy.
Další informace o~numerickém systému lze nalézt na straně \pageref{sec:dec}.
\end{enumerate}

Deklinační / konjugační tvary nebo koncovky podstatných jmen, přídavných jmen, zájmen, číslovek, sloves a~příslovcí jsou uvedeny v~závorkách za slovním druhem.

Hnízdovací šev \textbf{'|'} označuje místo ve slově, kde se deklinační koncovka připojuje ke slovu.

\blspace
  \dicEntry \dicTerm{heils|a{\textsuperscript{1}}} \dicPos{f} \dicFlx{(-u)}
\blspace

Pokud není hnízdovací šev \textbf{'|'} uveden, tak se koncovka připojuje k~samotnému heslu.

\blspace
  \dicEntry \dicTerm{mynd} \dicPos{f} \dicFlx{(-ar, -ir)}
\blspace

Pokud dochází v~ohýbání k~přehlásce nebo je ohýbání nepravidelné, pak jsou uvedeny celé tvary hesla.

\blspace
  \dicEntry \dicTerm{maður} \dicPos{m} \dicFlx{(manns, menn)}
\blspace

V~případě, že heslové slovo má více tvarů, jsou varianty koncovky odděleny lomítkem.

\blspace
  \dicEntry \dicTerm{beð|ur} \dicPos{m} \dicFlx{(-s / -jar, -ir)}
\blspace

\subsection*{Podstatná jména}

U~podstatných jmen je uveden tvar hesla pro 2.~pád jednotného čísla a~1.~pád množného čísla.

\blspace
  \dicEntry \dicTerm{hval|ur} \dicPos{m} \dicFlx{(-s, -ir)}
\blspace

kde '\textit{hvals}' je 2.~pád jednotného čísla a~'\textit{hvalir}' 1.~pád množného čísla.
Pokud je uveden pouze jeden tvar, jedná se o~2.~pád jednotného čísla a~znamená to rovněž, že heslové slovo se neohýbá v~množném čísle.

\blspace
  \dicEntry \dicTerm{heils|a{\textsuperscript{1}}} \dicPos{f} \dicFlx{(-u)}
\blspace

Pokud je použita zkratka pro množné číslo (\textbf{pl}), znamená to, že se heslo ohýbá pouze v~množném čísle.

\blspace
  \dicEntry \dicTerm{afar··kostir} \dicPos{m} \dicFlx{pl}
\blspace

Pokud je použita zkratka \textbf{indecl}, znamená to, že podstatné jméno je nesklonné.

\blspace
  \dicEntry \dicTerm{október} \dicPos{m} \dicFlx{indecl}
\blspace

\subsection*{Přídavná jména}

U~přídavných jmen jsou uvedeny v~závorkách tvary, které poukazují na nepravidelnosti ve skloňování.

Pokud přídavné jméno v~ženském rodě neztrácí koncové '\textit{-r}', je uveden tvar hesla v~ženském rodě.

\blspace
  \dicEntry \dicTerm{hýr} \dicPos{adj} \dicFlx{(f -)}
\blspace

Dále jsou uvedeny tvary ženského rodu, ve~kterých dochází k~přehlásce:

\blspace
  \dicEntry \dicTerm{reglu··|samur} \dicPos{adj} \dicFlx{(f -söm)}
\blspace

Pokud je přídavné jméno nesklonné, je uvedena zkratka \textbf{indecl}.

\blspace
  \dicEntry \dicTerm{tví··mála} \dicPos{adj} \dicFlx{indecl}
\blspace

Pokud má přídavné jméno nepravidelné tvary 2.~a~3.~stupně stupňování, jsou tyto tvary uvedeny v~závorce spolu se zkratkou pro 2.~stupeň (\textbf{comp}) a~3.~stupeň (\textbf{sup}).

\blspace
  \dicEntry \dicTerm{þunnur} \dicPos{adj} \dicFlx{(comp þynnri, sup þynnstur)}
\blspace

V~případě, že je uveden pouze jeden ze dvou zbývajících stupňů, znamená to, že heslo se stupňuje pouze v~uvedených stupních.

\blspace
  \dicEntry \dicTerm{syðri} \dicPos{adj} \dicFlx{comp (sup syðstur)}
\blspace

\subsection*{Slovesa}

Islandská slovesa se dělí na slabá, silná a~nepravidelná. Vzhledem k~tomuto dělení se liší počet slovních tvarů uvedených v~závorce za zkratkou \textbf{v}. Princip uvádění slovních tvarů spočívá v~tom, že uživatel
je schopen (více méně) z~těchto tvarů odvodit časování slovesa. V~tomto slovníku jsme
záměrně umístili některá slabá slovesa mezi silná (místo dvou slovesných tvarů je uvedeno pět tvarů) kvůli tomu, aby byl uživatel přesněji informován o~časování slovesa.

Slabá slovesa, která tvoří 1.~osobu jednotného čísla minulého času oznamovacího způsobu koncovkou '\textit{-aði}', mají v~závorce uvedenu jen jednu koncovku, jmenovitě '\textit{(-aði)}'.

\blspace
  \dicEntry \dicTerm{ætl|a} \dicPos{v} \dicFlx{(-aði)}
\blspace

U~zbývajících skupin slabých sloves se v~závorce nachází koncovka 1.~osoby jednotného čísla minulého času oznamovacího způsobu a~supinum.

\blspace
  \dicEntry \dicTerm{kenn|a} \dicPos{v} \dicFlx{(-di, -t)}
\blspace

Silná nebo nepravidelná slovesa mají v~závorce vždy pět tvarů, jmenovitě 1.~osobu jednotného čísla přítomného času oznamovacího způsobu, 1.~osobu jednotného čísla minulého času oznamovacího způsobu, 1.~osobu množného čísla minulého času oznamovacího způsobu, 1.~osobu jednotného čísla minulého času spojovacího způsobu a~supinum.

\blspace
  \dicEntry \dicTerm{grípa} \dicPos{v} \dicFlx{(gríp, greip, gripum, gripi, gripið)}
\blspace

Silná nebo nepravidelná slovesa vyskytující se pouze v~neosobním užití mají v~závorce 4 tvary, jmenovitě 3.~osobu jednotného čísla přítomného času oznamovacího způsobu, 3.~osobu jednotného čísla minulého času oznamovacího způsobu, 3.~osobu jednotného čísla minulého času spojovacího způsobu a~supinum.

\blspace
  \dicEntry \dicTerm{kross··|bregða} \dicPos{v} \dicFlx{(-bregður, -brá, -brygði, -brugðið)}
\blspace

U~některých významů v~heslovém slově se vyskytují gramatické informace, které popisují chování daného významu.

Zkratka \textbf{impers} označuje, že se daný význam slovesa vyskytuje v~neosobním užití. Pomocí slovního spojení je vyjádřeno užití slovesa.

\blspace
  \dicEntry \dicTerm{batn|a} \dicPos{v} \dicFlx{(-aði)} \dicTerm{e-m batnar} \dicFlx{impers} (kdo) se uzdravuje, (kdo) se zotavuje, (komu) je lépe
\blspace

Zkratka \textbf{refl} označuje mediopasivní tvary slovesa.

\blspace
  \dicEntry \dicTerm{að··lag|a} \dicPos{v} \dicFlx{(-aði)} \dicIdiom{aðlagast} {\dicPhraseIS{aðlagast e-u}} \dicFlx{refl} {přizpůsobit se (čemu)}
\blspace

\subsection*{Příslovce}

U~příslovcí se v~záhlaví nachází tvary stupňování. Tyto tvary jsou uvedeny v~závorce spolu se zkratkou pro 2.~stupeň (\textbf{comp}) a~3.~stupeň (\textbf{sup}).

\blspace
  \dicEntry \dicTerm{vel\textsuperscript{2}} \dicPos{adv} \dicFlx{(comp betur, sup best)}
\blspace

Nepravidelné tvary 2.~a~3.~stupně jsou uvedeny buď jako hesla odkazující na základní tvar v~1.~stupni (\textbf{pos}), anebo jako svébytná hesla. V~obou případech jsou uvedeny v~závorce zbývající stupně.

\blspace
  \dicEntry \dicTerm{síst} \dicPos{adv} \dicFlx{sup (pos varla, comp síður)}
\blspace

Pokud některý ze stupňů uveden není, znamená to, že pro dané slovo neexistuje.

\blspace
  \dicEntry \dicTerm{síðar} \dicPos{adv} \dicFlx{comp (sup síðast)}
\blspace

\section{Morfologická informace} \label{sec:dec}

V~záhlaví hesla se kromě deklinačních / konjugačních koncovek nachází také morfologické numerické informace, které odkazují na morfologický klíč (viz strana \pageref{sec:morpho}), kde jsou podrobně vypsány kompletní deklinační / konjugační tabulky.

\blspace
  \dicEntry \dicTerm{ap|i} \dicPos{m}[1] \dicFlx{(-a, -ar)}[8]
\blspace

Morfologická informace je zapsána podle následujícího schématu: \textbf{m{\textsubscript{1}} (-a, -ar){\textsubscript{8}}}

\blspace
\tableB
\blspace

Deklinační / konjugační třídy a~jejich počet nejsou v~islandštině ustálené a~slouží pouze pro rozlišení odlišných skloňování / časování.

Názvy podtříd slouží pro rozlišení odlišných skloňování / časování.

\subsection*{Práce s~morfologickým klíčem}

Pro vyhledání kompletních deklinačních / konjugačních tabulek slouží následující postup, který demonstrujeme na ukázce:

\blspace
  \dicEntry \dicTerm{kis|a} \dicPos{f}[1] \dicFlx{(-u, -ur)}[7]
\blspace

Nejdříve určíme slovní druh (v~případě podstatných jmen také rod), v~tomto případě podstatné jméno rodu ženského, a~navštívíme Morfologický klíč (str. \pageref{sec:morpho}) pro podstatná jména rodu ženského (str. \pageref{sec:morpho_f}).
Následně dohledáme třídu (v~tomto případě číslo {\textsubscript{\textbf{1}}}) a~následně dohledáme dodatečné numerické informace za závorkou (v~tomto případě číslo {\textsubscript{\textbf{7}}}).

\tableC

\section{Definice}

\subsection*{Základní tvar}

Základním tvarem definice je, že jednomu islandskému slovu odpovídá jeden nebo více českých významů.

\blspace
  \dicEntry \dicTerm{landa··fræð|i} \dicPos{f} \dicFlx{(-i)} \textCS{zeměpis, geografie}
\blspace

V~některých případech je v~české části použita závorka. V~prvním případě je závorka užívaná ke~zkrácení zápisu dvou slov ('\textit{pomoc}', '\textit{výpomoc}').

\blspace
  \dicEntry \dicTerm{að··stoð} \dicPos{f} \dicFlx{(-ar)} \textCS{(vý)pomoc, asistence}
\blspace

V~druhém případě se v~závorce nachází upřesňující informace, která zabraňuje dvojznačnosti.

\blspace
  \dicEntry \dicTerm{fíkj|a} \dicPos{f} \dicFlx{(-u, -ur)} \textCS{fík} {\footnotesize {\textCS{(plod)}}}
\blspace

V~mnoha případech je závorka v~české části použita pro příklad užití českého slova a~tím vymezení jeho významu.

\blspace
  \dicEntry \dicTerm{að·gengi··legur} \dicPos{adj} {přístupný, dostupný (informace ap.)}
\blspace

Lomítko \textbf{'/'} slouží v~české části významu k~oddělení více významů pojících se s~jedním nebo více slovy a~je použito v~české části následujícím způsobem:

\blspace
  \dicEntry \dicTerm{aðal··|gata} \dicPos{f} \dicFlx{(-götu, -götur)} \textCS{hlavní třída / ulice}
\blspace

Význam přečteme jako '\textit{hlavní třída, hlavní ulice}'.

\subsection*{Opis heslového slova}

V~případech, kdy islandské slovo nemá ekvivalent v~českém jazyce, jsou významy islandských slov opsány delším opisem a~zobrazeny menším písmem. Příkladem jsou slova z~lidových pověstí, kulinářské speciality nebo botanické pojmy.

\blspace
  \dicEntry \dicTerm{til··ber|i} \dicPos{m} \dicFlx{(-a, -ar)} {\footnotesize{pov.}} {\footnotesize{stvoření, které dojí krávy a~ovce hospodářů}}
\blspace

\subsection*{Slovní spojení}

V~definici hesla jsou kromě samotných významů uváděna také slovní spojení, která jsou sázena tučně. Po slovním spojení následuje překlad slovního spojení.

\blspace
  \dicEntry \dicTerm{lær|a} \dicPos{v} \dicFlx{(-ði, -t)} \textbf{læra e-ð utanbókar} {učit se (co) zpaměti, memorovat (co)}
\blspace

Závorky jsou použity na tu část slovního spojení, kterou je možné vypustit.

\blspace
  \dicEntry \dicTerm{ár\textsuperscript{1}} \dicPos{f} \dicFlx{(-ar, -ar)} \textbf{taka (of) djúpt í árinni} {\footnotesize{přen.}} {přehnat (co v~tvrzení)}
\blspace

Příklad můžeme přečíst jako '\textit{taka djúpt í árinni}' nebo '\textit{taka of djúpt í árinni}'.

Mezi slovní spojení v~tomto slovníku řadíme a) tvary heslového slova, b) příkladová spojení, c) ustálená slovní spojení a~d) rčení nebo přísloví.

\begin{enumerate}
\item Tvary heslového slova\\
Do této skupiny řadíme tvar množného čísla hesla (zkratka \textbf{pl}).
\blspace
\dicEntry \dicTerm{bót} \dicPos{f} \dicFlx{(-ar, bætur)} \textbf{bætur} \dicFlx{pl} {dávky, přídavky, příspěvky}\\

Dále do této skupiny řadíme tvar slovesa v~středním slovesném rodě (zkratka \textbf{refl}).\\

\dicEntry \dicTerm{finna} \dicPos{v} \dicFlx{(finn, fann, fundum, fyndi, fundið)} \dicIdiom{finnast} {\dicPhraseIS{finnast}} \dicFlx{refl} {být nalezen, najít se, nalézt se, nacházet se}
\blspace
\item Příkladová spojení\\
Příkladová spojení ilustrují komunikační chování heslových slov v~příslušných významech a~v~různých syntagmatech (např. sloveso s~určitou předložkou ap.). Jsou označena tučně modrou barvou a~odsazená pro snadnou orientaci jako svébytné heslo.
\blspace
\dicEntry \dicTerm{koma\textsuperscript{2}} \dicPos{v} \dicFlx{(kem, kom, komum, kæmi, komið)} \dicIdiom{koma}[að] {\dicPhraseIS{koma að e-m}} {přiblížit se ke (komu), přijít ke (komu)}
\blspace
\item Ustálená slovní spojení\\
Tento druh slovních spojení je charakterizován většinou přeneseným významem. Tato spojení jsou uváděna za posledním významem nebo příkladovým spojením a~jsou označená tučně modrou barvou a~odsazená pro snadnou orientaci jako svébytné heslo.
\blspace
\dicEntry \dicTerm{kast} \dicPos{n} \dicFlx{(-s, köst)} \dicIdiom{kast} {\dicPhraseIS{þegar til kastanna kemur}} {\footnotesize{přen.}} {když jde do tuhého}
\blspace
\item Rčení a~přísloví\\
Rčení a~přísloví se nacházejí vždy na konci definice a~jsou odlišená symbolem \dicsymProverb\ a~odsazená pro snadnou orientaci.
\blspace
\dicEntry \dicTerm{æfing} \dicPos{f} \dicFlx{(-ar, -ar)} \dicProverb{} \dicPhraseIS{Æfing skapar meistara.} {\footnotesize{přís.}} \textCS{Cvičení dělá mistra.}
\end{enumerate}

\section{Vícevýznamovost}

\subsection*{Vícevýznamová hesla}

Značné množství hesel má více než jeden význam. Významy jsou odlišeny arabskými číslicemi a~v~některých případech i~islandskými synonymy.

\blspace
  \dicEntry \dicTerm{land} \dicPos{n} \dicFlx{(-s, lönd)} {\textbf{1.}} \dicSynonym{þurrlendi} {souš, pevnina, země} {\textbf{2.}} \dicSynonym{árbakki} {břeh} {\textbf{3.}} \dicSynonym{ríki} {země, stát, vlast} {\textbf{4.}} \dicSynonym{landareign} {země, půda, pozemek}
\blspace

\subsection*{Vícevýznamová slovní spojení}

V~případě, že slovní spojení uvnitř definice má více významů, jsou jednotlivé významy odlišeny písmeny {\textbf{a., b., c.}} atd.

\blspace
  \dicEntry \dicTerm{drag} {\dicPos{n} \dicFlx{(-s, drög)} {\textbf{1.}} {mokřina, podmoklý terén;} {\textbf{2.}} \textbf{drög} {\footnotesize{pl}} {\textbf{a.}} \dicSynonym{uppsprettur} {prameny (řeky ap.)} {\textbf{b.}} \dicSynonym{undirbúningur} {náčrt, návrh.}}
\blspace

\section{Řazení významů a~slovních spojení v~definici}
\longsectionskip

V~definici jsou nejdříve seřazeny významy slova a~odlišeny arabskými číslicemi.

\subsection*{V rámci významu}

Slovní spojení jsou uváděna v~rámci jednotlivého významu, pokud se s~ním významově vážou.

\blspace
  \dicEntry \dicTerm{tím|i} \dicPos{m} \dicFlx{(-a, -ar)} {\textbf{1.}} \dicSynonym{tíð} {čas, doba}; \textbf{í þann tíma} \dicFlx{adv} \dicSynonym{þá\textsuperscript{2}} {pak} {\textbf{2.}} \dicSynonym{rétt stund} {(správný) čas / moment}
\blspace

\subsection*{Za posledním očíslovaným významem}

Pokud se slovní spojení neváže k~žádnému významu, je řazeno za posledním očíslovaným významem. Ustálená slovní spojení jsou uvedena symbolem \dicsymIdiom.

\blspace
  \dicEntry \dicTerm{finna} \dicPos{v} \dicFlx{(finn, fann, fundum, fyndi, fundið)} {\textbf{3.}} \dicSynonym{skynja} {cítit, vnímat} \dicIdiom{finna}[að] {\dicPhraseIS{finna að e-u}} {kritizovat (co), nacházet chyby na (čem)}
\blspace

\subsection*{Způsob řazení slovních spojení}

Slovní spojení, která se nachází za posledním očíslovaným významem, jsou řazena podle předložky / příslovce / zájmena, se kterou se pojí. Pro snadnou orientaci v~rozsáhlých definicích je uvedeno heslové slovo nebo tvar heslového slova plus předložka / příslovce / zájmeno a~barevně odlišeno modrou barvou.

\blspace
  \dicEntry \dicTerm{koma\textsuperscript{2}} \dicPos{v} \dicFlx{(kem, kom, komum, kæmi, komið)} \dicIdiom{koma}[út] {\dicPhraseIS{koma út}} {vyjít (kniha ap.), objevit se (na trhu ap.)}
\blspace

Poté následují slovní spojení v~slovesném středním rodě řazená opět podle předložky / příslovce / zájmena, se kterou se pojí:

\blspace
  \dicEntry \dicTerm{koma\textsuperscript{2}} \dicPos{v} \dicFlx{(kem, kom, komum, kæmi, komið)} \dicIdiom{komast} {\dicPhraseIS{komast}} \dicFlx{refl} {dostat se, dostávat se}
\blspace

Poté následují ustálená slovní spojení, která nebylo možné zařadit do předchozích položek:

\blspace
  \dicEntry \dicTerm{koma\textsuperscript{2}} \dicPos{v} \dicFlx{(kem, kom, komum, kæmi, komið)} \dicIdiom{koma} {\dicPhraseIS{(komdu) sæll (og blessaður)}} {ahoj, nazdar, dobrý den}
\blspace

Nakonec jsou uvedena přísloví a~rčení označená symbolem \dicsymProverb\:

\blspace
  \dicEntry \dicTerm{eik} \dicPos{n} \dicFlx{(-ar / -ur, -ur)} \dicProverb{} \dicPhraseIS{Ekki / eigi fellur eik við fyrsta högg.} {\footnotesize{přís.}} {Jednou ranou dub nespadne.}
\blspace

\section{Synonyma a~antonyma / hesla ke~srovnání}
\longsectionskip

\subsection*{Synonyma}

Synonyma nebo synonymní slovní spojení jsou uvedeny před opisem významu hesla. Jsou umístěna v~závorkách. Islandská synonyma slouží a) islandskému uživateli k~rozlišení českých významů b) českému uživateli jako dodatečná informace o~významu slova.

\blspace
  \dicEntry \dicTerm{staða} \dicPos{f} \dicFlx{(stöðu, stöður)} {\textbf{1.}} \dicSynonym{ástand} {situace, stav, poměry} {\textbf{2.}} \dicSynonym{embætti} {post, (pracovní) místo / pozice (ve firmě ap.)}
\blspace

V~případě, že synonymum lze nalézt jako heslo ve slovníku, je použita šipka '\dicsymSee\ '.

\blspace
  \dicEntry \dicTerm{bygging} {\dicPos{f} \dicFlx{(-ar, -ar)} {\textbf{1.}} \dicSynonym{hús} {budova, stavba, stavení} {\textbf{2.}} \textit{(það að byggja)} {(vý)stavba} }
\blspace

\subsection*{Antonyma / Hesla ke~srovnání}

Antonyma či hesla ke~srovnání jsou uvedena za opisem významu heslového slova a~jsou označena šipkou '$\shortuparrow$'.

\blspace
  \dicEntry \dicTerm{kaldur} \dicPos{adj} \dicFlx{(f köld)} {\textbf{1.}} {studený, chladný} \textit{($\shortuparrow$ heitur)}
\blspace

\section{Oborové a~stylové charakteristiky}

K~upřesnění významu heslového slova slouží oborové a~stylové charakteristiky.

\subsection*{Oborové charakteristiky}

Oborové charakteristiky zařazují islandské slovo do určitého oboru a~naznačují použití daného slova v~jazyce

\blspace
  \dicEntry \dicTerm{berg} \dicPos{n} \dicFlx{(-s, -)} {\textbf{1.}} {\footnotesize{geol.}} {hornina, kámen} {\textbf{2.}} \dicSynonym{klöpp\textsuperscript{1}} {skalní stěna, skalisko}
\blspace

a~vymezují význam českého slova.

\blspace
  \dicEntry \dicTerm{mús} \dicPos{f} \dicFlx{(-ar, mýs)} {\textbf{1.}} {\footnotesize{zool.}} {myš} (\textit{l. Mus}) {\textbf{2.}} {\footnotesize{poč.}} {myš}
\blspace

\subsection*{Stylové charakteristiky}

Stylové charakteristiky jsou uváděny v~případě, kdy význam heslového slova je možné použít pouze při určité příležitosti nebo je význam jazykově zabarven.

\blspace
  \dicEntry \dicTerm{gaur} \dicPos{m} \dicFlx{(-s, -ar)} {\textbf{1.}} {\footnotesize{hovor.}} {chlápek, maník, člověk}
\blspace

\section{Syntax}

Syntaktické informace popisují:

\begin{enumerate}
\item zdali je sloveso tranzitivní nebo intranzitivní a~s~jakým(i) pádem (pády) se pojí, pokud je tranzitivní,
\item zdali je podmět / předmět životný či neživotný,
\item zdali je sloveso osobní či neosobní a~v~jakém pádu je podmět, pokud je sloveso neosobní,
\item zdali je podmět gramatický,
\item zdali je sloveso použito ve středním rodě,
\item nejčastější příklady, které se pojí se slovesem,
\item se kterými příslovci nebo předložkami se sloveso pojí.
\end{enumerate}

Syntaktické informace se ve slovníku vyskytují ve třech formách: a) informace uvedené zkratkou, b) teoretické příklady a~c) praktické příklady.

\subsection*{Syntaktické informace uvedené zkratkou}

Zde máme na mysli informace o~pádě (\textbf{nom, acc, dat, gen}), slovesu v~středním slovesném rodě (\textbf{refl}), neosobním slovesu (\textbf{impers}), čísle (\textbf{sg, pl}), rozkazovacím způsobu (\textbf{imper}) a~spojovacím způsobu (\textbf{con}).

\subsection*{Teoretické příklady}

Teoretickým příkladem nazýváme příklad, kdy podmět nebo předmět je zastoupen zájmenem neurčitým v~pádě, se~kterým se slovo pojí a~sloveso je zpravidla v~infinitivu. V~případě, že podmět není v~1.~pádě (např. u~sloves neosobních), je sloveso uvedeno v~3.~osobě jednotného čísla přítomného času.
Ve slovních spojeních používáme zkratky pro tvary neurčitého zájmena \textbf{einhver} (někdo) a~\textbf{eitthvað} (něco). Tyto zkratky označují a) pád, b) číslo a~c) životnost nebo neživotnost podmětu / předmětu.\\
Následuje seznam zkratek neurčitých zájmen použitých ve slovníku: \\

\tableD
\blspace
Příklad použití zkratek s~neurčitými zájmeny:

\blspace
  \dicEntry \dicTerm{leit|a} \dicPos{v} \dicFlx{(-aði)} \textbf{leita að e-u / e-m} {hledat (co / koho)} \dicExampleIS{leita að lyklunum} \dicExampleCS{hledat klíče}
\blspace

Příklad tedy přečteme jako '\textit{leita að einhverju}' (\textit{hledat co}) a~'\textit{leita að einhverjum}' (\textit{hledat koho}). V~tomto případě se sloveso pojí s~3.~pádem a~předmět může být životný i~neživotný.

\subsection*{Praktické příklady}

Ve slovníku se vyskytuje velké množství praktických příkladů a~jejich překladů. Praktický příklad se liší od příkladu teoretického tím, že místo zájmen neurčitých jsou uvedena frekventovaná slova z~běžného projevu. Praktické příklady slouží k~ilustraci použití
heslového slova.\\
Praktickým příkladem může být a) (složené) slovo, b) slovní spojení nebo c) věta.

\begin{enumerate}
\item Slovo jako praktický příklad:
\blspace
\dicEntry \dicTerm{hett|a} \dicPos{f} \dicFlx{(-u, -ur)} \textbf{1.} {kapuce, kápě} \dicExampleIS{hettuúlpa} \dicExampleCS{bunda s~kapucí}
\blspace
\item Slovní spojení jako praktický příklad:
\blspace
\dicEntry \dicTerm{tím|i} \dicPos{m} \dicFlx{(-a, -ar)} \textbf{3.} \dicSynonym{klukkustund} {hodina} {\footnotesize{(šedesát minut)}} \dicExampleIS{tveggja tíma gangur} \dicExampleCS{dvouhodinový pochod}
\blspace
\item Věta jako praktický příklad:
\blspace
\dicEntry \dicTerm{mæt|a} \dicPos{v} \dicFlx{(-ti, -t)} \dicIdiom{mætast} \dicPhraseIS{mætast} \dicSynonym*{ná saman} \dicFlx{refl} {setkat se, potkat se} \dicExampleIS{Þau mættust á gangstéttinni.} \dicExampleCS{Potkali se na chodníku.}
\end{enumerate}

\section{Fotografie a~ilustrace}

Ve slovníku je publikováno 630 fotografií rostlin, živočichů a~fotografií, která dokumentují místopisně heslová slova (např. kulinářské speciality).
Fotografie jsou publikovány pod veřejnými licencemi. Jména autorů a~názvy licencí lze nalézt v~Seznamu autorů fotografií a~ilustrací na straně \pageref{sec:photo}. U~hesla a~konkrétně u~daného významu, ke kterému se fotografie vztahuje, je použit symbol fotoaparátu \dicsymPhoto. Fotografie je umístěna buď pod heslem, anebo v~horní části stránky, kde se heslo nachází, anebo na stránce (stránkách) následující(ch).

%\par\begin{center}\setlength\fboxsep{0pt}\setlength\fboxrule{0.5pt}
%\fbox{\includegraphics[width=6cm]{ds_image_posthus_0_2.jpg}}\end{center}
%\par\begin{center}\footnotesize {Pósthús}\end{center}

\end{multicols}

% ---------------------------------------- List of abbreviations and symbols -

\chapter{Seznam použitých zkratek a symbolů}
\begin{multicols}{2}
\begin{description}[leftmargin=!,labelwidth=1.5cm]
%[\breaklabel\setleftmargin{40pt}\setlabelstyle{\bfseries}
%\addtolength{\itemsep}{-0.5\baselineskip}]
\item[{abb}] {zkrácený}
\item[{acc}] {4. pád, akuzativ}
\item[{adj}] {přídavné jméno}
\item[{adv}] {příslovce}
\item[{akt}] {rod činný, aktivum}
\item[{anat.}] {anatomie}
\item[{angl.}] {anglicky}
\item[{astro.}] {astronomie, výzkum vesmíru}
\item[{ap.}] {a podobně}
\item[{básn.}] {básnický, poetický výraz}
\item[{biol.}] {biologie}
\item[{bot.}] {botanika}

\item[{comp}] {2. stupeň stupňování, komparativ}
\item[{con}] {spojovací způsob, konjunktiv}
\item[{conj}] {spojka}
\item[{dat}] {3. pád, dativ}
\item[{def}] {určitý tvar}
\item[{dem}] {zájmeno ukazovací}
\item[{dět.}] {dětsky}

\item[{e-að}] {co (sg nom neživotný) (\textit{eitthvað})}
\item[{e-ð}] {co (sg acc neživotný) (\textit{eitthvað})}
\item[{e-ir}] {kdo (pl nom životný) (\textit{einhverjir})}
\item[{e-ja}] {koho (pl acc životný) (\textit{einhverja})}
\item[{e-jum}] {komu (pl dat životný) (\textit{einhverjum})}
\item[{e-m}] {komu (sg dat životný) (\textit{einhverjum})}
\item[{e-n}] {koho (sg acc životný) (\textit{einhvern})}
\item[{e-r}] {kdo (sg nom životný) (\textit{einhver})}
\item[{e-rra}] {koho (pl gen životný) (\textit{einhverra})}
\item[{e-rs}] {koho (sg gen životný) (\textit{einhvers})}
\item[{e-s}] {čeho (sg gen neživotný) (\textit{einhvers})}
\item[{e-u}] {čemu (sg dat neživotný) (\textit{einhverju})}

\item[{ekon.}] {ekonomika, obchod}
\item[{elek.}] {elektřina}
\item[{f}] {rod ženský}
\item[{filos.}] {filosofie, logika}
\item[{form.}] {formálně}
\item[{fyz.}] {fyzika}
\item[{gen}] {2. pád, genitiv}
\item[{geog.}] {zeměpis, geografie}
\item[{geol.}] {geologie}
\item[{han.}] {hanlivě, pejorativně}
\item[{hist.}] {historie}
\item[{hovor.}] {hovorově}
\item[{hrub.}] {hrubě}
\item[{hud.}] {hudba}
\item[{imper}] {rozkazovací způsob, imperativ}
\item[{impers}] {sloveso neosobní}
\item[{ind}] {oznamovací způsob, indikativ}
\item[{indecl}] {nesklonný}
\item[{indef}] {neurčitý tvar, (pron +) zájmeno neurčité}
\item[{inf}] {infinitiv}
\item[{int}] {tázací zájmeno}
\item[{inter}] {citoslovce}
\item[{chem.}] {chemie}
\item[{jaz.}] {jazykověda}
\item[{kulin.}] {kulinářství, vaření}
\item[{l.}] {latinsky}
\item[{let.}] {letectví}
\item[{lit.}] {literatura, vydavatelství}
\item[{m}] {rod mužský}
\item[{mat.}] {matematika}
\item[{med. }] {lékařství}
\item[{meteo.}] {meteorologie}
\item[{myt.}] {mytologie}
\item[{n}] {rod střední}
\item[{náb.}] {náboženství}
\item[{nám.}] {námořnictví, rybolov}
\item[{nom}] {1. pád, nominativ}
\item[{num}] {číslovka}
\item[{ord}] {řadová číslovka}
\item[{p}] {osoba}
\item[{part}] {částice}
\item[{pers}] {(v +) osoba, (pron +) zájmeno osobní}
\item[{pf}] {minulý čas}
\item[{pl}] {množné číslo, plurál}
\item[{poč.}] {informatika}
\item[{pol.}] {politika, politologie}
\item[{poss}] {zájmeno přivlastňovací}
\item[{pos}] {1. stupeň stupňování, pozitiv}
\item[{pov.}] {lidové pověsti, folkloristika}
\item[{pp}] {příčestí minulé}
\item[{praes}] {přítomný čas}
\item[{práv.}] {právnictví, soudnictví}
\item[{predp}] {předpona}
\item[{prep}] {předložka}
\item[{presp}] {příčestí přítomné}
\item[{pron}] {zájmeno}
\item[{prop}] {vlastní jméno, proprium}
\item[{přen.}] {přeneseně}
\item[{přís.}] {přísloví}
\item[{psych.}] {psychologie}
\item[{refl}] {(v +) médium, (pron +) zvratné zájmeno}
\item[{rel}] {vztažné zájmeno}
\item[{sg}] {jednotné číslo, singulár}
\item[{slang.}] {slang}
\item[{sport.}] {sport}
\item[{stav.}] {stavebnictví, architektura}
\item[{subs}] {podstatné jméno (bez rodu)}
\item[{sup}] {3. stupeň stupňování, superlativ}
\item[{supin}] {supinum}
\item[{škol.}] {školství}
\item[{techn.}] {technika, mechanika}
\item[{v}] {sloveso}
\item[{voj.}] {vojenství}
\item[{zast.}] {zastarale}
\item[{zkr}] {zkratka}
\item[{zool.}] {zoologie}

\item[{;}] {středník uprostřed definice odděluje věty; středník na konci definice odděluje významy}

\item[{(...)}] {v kulatých závorkách jsou uvedeny doplňkové informace}
\item[{[...]}] {v hranatých závorkách je uveden fonetický zápis výslovnosti}\footnote{Vysvětlení fonetických symbolů je uvedeno v kapitole Fonetika \ref{sec:phon_phonems} na straně \pageref{sec:phon_phonems}}
\item[{|}] {naznačuje místo, kde se deklinační koncovka připojuje ke slovnímu základu}
\item[{·}] {vedlejší dělení složeného slova}
\item[{··}] {hlavní dělení složeného slova}
\item[{+}] {plus}
\item[{/}] {lomítko označuje více možností}
\item[{\dicsymSee}] {šipka odkazuje na heslo ve slovníku}
\item[{\dicsymCompare}] {šipka odkazuje na heslo ve slovníku ke srovnání}
\item[{\dicsymExampleIS}] {uvozuje islandský příklad}
\item[{\dicsymExampleCS}] {uvozuje český překlad příkladu}
\item[{\dicsymIdiom}] {uvozuje slovní spojení}
\item[{\dicsymFrequent}] {heslo patří mezi frekventované výrazy}

\item[{\dicsymPhoto}] {fotografie, ilustrace}
\item[{\dicsymProverb}] {přísloví, rčení}

\end{description}
\end{multicols}

\chapter{Skrá yfir skammstafanir og tákn}
\begin{multicols}{2}
\begin{description}[leftmargin=!,labelwidth=1.5cm]
%[\breaklabel\setleftmargin{40pt}\setlabelstyle{\bfseries}
%\addtolength{\itemsep}{-0.5\baselineskip}]
\item[{abb}] {stýfður}
\item[{acc}] {þolfall}
\item[{adj}] {lýsingarorð}
\item[{adv}] {atviksorð}
\item[{akt}] {germynd}
\item[{anat.}] {líffærafræði}
\item[{angl.}] {enska}
\item[{astro.}] {geimvísindi, stjörnufræði}
\item[{ap.}] {e.þ.h.}
\item[{básn.}] {ljóðrænn}
\item[{biol.}] {líffræði}
\item[{bot.}] {grasafræði}

\item[{comp}] {miðstig}
\item[{con}] {viðtengingarháttur}
\item[{conj}] {samtenging}
\item[{dat}] {þágufall}
\item[{def}] {með greini}
\item[{dem}] {ábendingarfornafn}
\item[{dět.}] {barnamál}

\item[{e-að}] {(\textit{eitthvað})}
\item[{e-ð}] {(\textit{eitthvað})}
\item[{e-ir}] {(\textit{einhverjir})}
\item[{e-ja}] {(\textit{einhverja})}
\item[{e-jum}] {(\textit{einhverjum})}
\item[{e-m}] {(\textit{einhverjum})}
\item[{e-n}] {(\textit{einhvern})}
\item[{e-r}] { (\textit{einhver})}
\item[{e-rra}] {(\textit{einhverra})}
\item[{e-rs}] {(\textit{einhvers})}
\item[{e-s}] {(\textit{einhvers})}
\item[{e-u}] { (\textit{einhverju})}

\item[{ekon.}] {hagfræði, viðskipti}
\item[{elek.}] {rafmagn}
\item[{f}] {kvenkyn}
\item[{filos.}] {heimspeki, rökfræði}
\item[{form.}] {formlegt mál}
\item[{fyz.}] {eðlisfræði}
\item[{gen}] {eignarfall}
\item[{geog.}] {landafræði}
\item[{geol.}] {jarðfræði}
\item[{han.}] {niðrandi}
\item[{hist.}] {sagnfræði}
\item[{hovor.}] {talmál}
\item[{hrub.}] {gróft}
\item[{hud.}] {tónlist}
\item[{imper}] {boðháttur}
\item[{impers}] {ópersónuleg sögn}
\item[{ind}] {framsöguháttur}
\item[{indecl}] {óbeygjanlegur}
\item[{indef}] {óákveðinn, (pron +) óákveðið fornafn}
\item[{inf}] {nafnháttur}
\item[{int}] {spurnigarfornafn}
\item[{inter}] {upphrópun}
\item[{chem.}] {efnafræði}
\item[{jaz.}] {málfræði, málvísindi}
\item[{kulin.}] {eldhús}
\item[{l.}] {latína}
\item[{let.}] {flug, flugmál}
\item[{lit.}] {skáldskapur}
\item[{m}] {nafnorð, karlkyn}
\item[{mat.}] {stærðfræði}
\item[{med. }] {læknisfræði}
\item[{meteo.}] {veðurfræði}
\item[{myt.}] {goðafræði}
\item[{n}] {nafnorð, hvorugkyn}
\item[{náb.}] {trúfræði}
\item[{nám.}] {sjómennska}
\item[{nom}] {nefnifall}
\item[{num}] {töluorð}
\item[{ord}] {raðtala}
\item[{p}] {persóna}
\item[{part}] {ögn}
\item[{pers}] {(v +) persóna, (pron +) persónufornafn}
\item[{pf}] {þátíð}
\item[{pl}] {fleirtala}
\item[{poč.}] {tölvufræði, tölvur}
\item[{pol.}] {stjórnmál, stjórnmálafræði}
\item[{poss}] {eignarfornafn}
\item[{pos}] {frumstig}
\item[{pov.}] {þjóðtrú}
\item[{pp}] {lýsingarháttur þátíðar}
\item[{praes}] {nútíð}
\item[{práv.}] {lögfræði, dómsmál}
\item[{predp}] {forskeyti}
\item[{prep}] {forsetning}
\item[{presp}] {lýsingarháttur nútíðar}
\item[{pron}] {fornafn}
\item[{prop}] {sérnafn}
\item[{přen.}] {afleidd merking}
\item[{přís.}] {málsháttur}
\item[{psych.}] {sálfræði}
\item[{refl}] {(v +) miðmynd, (pron +) afturbeygt fornafn}
\item[{rel}] {tilvísunarfornafn}
\item[{sg}] {eintala}
\item[{slang.}] {slangur}
\item[{sport.}] {íþróttir}
\item[{stav.}] {húsasmíði, byggingarlist}
\item[{subs}] {nafnorð (án kyns)}
\item[{sup}] {efstastig}
\item[{supin}] {sagnbót}
\item[{škol.}] {skólamál}
\item[{techn.}] {tækni, verkfræði}
\item[{v}] {sagnorð}
\item[{voj.}] {her}
\item[{zast.}] {úrelt mál}
\item[{zkr}] {skammstöfun}
\item[{zool.}] {dýrafræði}

\item[{;}] {semíkomma í miðri skýringu aðgreinir setningu; semíkomma í enda skýringar aðgreinir mismunandi merkingarsvið}

\item[{(...)}] {innan sviga eru viðbótarupplýsingar}
\item[{[...]}] {innan hornklofa er hljóðritaður framburður}
\item[{|}] {merkir stað, þar sem beygingarending bætist við stofn}
\item[{·}] {aukasamskeyti samsetts orðs}
\item[{··}] {aðalsamskeyti samsetts orðs}
\item[{+}] {plús}
\item[{/}] {skástrik táknar fleiri möguleika}
\item[{\dicsymSee}] {ör vísar til uppflettiorðs í orðabókinni}
\item[{\dicsymCompare}] {ör vísar til uppflettiorðs í bókinni til samanburðar}
\item[{\dicsymExampleIS}] {kemur á undan dæmi á íslensku}
\item[{\dicsymExampleCS}] {kemur á undan tékkneskri þýðingu á dæminu}
\item[{\dicsymIdiom}] {orð með háa tíðni}

\item[{\dicsymPhoto}] {ljósmynd}
\item[{\dicsymProverb}] {málsháttur}

\end{description}
\end{multicols}

% ------------------------------------------------- Headword part and alphabet

\chapter{Heslová část}
{\renewcommand{\arraystretch}{2}%
%\topskip0pt
\vstretch{1}
{\setlength{\tabcolsep}{1.5em}
\begin{table}[htpb]
\centering

\begin{tabular}{ccccccc} \hline
A a &  & G g &  & O o &  & V v \\
Á á  &  & H h &  & Ó ó &  & W w \\
B b &  & I i &  & P p &  & X x \\
C c &  & Í í &  & Q q &  & Y y \\
D d &  & J j &  & R r &  & Ý ý \\
Ð ð &  & K k &  & S s &  & Z z \\
E e &  & L l &  & T t &  & Þ þ \\
É é  &  & M m &  & U u &  & Æ æ \\
F f &  & N n &  & Ú ú &  & Ö ö \\
\hline
\end{tabular}

%\caption{\textbf{Islandská abeceda}}
\label{my-islandska-abeceda}
\end{table}
}
\vstretch{1}
{\renewcommand{\arraystretch}{1}%

\clearpage

\fi

% =========================================================== THE DICTIONARY =

\ifinputletters

\cleardoublepage

\dictionarygeometry
\pagestyle{myheadings}

\newcommand*{\alphabetF}{%
  a,aa,b,c,d,e,ee,f,g,h,i,ii,j,k,l,m,n,o,oo,p,r,s1,s2,t,u,uu,v,w,x,y,yy,th,ae,oe}
\ifPDFtest
\dicLetter{o}{letter18}
\dicEntry[o] \dicTerm{o} \dicsymFrequent\  \dicIPA{{\textopeno}{\textlengthmark}} \dicPos{inter} \textbf{1.} \dicIndirectTranslationCS{citoslovce vyjadřující malý zájem (o~otázku ap.)}  \textbf{2.} \dicIndirectTranslationCS{citoslovce vyjadřující nadšení, obdiv}  \textbf{3.} \dicIndirectTranslationCS{málo významné slovo používané k~zdůraznění před jinými citoslovci} \dicExampleIS{o~jæja} \dicExampleCS{jojo}
\dicEntry[o.fl.] \dicTerm{o.fl.} \dicPos{zkr} \dicPhraseIS{og fleiri, og fleira} \dicFlx{adv} \dicDirectTranslationCS{a~podobně}
\dicEntry[o.m.fl.] \dicTerm{o.m.fl.} \dicPos{zkr} \dicPhraseIS{og margt fleira} \dicFlx{adv} \dicDirectTranslationCS{a~mnohem více}
\dicEntry[o.s.frv.] \dicTerm{o.s.frv.} \dicPos{zkr} \dicPhraseIS{og svo framvegis} \dicFlx{adv} \dicDirectTranslationCS{a~tak dále}
\dicEntry[o.þ.h.] \dicTerm{o.þ.h.} \dicPos{zkr} \dicPhraseIS{og þess háttar} \dicFlx{adv} \dicDirectTranslationCS{a~podobně}
\dicEntry[o.þ.l.] \dicTerm{o.þ.l.} \dicPos{zkr} \dicPhraseIS{og því um líkt} \dicFlx{adv} \dicDirectTranslationCS{a~podobně}
\dicEntry[obbi] \dicTerm{obb|i} \dicIPA{{\textopeno}{\textsubring{b}}{\textlengthmark}{\textsci}} \dicPos{m}[1] \dicFlx{(‑a)}[3] \dicDirectTranslationCS{větší část};  \dicPhraseIS{obbinn af e‑u} \dicDirectTranslationCS{větší část (čeho)}
\dicEntry[obláta] \dicTerm{oblát|a} \dicIPA{{\textopeno}{\textsubring{b}}{l}{au}{\textsubring{d}}{a}} \dicPos{f}[1] \dicFlx{(‑u, ‑ur)}[19] \dicFieldCat{náb.} \dicDirectTranslationCS{hostie}
\dicEntry[oddaaðstaða] \dicTerm{odda··að·|staða} \dicIPA{{\textopeno}{\textsubring{d}}{\textlengthmark}{a}{a}{ð}{s}{\textsubring{d}}{a}{ð}{a}} \dicPos{f}[1] \dicFlx{(‑stöðu)}[2] \dicDirectTranslationCS{klíčové postavení}
\dicEntry[oddatala] \dicTerm{odda··|tala} \dicIPA{{\textopeno}{\textsubring{d}}{\textlengthmark}{a}{t\smash{\textsuperscript{h}}}{a}{l}{a}} \dicPos{f}[1] \dicFlx{(‑tölu, ‑tölur)}[14] \dicFieldCat{mat.} \dicDirectTranslationCS{liché číslo} \dicAntonym*{jöfn tala}
\dicEntry[oddbogi] \dicTerm{odd··bog|i} \dicIPA{{\textopeno}{\textsubring{d}}{\textsubring{b}}{\textopeno i}{j}{\textsci}} \dicPos{m}[1] \dicFlx{(‑a, ‑ar)}[1] \dicDirectTranslationCS{lomený oblouk}
\dicEntry[oddi] \dicTerm{odd|i} \dicIPA{{\textopeno}{\textsubring{d}}{\textlengthmark}{\textsci}} \dicPos{m}[1] \dicFlx{(‑a, ‑ar)}[1] \textbf{1.} \dicSynonym{tangi} \dicDirectTranslationCS{výběžek, ostroh (země ap.)} \dicExampleIS{Húsið stendur úti á oddanum.} \dicExampleCS{Dům stojí na výběžku.}  \textbf{2.} \dicSynonym{oddur} \dicDirectTranslationCS{hrot, špice};  \dicPhraseIS{það skerst í odda með e‑jum} \dicFlx{refl} \dicLangCat{přen.} \dicDirectTranslationCS{(kdo) je ve sporu}
\dicEntry[oddmjór] \dicTerm{odd··mjór} \dicIPA{{\textopeno}{\textsubring{d}}{m}{j}{ou}{\textsubring{r}}} \dicPos{adj}[4]\dicFlx{}[-5] \dicDirectTranslationCS{špičatý, zašpičatělý}
\dicEntry[oddur] \dicTerm{odd|ur} \dicIPA{{\textopeno}{\textsubring{d}}{\textlengthmark}{\textscy}{\textsubring{r}}} \dicPos{m}[6] \dicFlx{(‑s, ‑ar)}[4] \dicDirectTranslationCS{hrot, špička, osten};  \dicPhraseIS{brjóta odd af oflæti sínu} \dicLangCat{přen.} \dicDirectTranslationCS{pokořit se, ustoupit};  \dicPhraseIS{leika á als oddi} \dicLangCat{přen.} \dicDirectTranslationCS{být radostí bez sebe};  \dicPhraseIS{með oddi og egg(ju)} \dicFlx{adv} \dicLangCat{přen.} \dicDirectTranslationCS{ze všech sil};  \dicPhraseIS{setja e‑ð á oddinn} \dicDirectTranslationCS{dávat (čemu) prioritu}
\dicEntry[oddviti] \dicTerm{odd··vit|i} \dicIPA{{\textopeno}{\textsubring{d}}{v}{\textsci}{\textsubring{d}}{\textsci}} \dicPos{m}[1] \dicFlx{(‑a, ‑ar)}[1] \dicSynonym{formaður} \dicDirectTranslationCS{přednosta, přednostka, předák, předačka}
\dicEntry[of] \dicTerm{of\smash{\textsuperscript{1}}} \dicsymFrequent\  \dicIPA{{\textopeno}{\textlengthmark}{f}} \dicPos{adv} \dicDirectTranslationCS{příliš, nadmíru, moc} \dicExampleIS{of seint} \dicExampleCS{příliš pozdě}
\dicEntry[of] \dicTerm{of-\smash{\textsuperscript{2}}} \dicIPA{{\textopeno}{\textlengthmark}{f}} \dicPos{predp} \dicDirectTranslationCS{pře-, nad-} \dicExampleIS{ofneysla} \dicExampleCS{nadměrná spotřeba}
\dicEntry[ofan] \dicTerm{ofan} \dicsymFrequent\  \dicIPA{{\textopeno}{\textlengthmark}{v}{a}{\textsubring{n}}} \dicPos{prep\,/\addthin adv} \dicFlx{gen\,/\addthin acc} \textbf{1.} \dicFlx{prep acc} \dicDirectTranslationCS{(ze) shora, dolů z} \dicIndirectTranslationCS{(o~pohybu dolů)} \dicExampleIS{fara ofan stigann} \dicExampleCS{jít (dolů) ze schodů}  \textbf{2.} \dicFlx{prep gen} \dicDirectTranslationCS{nad} \dicExampleIS{ofan götunnar} \dicExampleCS{nad ulicí} \dicAntonym{neðan}  \textbf{3.} \dicFlx{adv} \dicDirectTranslationCS{dolů, shora, seshora};  \dicIdiom{að}[ofan]{ \dicPhraseIS{að ofan}} \dicFlx{adv} {\textbf{a.}} \dicDirectTranslationCS{seshora, shora, svrchu};  {\textbf{b.}} \dicDirectTranslationCS{nahoře, na horní straně};  {\textbf{c.}} \dicDirectTranslationCS{výše (v~textu ap.), nahoře};  \dicIdiom{ofan}[að]{ \dicPhraseIS{ofan að e‑u}} \dicFlx{prep} \dicDirectTranslationCS{dolů k~(čemu)}; { \dicPhraseIS{ofan að}} \dicFlx{adv} \dicDirectTranslationCS{shora, seshora};  \dicIdiom{ofan}[af]{ \dicPhraseIS{ofan af e‑u}} \dicFlx{prep} \dicDirectTranslationCS{shora z~(čeho)};  \dicIdiom{ofan}[á]{ \dicPhraseIS{ofan á e‑u}} \dicFlx{prep} \dicDirectTranslationCS{navrchu (čeho), nahoře na (čem)}; { \dicPhraseIS{ofan á e‑ð}} \dicFlx{prep} \dicDirectTranslationCS{navrch (čeho), nahoru na (co)};  \dicIdiom{ofan}[eftir]{ \dicPhraseIS{ofan eftir e‑u}} \dicFlx{prep} \dicDirectTranslationCS{dolů po (čem)}; { \dicPhraseIS{ofan eftir}} \dicFlx{adv} \dicDirectTranslationCS{dolů} \dicIndirectTranslationCS{(o~místě, které leží níže než výchozí bod)};  \dicIdiom{ofan}[frá]{ \dicPhraseIS{ofan frá e‑u}} \dicFlx{prep} \dicDirectTranslationCS{seshora z~(čeho)}; { \dicPhraseIS{ofan frá}} \dicFlx{adv} \dicDirectTranslationCS{shora, seshora};  \dicIdiom{ofan}[í]{ \dicPhraseIS{ofan í e‑ð, oní e‑ð}} \dicFlx{prep} \dicDirectTranslationCS{dolů do (čeho)}; { \dicPhraseIS{ofan í e‑u, oní e‑u}} \dicFlx{prep} \dicDirectTranslationCS{dole v~(čem)}; { \dicPhraseIS{ofan í, oní}} \dicFlx{adv} \dicDirectTranslationCS{dovnitř (vany ap.)};  \dicIdiom{ofan}[til]{ \dicPhraseIS{ofan til}} \dicFlx{adv} \dicDirectTranslationCS{v~horní části, nahoře};  \dicIdiom{ofan}[um]{ \dicPhraseIS{ofan um e‑ð}} \dicFlx{prep} \dicDirectTranslationCS{dolů skrz (co), dolů do (čeho)};  \dicIdiom{ofan}[úr]{ \dicPhraseIS{ofan úr e‑u}} \dicFlx{prep} \dicDirectTranslationCS{(dolů) z~(čeho)};  \dicIdiom{ofan}[við]{ \dicPhraseIS{ofan við e‑ð}} \dicFlx{prep} \dicDirectTranslationCS{nad (čím)} \dicIndirectTranslationCS{(místně)}
\dicEntry[ofanálag] \dicTerm{ofan··á·lag} \dicIPA{{\textopeno}{\textlengthmark}{v}{a}{n}{au}{l}{a}{x}} \dicPos{n}[2] \dicFlx{(‑s)}[2] \dicDirectTranslationCS{nádavek, přídavek};  \dicPhraseIS{í ofanálag} \dicFlx{adv} \dicDirectTranslationCS{nádavkem, navíc}
\dicEntry[ofangreindur] \dicTerm{ofan··greindur} \dicIPA{{\textopeno}{\textlengthmark}{v}{a}{\ng}{\r{g}}{r}{ei}{n}{\textsubring{d}}{\textscy}{\textsubring{r}}} \dicPos{adj}[2]\dicFlx{}[-17] \dicDirectTranslationCS{výše uvedený\,/\addthin zmíněný}
\dicEntry[ofanígjöf] \dicTerm{ofan·í··|gjöf} \dicIPA{{\textopeno}{\textlengthmark}{v}{a}{n}{i}{\r{\textObardotlessj}}{\oe}{f}} \dicPos{f}[7] \dicFlx{(‑gjafar)}[19] \textbf{1.} \dicSynonym{ávítur} \dicDirectTranslationCS{(po)kárání, důtka}  \textbf{2.} \dicSynonym{uppbót} \dicDirectTranslationCS{kompenzace, náhrada}
\dicEntry[ofanverður] \dicTerm{ofan··verður} \dicIPA{{\textopeno}{\textlengthmark}{v}{a}{n}{v}{\textepsilon}{r}{ð}{\textscy}{\textsubring{r}}} \dicPos{adj}[2]\dicFlx{}[-4] \textbf{1.} \dicSynonym*{sem heyrir til efra} \dicDirectTranslationCS{horní, vrchní}  \textbf{2.} \dicSynonym*{sem heyrir til síðara hluta tímabils} \dicDirectTranslationCS{pozdní} \dicExampleIS{ofanverður vetur} \dicExampleCS{pozdní zima} \dicAntonym{öndverður};  \dicPhraseIS{ofanverð nítjánda öld} \dicDirectTranslationCS{druhá polovina 19. století};  \dicPhraseIS{að ofanverðu} \dicFlx{adv} \dicDirectTranslationCS{nahoře, seshora}
\dicEntry[ofar] \dicTerm{ofar} \dicsymFrequent\  \dicIPA{{\textopeno}{\textlengthmark}{v}{a}{\textsubring{r}}} \dicPos{adv} \dicFlx{comp (pos uppi, sup efst)} \textbf{1.} \dicSynonym{hærra} \dicDirectTranslationCS{výš(e)} \dicExampleIS{ofar í fjallinu} \dicExampleCS{výše v~horách}  \textbf{2.} \dicSynonym*{fyrir ofan} \dicDirectTranslationCS{nad}
\dicEntry[ofaukinn] \dicTerm{of··aukinn} \dicIPA{{\textopeno}{\textlengthmark}{v}{\oe i}{\r{\textObardotlessj}}{\textsci}{\textsubring{n}}} \dicPos{adj}[6]\dicFlx{}[-6] \dicDirectTranslationCS{přehnaný, (jsoucí) navíc};  \dicPhraseIS{e‑u er ofaukið} \dicFlx{impers} \dicDirectTranslationCS{(co) je přebytečné, (co) je navíc}
\dicEntry[ofbauð] \dicTerm{of··bauð} \dicIPA{{\textopeno}{v}{\textsubring{b}}{\oe i}{\texttheta}} \dicPos{v} \dicFlx{ind pf sg 1 pers} \dicLink{ofbjóða}
\dicEntry[ofbeldi] \dicTerm{of··beldi} \dicIPA{{\textopeno}{v}{\textsubring{b}}{\textepsilon}{l}{\textsubring{d}}{\textsci}} \dicPos{n}[2] \dicFlx{(‑s)}[20] \dicSynonym{ofríki} \dicDirectTranslationCS{násilí, teror, agrese}
\dicEntry[ofbeldishneigður] \dicTerm{of·beldis··hneigður} \dicIPA{{\textopeno}{v}{\textsubring{b}}{\textepsilon}{l}{\textsubring{d}}{\textsci}{s}{\textsubring{n}}{ei}{\textbabygamma}{ð}{\textscy}{\textsubring{r}}} \dicPos{adj}[2]\dicFlx{}[-1] \dicDirectTranslationCS{násilnický}
\dicEntry[ofbeldismaður] \dicTerm{of·beldis··|maður} \dicIPA{{\textopeno}{v}{\textsubring{b}}{\textepsilon}{l}{\textsubring{d}}{\textsci}{s}{m}{a}{ð}{\textscy}{\textsubring{r}}} \dicPos{m}[13] \dicFlx{(‑manns, ‑menn)}[2] \dicDirectTranslationCS{násilník, násilnice, agresor(ka)}
\dicEntry[ofbeldisseggur] \dicTerm{of·beldis··segg|ur} \dicIPA{{\textopeno}{v}{\textsubring{b}}{\textepsilon}{l}{\textsubring{d}}{\textsci}{s}{\textepsilon}{\r{g}}{\textscy}{\textsubring{r}}} \dicPos{m}[9] \dicFlx{(‑s, ‑ir)}[15] \dicDirectTranslationCS{násilník, násilnice, agresor(ka)}
\dicEntry[ofbeldisverk] \dicTerm{of·beldis··verk} \dicIPA{{\textopeno}{v}{\textsubring{b}}{\textepsilon}{l}{\textsubring{d}}{\textsci}{s}{v}{\textepsilon}{\textsubring{r}}{\r{g}}} \dicPos{n}[2] \dicFlx{(‑s, ‑)}[5] \dicDirectTranslationCS{násilný čin, násilí}
\dicEntry[ofbirta] \dicTerm{of··birt|a} \dicIPA{{\textopeno}{v}{\textsubring{b}}{\textsci}{\textsubring{r}}{\textsubring{d}}{a}} \dicPos{f}[1] \dicFlx{(‑u)}[5] \dicDirectTranslationCS{oslnění};  \dicPhraseIS{fá ofbirtu í augun} \dicDirectTranslationCS{být oslněn\,/\addthin oslepen}
\dicEntry[ofbjóða] \dicTerm{of··|bjóða} \dicIPA{{\textopeno}{v}{\textsubring{b}}{j}{ou}{ð}{a}} \dicPos{v}[6] \dicFlx{(‑býð, ‑bauð, ‑buðum, ‑byði, ‑boðið)}[98] \dicFlx{dat} \textbf{1.} \dicPhraseIS{e‑m ofbýður e‑að} \dicFlx{impers} \dicDirectTranslationCS{(koho co) šokuje, (koho co) děsí}  \textbf{2.} \dicPhraseIS{ofbjóða e‑m} \dicDirectTranslationCS{přepínat (koho), přetěžovat (koho)} \dicExampleIS{ofbjóða sjálfum sér} \dicExampleCS{přetěžovat se}
\dicEntry[ofboð] \dicTerm{of··boð} \dicIPA{{\textopeno}{v}{\textsubring{b}}{\textopeno}{\texttheta}} \dicPos{n}[2] \dicFlx{(‑s)}[2] \dicSynonym{angist} \dicDirectTranslationCS{šok, úděs, hrůza};  \dicPhraseIS{í dauðans ofboði} \dicFlx{adv} \dicDirectTranslationCS{v~panice, s~hrůzou v~očích}
\dicEntry[ofboðið] \dicTerm{of··boðið} \dicIPA{{\textopeno}{v}{\textsubring{b}}{\textopeno}{ð}{\textsci}{\texttheta}} \dicPos{v} \dicFlx{supin} \dicLink{ofbjóða}
\dicEntry[ofboðslegur] \dicTerm{of·boðs··legur} \dicIPA{{\textopeno}{v}{\textsubring{b}}{\textopeno}{ð}{s}{l}{\textepsilon}{\textbabygamma}{\textscy}{\textsubring{r}}} \dicPos{adj}[1]\dicFlx{}[-8] \dicSynonym{skelfilegur} \dicDirectTranslationCS{příšerný, děsivý} \dicExampleIS{ofboðsleg læti} \dicExampleCS{příšerný kravál}
\dicEntry[ofborga] \dicTerm{of··borg|a} \dicIPA{{\textopeno}{v}{\textsubring{b}}{\textopeno}{r}{\r{g}}{a}} \dicPos{v}[1] \dicFlx{(‑aði)}[1] \dicFlx{acc} \dicDirectTranslationCS{přeplatit, přeplácet}
\dicEntry[ofbuðum] \dicTerm{of··buðum} \dicIPA{{\textopeno}{v}{\textsubring{b}}{\textscy}{ð}{\textscy}{\textsubring{m}}} \dicPos{v} \dicFlx{ind pf pl 1 pers} \dicLink{ofbjóða}
\dicEntry[ofbyði] \dicTerm{of··byði} \dicIPA{{\textopeno}{v}{\textsubring{b}}{\textsci}{ð}{\textsci}} \dicPos{v} \dicFlx{con pf sg 1 pers} \dicLink{ofbjóða}
\dicEntry[ofbýð] \dicTerm{of··býð} \dicIPA{{\textopeno}{v}{\textsubring{b}}{i}{\texttheta}} \dicPos{v} \dicFlx{ind praes sg 1 pers} \dicLink{ofbjóða}
\dicEntry[ofdirfska] \dicTerm{of··dirfsk|a} \dicIPA{{\textopeno}{v}{\textsubring{d}}{\textsci}{\textsubring{r}}{s}{\r{g}}{a}} \dicPos{f}[1] \dicFlx{(‑u)}[5] \dicDirectTranslationCS{přílišná troufalost}
\dicEntry[ofdrykkja] \dicTerm{of··drykkj|a} \dicIPA{{\textopeno}{v}{\textsubring{d}}{r}{\textsci}{h}{\r{\textObardotlessj}}{a}} \dicPos{f}[1] \dicFlx{(‑u)}[5] \dicDirectTranslationCS{nadměrné pití (alkoholu)}
\dicEntry[ofdrykkjumaður] \dicTerm{of·drykkju··|maður} \dicIPA{{\textopeno}{v}{\textsubring{d}}{r}{\textsci}{h}{\r{\textObardotlessj}}{\textscy}{m}{a}{ð}{\textscy}{\textsubring{r}}} \dicPos{m}[13] \dicFlx{(‑manns, ‑menn)}[2] \dicSynonym{bytta} \dicDirectTranslationCS{opilec, opilkyně, alkoholik, alkoholička}
\dicEntry[offita] \dicTerm{of··fit|a} \dicIPA{{\textopeno}{f}{\textlengthmark}{\textsci}{\textsubring{d}}{a}} \dicPos{f}[1] \dicFlx{(‑u)}[5] \dicFieldCat{med.} \dicDirectTranslationCS{otylost, obezita}
\dicEntry[offjárfesting] \dicTerm{of··fjár·fest·ing} \dicIPA{{\textopeno}{f}{\textlengthmark}{j}{au}{\textsubring{r}}{f}{\textepsilon}{s}{\textsubring{d}}{i}{\ng}{\r{g}}} \dicPos{f}[4] \dicFlx{(‑ar, ‑ar)}[5] \dicFieldCat{ekon.} \dicDirectTranslationCS{nadinvestice}
\dicEntry[offjölgun] \dicTerm{of··fjölg|un} \dicIPA{{\textopeno}{f}{\textlengthmark}{j}{\oe}{l}{\r{g}}{\textscy}{\textsubring{n}}} \dicPos{f}[7] \dicFlx{(‑unar)}[12] \dicDirectTranslationCS{přelidnění, přemnožení}
\dicEntry[offramleiðsla] \dicTerm{of··fram·leiðsl|a} \dicIPA{{\textopeno}{f}{\textlengthmark}{r}{a}{m}{l}{ei}{ð}{s}{\textsubring{d}}{l}{a}} \dicPos{f}[1] \dicFlx{(‑u)}[5] \dicDirectTranslationCS{nadvýroba, nadprodukce}
\dicEntry[ofhitna] \dicTerm{of··hitn|a} \dicIPA{{\textopeno}{\textlengthmark}{f}{h}{\textsci}{h}{\textsubring{d}}{n}{a}} \dicPos{v}[1] \dicFlx{(‑aði)}[72] \dicDirectTranslationCS{přehřát se, příliš se zahřát} \dicExampleIS{Vélin hefur ofhitnað.} \dicExampleCS{Motor je přehřátý.}
\dicEntry[ofhitnun] \dicTerm{of··hitn|un} \dicIPA{{\textopeno}{\textlengthmark}{f}{h}{\textsci}{h}{\textsubring{d}}{n}{\textscy}{\textsubring{n}}} \dicPos{f}[7] \dicFlx{(‑unar)}[9] \dicDirectTranslationCS{přehřátí}
\dicEntry[ofhleðsla] \dicTerm{of··hleðsl|a} \dicIPA{{\textopeno}{f}{\textsubring{l}}{\textepsilon}{ð}{s}{\textsubring{d}}{l}{a}} \dicPos{f}[1] \dicFlx{(‑u)}[5] \textbf{1.} \dicDirectTranslationCS{přetížení (vozu ap.)}  \textbf{2.} \dicDirectTranslationCS{přetěžování (žáka ap.)}
\dicEntry[ofið] \dicTerm{ofið} \dicIPA{{\textopeno}{\textlengthmark}{v}{\textsci}{\texttheta}} \dicPos{v} \dicFlx{supin} \dicLink{vefa}
\dicEntry[ofjarl] \dicTerm{of··jarl} \dicIPA{{\textopeno}{v}{j}{a}{r}{\textsubring{d}}{\textsubring{l}}} \dicPos{m}[4] \dicFlx{(‑s, ‑ar)}[4] \dicIndirectTranslationCS{člověk nad (čí) síly};  \dicPhraseIS{hann var mér ofjarl} \dicDirectTranslationCS{byl nad mé síly, byl na mě příliš silný}
\dicEntry[ofkæla] \dicTerm{of··kæl|a} \dicIPA{{\textopeno}{f}{c\smash{\textsuperscript{h}}}{a}{i}{l}{a}} \dicPos{v}[2] \dicFlx{(‑di, ‑t)}[141] \dicDirectTranslationCS{příliš mrazit\,/\addthin chladit (lednička ap.)};  \dicPhraseIS{ofkæla sig} \dicDirectTranslationCS{prochladnout, prostydnout, prostudit se, podchladit se};  \dicIdiom{ofkælast}{ \dicPhraseIS{ofkælast}} \dicFlx{refl} \dicDirectTranslationCS{prochladnout, prostydnout, prostudit se, podchladit se}
\dicEntry[ofkæling] \dicTerm{of··kæl·ing} \dicIPA{{\textopeno}{f}{c\smash{\textsuperscript{h}}}{a}{i}{l}{i}{\ng}{\r{g}}} \dicPos{f}[4] \dicFlx{(‑ar)}[7] \dicDirectTranslationCS{prochladnutí, podchlazení}
\dicEntry[oflátungur] \dicTerm{of··lát·ung|ur} \dicIPA{{\textopeno}{v}{l}{au}{\textsubring{d}}{u}{\ng}{\r{g}}{\textscy}{\textsubring{r}}} \dicPos{m}[6] \dicFlx{(‑s, ‑ar)}[8] \dicDirectTranslationCS{nadutec, ješita}
\dicEntry[oflof] \dicTerm{of··lof} \dicIPA{{\textopeno}{v}{l}{\textopeno}{f}} \dicPos{n}[2] \dicFlx{(‑s)}[2] \dicDirectTranslationCS{přechválení, přecenění};  \dicPhraseIS{bera oflof á e‑n} \dicDirectTranslationCS{přechválit (koho)}
\dicEntry[oflæti] \dicTerm{of··læti} \dicIPA{{\textopeno}{v}{l}{a}{i}{\textsubring{d}}{\textsci}} \dicPos{n}[2] \dicFlx{(‑s)}[20] \textbf{1.} \dicSynonym{dramb} \dicDirectTranslationCS{ješitnost, nadutost}  \textbf{2.} \dicFieldCat{psych.} \dicDirectTranslationCS{mánie}
\dicEntry[ofmetnaður] \dicTerm{of··met·nað|ur} \dicIPA{{\textopeno}{v}{m}{\textepsilon}{h}{\textsubring{d}}{n}{a}{ð}{\textscy}{\textsubring{r}}} \dicPos{m}[10] \dicFlx{(‑ar)}[9] \dicSynonym{hroki} \dicDirectTranslationCS{nadutost, namyšlenost, pýcha}
\dicEntry[ofmetnast] \dicTerm{of··metn|ast} \dicIPA{{\textopeno}{v}{m}{\textepsilon}{h}{\textsubring{d}}{n}{a}{s}{\textsubring{d}}} \dicPos{v}[1] \dicFlx{(‑aðist)}[91] \dicFlx{refl} \dicDirectTranslationCS{zpychnout, stát se namyšleným} \dicExampleIS{ofmetnast af afrekum sínum} \dicExampleCS{zpychnout z~toho, čeho člověk dosáhl}
\dicEntry[ofmæltur] \dicTerm{of··mæltur} \dicIPA{{\textopeno}{v}{m}{a}{i}{\textsubring{l}}{\textsubring{d}}{\textscy}{\textsubring{r}}} \dicPos{adj}[1]\dicFlx{}[-13] \dicSynonym{ýktur} \dicDirectTranslationCS{přehnaný, zveličený}
\dicEntry[ofn] \dicTerm{ofn} \dicIPA{{\textopeno}{\textsubring{b}}{\textsubring{n}}} \dicPos{m}[4] \dicFlx{(‑s, ‑ar)}[3] \textbf{1.} \dicSynonym*{olíuofn} \dicDirectTranslationCS{kamna, sporák}  \textbf{2.} \dicSynonym{bakaraofn} \dicDirectTranslationCS{trouba, pec}  \textbf{3.} \dicSynonym*{miðstöðvarofn} \dicDirectTranslationCS{radiátor, topné těleso}  \textbf{4.} \dicSynonym*{bræðsluofn} \dicDirectTranslationCS{vysoká pec} \dicIndirectTranslationCS{(pec k~výrobě surového železa)}  \textbf{5.} \dicSynonym*{brennsluofn} \dicDirectTranslationCS{keramická pec}
\dicEntry[ofneysla] \dicTerm{of··neysl|a} \dicIPA{{\textopeno}{v}{n}{ei}{s}{\textsubring{d}}{l}{a}} \dicPos{f}[1] \dicFlx{(‑u)}[5] \dicDirectTranslationCS{nadměrná spotřeba, nadužívání}
\dicEntry[ofnæmi] \dicTerm{of··næmi} \dicIPA{{\textopeno}{v}{n}{a}{i}{m}{\textsci}} \dicPos{n}[2] \dicFlx{(‑s)}[20] \dicFieldCat{med.} \dicDirectTranslationCS{alergie}
\dicEntry[ofnæmisviðbrögð] \dicTerm{of·næmis··við·brögð} \dicIPA{{\textopeno}{v}{n}{a}{i}{m}{\textsci}{s}{v}{\textsci}{ð}{\textsubring{b}}{r}{\oe}{\textbabygamma}{\texttheta}} \dicPos{n}[2] \dicFlx{pl}[9] \dicDirectTranslationCS{alergická reakce}
\dicEntry[ofraun] \dicTerm{of··raun} \dicIPA{{\textopeno}{v}{r}{\oe i}{\textsubring{n}}} \dicPos{f}[7] \dicFlx{(‑ar)}[3] \dicPhraseIS{e‑að er e‑m ofraun} \dicDirectTranslationCS{(co) je nad (čí) síly}
\dicEntry[ofreyna] \dicTerm{of··reyn|a} \dicIPA{{\textopeno}{v}{r}{ei}{n}{a}} \dicPos{v}[2] \dicFlx{(‑di, ‑t)}[141] \dicDirectTranslationCS{přepínat, přetěžovat} \dicExampleIS{ofreyna augun} \dicExampleCS{přetěžovat oči};  \dicPhraseIS{ofreyna sig á e‑u} \dicDirectTranslationCS{přepínat se v~(čem)}
\dicEntry[ofríki] \dicTerm{of··ríki} \dicIPA{{\textopeno}{v}{r}{i}{\r{\textObardotlessj}}{\textsci}} \dicPos{n}[2] \dicFlx{(‑s)}[20] \dicSynonym{yfirgangur} \dicDirectTranslationCS{tyranie, krutovláda} \dicExampleIS{ofríki stjórnvalda} \dicExampleCS{tyranie vlády}
\dicEntry[ofsa] \dicTerm{ofsa} \dicIPA{{\textopeno}{f}{s}{a}} \dicPos{adv} \dicLangCat{hovor.} \dicDirectTranslationCS{strašně, hrozně, super, mega, ultra} \dicExampleIS{ofsa skemmtilegur} \dicExampleCS{strašně zábavný}
\dicEntry[ofsafenginn] \dicTerm{ofsa··fenginn} \dicIPA{{\textopeno}{f}{s}{a}{f}{ei}{\textltailn}{\r{\textObardotlessj}}{\textsci}{\textsubring{n}}} \dicPos{adj}[6]\dicFlx{}[-2] \textbf{1.} \dicDirectTranslationCS{zuřivý, zběsilý}  \textbf{2.} \dicSynonym{ofsalegur} \dicDirectTranslationCS{prudký, bouřlivý} \dicExampleIS{ofsafengin viðbrögð} \dicExampleCS{prudké reakce}
\dicEntry[ofsahræðsla] \dicTerm{ofsa··hræðsl|a} \dicIPA{{\textopeno}{f}{s}{a}{\textsubring{r}}{a}{i}{ð}{s}{\textsubring{d}}{l}{a}} \dicPos{f}[1] \dicFlx{(‑u)}[5] \dicDirectTranslationCS{panika}
\dicEntry[ofsakæti] \dicTerm{ofsa··kæt|i} \dicIPA{{\textopeno}{f}{s}{a}{c\smash{\textsuperscript{h}}}{a}{i}{\textsubring{d}}{\textsci}} \dicPos{f}[3] \dicFlx{(‑i)}[3] \dicDirectTranslationCS{euforie}
\dicEntry[ofsalega] \dicTerm{ofsa··lega} \dicIPA{{\textopeno}{f}{s}{a}{l}{\textepsilon}{\textbabygamma}{a}} \dicPos{adv} \dicDirectTranslationCS{úžasně, fantasticky, super} \dicExampleIS{Það var ofsalega gaman.} \dicExampleCS{Byla to super legrace.}
\begin{xtolerant}{}{1pt}
\dicEntry[ofsalegur] \dicTerm{ofsa··legur} \dicIPA{{\textopeno}{f}{s}{a}{l}{\textepsilon}{\textbabygamma}{\textscy}{\textsubring{r}}} \dicPos{adj}[1]\dicFlx{}[-8] \dicSynonym{ofsafenginn} \dicDirectTranslationCS{silný, moc\-ný, prudký} \dicExampleIS{ofsalegur magaverkur} \dicExampleCS{silná bolest břicha}
\end{xtolerant}
\dicEntry[ofsarok] \dicTerm{ofsa··rok} \dicIPA{{\textopeno}{f}{s}{a}{r}{\textopeno}{\r{g}}} \dicPos{n}[2] \dicFlx{(‑s)}[2] \dicSynonym{ofsaveður} \dicDirectTranslationCS{vichřice, prudká bouře}
\dicEntry[ofsaveður] \dicTerm{ofsa··veður} \dicIPA{{\textopeno}{f}{s}{a}{v}{\textepsilon}{ð}{\textscy}{\textsubring{r}}} \dicPos{n}[2] \dicFlx{(‑s)}[28] \dicDirectTranslationCS{mohutná vichřice} \dicIndirectTranslationCS{(Beaufort 11) (28,5--32,6~m/s)}
\dicEntry[ofsi] \dicTerm{ofs|i} \dicIPA{{\textopeno}{f}{s}{\textsci}} \dicPos{m}[1] \dicFlx{(‑a, ‑ar)}[1] \textbf{1.} \dicSynonym{frekja} \dicDirectTranslationCS{bouřlivost, zběsilost} \dicExampleIS{vera fullur ofsa} \dicExampleCS{být naprosto zběsilý}  \textbf{2.} \dicSynonym{hvassviðri} \dicDirectTranslationCS{vichřice, bouře, bouřlivý vítr}
\dicEntry[ofsjónir] \dicTerm{of··sjónir} \dicIPA{{\textopeno}{f}{s}{j}{ou}{n}{\textsci}{\textsubring{r}}} \dicPos{f}[7] \dicFlx{pl}[2] \dicDirectTranslationCS{halucinace, vidiny};  \dicPhraseIS{sjá ofsjónum yfir e‑u} \dicLangCat{přen.} \dicDirectTranslationCS{zezelenat závistí kvůli (čemu)}
\dicEntry[ofskynjun] \dicTerm{of··skynj|un} \dicIPA{{\textopeno}{f}{s}{\r{\textObardotlessj}}{\textsci}{n}{j}{\textscy}{\textsubring{n}}} \dicPos{f}[7] \dicFlx{(‑unar, ‑anir)}[8] \dicDirectTranslationCS{přízrak, přelud, mámení, halucinace}
\dicEntry[ofsókn] \dicTerm{of··sókn} \dicIPA{{\textopeno}{f}{s}{ou}{h}{\r{g}}{\textsubring{n}}} \dicPos{f}[7] \dicFlx{(‑ar, ‑ir)}[1] \dicDirectTranslationCS{pronásledování, perzekuce} \dicExampleIS{stöðugar ofsóknir gegn svertingjum} \dicExampleCS{neustálé perzekuce černochů}
\dicEntry[ofsóknaræði] \dicTerm{of·sóknar··æði} \dicIPA{{\textopeno}{f}{s}{ou}{h}{\r{g}}{n}{a}{r}{a}{i}{ð}{\textsci}} \dicPos{n}[2] \dicFlx{(‑s)}[20] \dicFieldCat{psych.} \dicDirectTranslationCS{paranoia}
\dicEntry[ofsótt] \dicTerm{of··sótt} \dicIPA{{\textopeno}{f}{s}{ou}{h}{\textsubring{d}}} \dicPos{v} \dicFlx{supin} \dicLink{ofsækja}
\dicEntry[ofsótti] \dicTerm{of··sótti} \dicIPA{{\textopeno}{f}{s}{ou}{h}{\textsubring{d}}{\textsci}} \dicPos{v} \dicFlx{ind pf sg 1 pers} \dicLink{ofsækja}
\dicEntry[ofsóttum] \dicTerm{of··sóttum} \dicIPA{{\textopeno}{f}{s}{ou}{h}{\textsubring{d}}{\textscy}{\textsubring{m}}} \dicPos{v} \dicFlx{ind pf pl 1 pers} \dicLink{ofsækja}
\dicEntry[ofstopafullur] \dicTerm{of·stopa··|fullur} \dicIPA{{\textopeno}{f}{s}{\textsubring{d}}{\textopeno}{\textsubring{b}}{a}{f}{\textscy}{\textsubring{d}}{l}{\textscy}{\textsubring{r}}} \dicPos{adj}[10] \dicFlx{(comp ‑fyllri, sup ‑fyll\-st\-ur)}[7] \dicDirectTranslationCS{agresivní, násilnický}
\dicEntry[ofstopamaður] \textls[15]{\dicTerm{of·stopa··|maður} \dicIPA{{\textopeno}{f}{s}{\textsubring{d}}{\textopeno}{\textsubring{b}}{a}{m}{a}{ð}{\textscy}{\textsubring{r}}} \dicPos{m}[13] \dicFlx{(‑manns, ‑menn)}[2]} \dicDirectTranslationCS{agresivní člověk, násilník, násilnice}
\dicEntry[ofstopi] \dicTerm{of··stop|i} \dicIPA{{\textopeno}{f}{s}{\textsubring{d}}{\textopeno}{\textsubring{b}}{\textsci}} \dicPos{m}[1] \dicFlx{(‑a, ‑ar)}[1] \textbf{1.} \dicSynonym{ofsi} \dicDirectTranslationCS{agrese, brutalita}  \textbf{2.} \dicSynonym{ofstopamaður} \dicDirectTranslationCS{agresivní\,/\addthin brutální člověk}
\dicEntry[ofstæki] \dicTerm{of··stæki} \dicIPA{{\textopeno}{f}{s}{\textsubring{d}}{a}{i}{\r{\textObardotlessj}}{\textsci}} \dicPos{n}[2] \dicFlx{(‑s)}[20] \dicDirectTranslationCS{fanatismus} \dicExampleIS{e‑að liggur við ofstæki} \dicExampleCS{(co) hraničí s~fanatismem}
\dicEntry[ofstækisfullur] \dicTerm{of·stækis··|fullur} \dicIPA{{\textopeno}{f}{s}{\textsubring{d}}{a}{i}{\r{\textObardotlessj}}{\textsci}{s}{f}{\textscy}{\textsubring{d}}{l}{\textscy}{\textsubring{r}}} \dicPos{adj}[10] \dicFlx{(comp ‑fyllri, sup ‑fyllstur)}[7] \dicDirectTranslationCS{fanatický}
\dicEntry[ofstækismaður] \dicTerm{of·stækis··|maður} \dicIPA{{\textopeno}{f}{s}{\textsubring{d}}{a}{i}{\r{\textObardotlessj}}{\textsci}{s}{m}{a}{ð}{\textscy}{\textsubring{r}}} \dicPos{m}[13] \dicFlx{(‑manns, ‑menn)}[2] \dicDirectTranslationCS{fanatik, fanatička}
\dicEntry[ofsæki] \dicTerm{of··sæki} \dicIPA{{\textopeno}{f}{s}{a}{i}{\r{\textObardotlessj}}{\textsci}} \dicPos{v} \dicFlx{ind praes sg 1 pers} \dicLink{ofsækja}
\dicEntry[ofsækja] \dicTerm{of··|sækja} \dicIPA{{\textopeno}{f}{s}{a}{i}{\r{\textObardotlessj}}{a}} \dicPos{v}[5] \dicFlx{(‑sæki, ‑sótti, ‑sóttum, ‑sækti, ‑sótt)}[7] \dicFlx{acc} \dicDirectTranslationCS{pronásledovat, perzekvovat}
\dicEntry[ofsækti] \dicTerm{of··sækti} \dicIPA{{\textopeno}{f}{s}{a}{i}{x}{\textsubring{d}}{\textsci}} \dicPos{v} \dicFlx{con pf sg 1 pers} \dicLink{ofsækja}
\dicEntry[ofsæll] \dicTerm{of··sæll} \dicIPA{{\textopeno}{f}{s}{a}{i}{\textsubring{d}}{\textsubring{l}}} \dicPos{adj}[8]\dicFlx{}[-2] \dicDirectTranslationCS{příliš obdařený\,/\addthin nadaný} \dicIndirectTranslationCS{(používané zejména v~záporu)} \dicExampleIS{vera ekki ofsæll af e‑u} \dicExampleCS{nebýt příliš obdařený (čím)}
\dicEntry[oft] \dicTerm{oft} \dicsymFrequent\  \dicIPA{{\textopeno}{f}{\textsubring{d}}} \dicPos{adv} \dicFlx{(comp ‑ar, sup ‑ast)} \dicDirectTranslationCS{často, mnohokrát, nejednou} \dicExampleIS{Við förum oft til útlanda.} \dicExampleCS{Jezdíme často do ciziny.};  \dicPhraseIS{oft og tíðum, oft á tíðum} \dicFlx{adv} \dicDirectTranslationCS{poměrně\,/\addthin docela často}
\dicEntry[oftar] \dicTerm{oft|ar} \dicIPA{{\textopeno}{f}{\textsubring{d}}{a}{\textsubring{r}}} \dicPos{adv} \dicFlx{comp (pos ‑, sup ‑ast)} \dicSynonym{meir} \dicDirectTranslationCS{častěji};  \dicPhraseIS{oftar en ekki} \dicFlx{adv} \dicDirectTranslationCS{většinou, nejčastěji};  \dicPhraseIS{einu sinni en oftar} \dicFlx{adv} \dicDirectTranslationCS{jako tolikrát\,/\addthin obvykle}
\dicEntry[oftast] \dicTerm{oft|ast} \dicIPA{{\textopeno}{f}{\textsubring{d}}{a}{s}{\textsubring{d}}} \dicPos{adv} \dicFlx{sup (pos ‑, comp ‑ar)} \dicSynonym{tíðast} \dicDirectTranslationCS{nejčastěji}
\dicEntry[oftraust] \dicTerm{of··traust} \dicIPA{{\textopeno}{f}{t\smash{\textsuperscript{h}}}{r}{\oe i}{s}{\textsubring{d}}} \dicPos{n}[2] \dicFlx{(‑s)}[2] \dicDirectTranslationCS{přehnaná důvěra};  \dicPhraseIS{bera oftraust til e‑rs} \dicDirectTranslationCS{příliš (komu) důvěřovat}
\dicEntry[ofur] \dicTerm{ofur-} \dicIPA{{\textopeno}{\textlengthmark}{v}{\textscy}{r}} \dicPos{predp} \dicDirectTranslationCS{nad-, supra-}
\dicEntry[ofurást] \dicTerm{ofur··ást} \dicIPA{{\textopeno}{\textlengthmark}{v}{\textscy}{r}{au}{s}{\textsubring{d}}} \dicPos{f}[7] \dicFlx{(‑ar)}[3] \dicDirectTranslationCS{slepá láska}
\dicEntry[ofurefli] \dicTerm{ofur··efli} \dicIPA{{\textopeno}{\textlengthmark}{v}{\textscy}{r}{\textepsilon}{\textsubring{b}}{l}{\textsci}} \dicPos{n}[2] \dicFlx{(‑s)}[20] \dicDirectTranslationCS{přesila, převaha};  \dicPhraseIS{eiga við ofurefli að etja} \dicDirectTranslationCS{muset čelit převaze}
\dicEntry[ofurhugi] \dicTerm{ofur··hug|i} \dicIPA{{\textopeno}{\textlengthmark}{v}{\textscy}{\textsubring{r}}{h}{\textscy}{i}{j}{\textsci}} \dicPos{m}[1] \dicFlx{(‑a, ‑ar)}[1] \dicDirectTranslationCS{odvážlivec, opovážlivec, opovážlivkyně}
\dicEntry[ofurleiðari] \dicTerm{ofur··leið·ar|i} \dicIPA{{\textopeno}{\textlengthmark}{v}{\textscy}{r}{l}{ei}{ð}{a}{r}{\textsci}} \dicPos{m}[1] \dicFlx{(‑a, ‑ar)}[13] \dicFieldCat{fyz.} \dicDirectTranslationCS{supravodič}
\dicEntry[ofurleiðni] \dicTerm{ofur··leiðn|i} \dicIPA{{\textopeno}{\textlengthmark}{v}{\textscy}{r}{l}{ei}{ð}{n}{\textsci}} \dicPos{f}[3] \dicFlx{(‑i)}[3] \dicFieldCat{fyz.} \dicDirectTranslationCS{supravodivost}
\dicEntry[ofurlítill] \dicTerm{ofur··lítill} \dicIPA{{\textopeno}{\textlengthmark}{v}{\textscy}{r}{l}{i}{\textsubring{d}}{\textsci}{\textsubring{d}}{\textsubring{l}}} \dicPos{adj}[11]\dicFlx{}[-11] \dicDirectTranslationCS{malinký, maličký}
\dicEntry[ofurmannlegur] \dicTerm{ofur··mann·legur} \dicIPA{{\textopeno}{\textlengthmark}{v}{\textscy}{r}{m}{a}{n}{l}{\textepsilon}{\textbabygamma}{\textscy}{\textsubring{r}}} \dicPos{adj}[1]\dicFlx{}[-8] \dicDirectTranslationCS{nadlidský}
\dicEntry[ofurmenni] \dicTerm{ofur··menni} \dicIPA{{\textopeno}{\textlengthmark}{v}{\textscy}{r}{m}{\textepsilon}{n}{\textsci}} \dicPos{n}[2] \dicFlx{(‑s, ‑)}[14] \dicDirectTranslationCS{nadčlověk, superman}
\dicEntry[ofurselja] \dicTerm{ofur··sel|ja} \dicIPA{{\textopeno}{\textlengthmark}{v}{\textscy}{\textsubring{r}}{s}{\textepsilon}{l}{j}{a}} \dicPos{v}[4] \dicFlx{(‑di, ‑t)}[23] \dicFlx{acc} \dicSynonym{afhenda} \dicDirectTranslationCS{zaprodat} \dicExampleIS{ofurselja land sitt} \dicExampleCS{zaprodat svoji zemi}
\dicEntry[ofursti] \dicTerm{ofurst|i} \dicIPA{{\textopeno}{\textlengthmark}{v}{\textscy}{\textsubring{r}}{s}{\textsubring{d}}{\textsci}} \dicPos{m}[1] \dicFlx{(‑a, ‑ar)}[1] \dicFieldCat{voj.} \dicDirectTranslationCS{plukovník, plukovnice}
\dicEntry[ofurölvi] \dicTerm{ofur··ölvi} \dicIPA{{\textopeno}{\textlengthmark}{v}{\textscy}{r}{\oe}{l}{v}{\textsci}} \dicPos{adj}[13] \dicFlx{indecl}[1] \dicDirectTranslationCS{(jsoucí) namol, zpitý}
\dicEntry[ofveiði] \dicTerm{of··veið|i} \dicIPA{{\textopeno}{v}{\textlengthmark}{ei}{ð}{\textsci}} \dicPos{f}[2] \dicFlx{(‑i\,/\addthin ‑ar)}[4] \dicDirectTranslationCS{nadměrný rybolov}
\dicEntry[ofviða] \dicTerm{of··viða} \dicIPA{{\textopeno}{v}{\textlengthmark}{\textsci}{ð}{a}} \dicPos{adj}[13] \dicFlx{indecl}[1] \dicPhraseIS{e‑m er e‑að ofviða} \dicFlx{impers} \dicDirectTranslationCS{(co) je nad (čí) síly}
\dicEntry[ofviðri] \dicTerm{of··viðri} \dicIPA{{\textopeno}{v}{\textlengthmark}{\textsci}{ð}{r}{\textsci}} \dicPos{n}[2] \dicFlx{(‑s, ‑)}[14] \dicDirectTranslationCS{silná bouře, vichřice}
\dicEntry[ofvirki] \dicTerm{of··virk|i} \dicIPA{{\textopeno}{v}{\textlengthmark}{\textsci}{\textsubring{r}}{\r{\textObardotlessj}}{\textsci}} \dicPos{m}[1] \dicFlx{(‑ja, ‑jar)}[14] \dicFieldCat{psych.} \dicDirectTranslationCS{hyperaktivní člověk}
\dicEntry[ofvirkni] \dicTerm{of··virkn|i} \dicIPA{{\textopeno}{v}{\textlengthmark}{\textsci}{\textsubring{r}}{\r{g}}{n}{\textsci}} \dicPos{f}[3] \dicFlx{(‑i)}[3] \dicFieldCat{psych.} \dicDirectTranslationCS{hyperaktivita}
\dicEntry[ofvirkur] \dicTerm{of··virkur} \dicIPA{{\textopeno}{v}{\textlengthmark}{\textsci}{\textsubring{r}}{\r{g}}{\textscy}{\textsubring{r}}} \dicPos{adj}[1]\dicFlx{}[-1] \dicFieldCat{psych.} \dicDirectTranslationCS{hyperaktivní} \dicExampleIS{ofvirkur nemandi} \dicExampleCS{hyperaktivní žák}
\dicEntry[ofviti] \dicTerm{of··vit|i} \dicIPA{{\textopeno}{v}{\textlengthmark}{\textsci}{\textsubring{d}}{\textsci}} \dicPos{m}[1] \dicFlx{(‑a, ‑ar)}[1] \dicDirectTranslationCS{génius}
\dicEntry[ofvæni] \dicTerm{of··væni} \dicIPA{{\textopeno}{v}{\textlengthmark}{a}{i}{n}{\textsci}} \dicPos{n}[2] \dicFlx{(‑s)}[20] \dicPhraseIS{bíða í ofvæni eftir e‑u} \dicDirectTranslationCS{napjatě na (co) čekat}
\dicEntry[ofvöxtur] \dicTerm{of··|vöxtur} \dicIPA{{\textopeno}{v}{\textlengthmark}{\oe}{x}{s}{\textsubring{d}}{\textscy}{\textsubring{r}}} \dicPos{m}[11] \dicFlx{(‑vaxtar)}[6] \textbf{1.} \dicDirectTranslationCS{nadměrný růst}  \textbf{2.} \dicFieldCat{med.} \dicDirectTranslationCS{hypertrofie}
\dicEntry[ofþensla] \dicTerm{of··þensl|a} \dicIPA{{\textopeno}{f}{\texttheta}{\textepsilon}{n}{s}{\textsubring{d}}{l}{a}} \dicPos{f}[1] \dicFlx{(‑u)}[5] \dicDirectTranslationCS{nadměrné rozpínání}
\dicEntry[ofþreyta] \dicTerm{of··þreyt|a\smash{\textsuperscript{1}}} \dicIPA{{\textopeno}{f}{\texttheta}{r}{ei}{\textsubring{d}}{a}} \dicPos{f}[1] \dicFlx{(‑u)}[5] \dicDirectTranslationCS{přepínání, přílišná námaha}
\dicEntry[ofþreyta] \dicTerm{of··þreyt|a\smash{\textsuperscript{2}}} \dicIPA{{\textopeno}{f}{\texttheta}{r}{ei}{\textsubring{d}}{a}} \dicPos{v}[2] \dicFlx{(‑ti, ‑t)}[54] \dicFlx{acc} \dicDirectTranslationCS{přepínat, vyčerpávat, vysilovat};  \dicPhraseIS{ofþreyta sig á e‑u} \dicDirectTranslationCS{vysílit se (čím)}
\dicEntry[og] \dicTerm{og\smash{\textsuperscript{1}}} \dicIPA{{\textopeno}{\textlengthmark}{x}} \dicPos{adv} \dicSynonym{einnig} \dicDirectTranslationCS{také, též, rovněž} \dicExampleIS{Hann var þar og.} \dicExampleCS{Taky tam byl.}
\dicEntry[og] \dicTerm{og\smash{\textsuperscript{2}}} \dicsymFrequent\  \dicIPA{{\textopeno}{\textlengthmark}{x}} \dicPos{conj} \dicDirectTranslationCS{a} \dicExampleIS{ég og þú} \dicExampleCS{já a~ty}
\dicEntry[ogonek] \dicTerm{ogonek} \dicIPA{{\textopeno}{\textlengthmark}{\textbabygamma}{\textopeno}{n}{\textepsilon}{\r{g}}} \dicPos{n}[2] \dicFlx{(‑s, ‑)}[5] \dicFieldCat{jaz.} \dicDirectTranslationCS{ocásek, obrácený háček} \dicIndirectTranslationCS{(diakritické znaménko)}
\dicEntry[ok] \dicTerm{ok} \dicIPA{{\textopeno}{\textlengthmark}{\r{g}}} \dicPos{n}[2] \dicFlx{(‑s, ‑)}[5] \textbf{1.} \dicDirectTranslationCS{jho, jařmo (zapřažení)}  \textbf{2.} \dicSynonym{kúgun} \dicDirectTranslationCS{jařmo, otroctví, poroba}
\dicEntry[okkar] \dicTerm{okkar\smash{\textsuperscript{1}}} \dicIPA{{\textopeno}{h}{\r{g}}{a}{\textsubring{r}}} \dicPos{pron}[17] \dicFlx{poss} \dicDirectTranslationCS{náš, naše, naši}
\dicEntry[okkar] \dicTerm{okkar\smash{\textsuperscript{2}}} \dicIPA{{\textopeno}{h}{\r{g}}{a}{\textsubring{r}}} \dicPos{pron}[17] \dicFlx{pers} \dicFlx{gen pl} \dicExampleIS{til okkar} \dicExampleCS{k~nám} \dicLink{við\smash{\textsuperscript{1}}}
\dicEntry[okkitíska] \dicTerm{okkitísk|a} \dicIPA{{\textopeno}{h}{\r{\textObardotlessj}}{\textsci}{\textsubring{d}}{i}{s}{\r{g}}{a}} \dicPos{f}[1] \dicFlx{(‑u)}[5] \dicLink{próvensalska}
\dicEntry[okkur] \dicTerm{okkur\smash{\textsuperscript{1}}} \dicIPA{{\textopeno}{h}{\r{g}}{\textscy}{\textsubring{r}}} \dicPos{n}[2] \dicFlx{(‑s)}[28] \dicDirectTranslationCS{okr}
\dicEntry[okkur] \dicTerm{okkur\smash{\textsuperscript{2}}} \dicIPA{{\textopeno}{h}{\r{g}}{\textscy}{\textsubring{r}}} \dicPos{pron}[17] \dicFlx{pers} \dicFlx{dat acc pl} \dicDirectTranslationCS{nám, nás} \dicExampleIS{frá okkur} \dicExampleCS{od nás} \dicLink{við\smash{\textsuperscript{1}}}
\dicEntry[okra] \dicTerm{okr|a} \dicIPA{{\textopeno}{\textlengthmark}{\r{g}}{r}{a}} \dicPos{v}[1] \dicFlx{(‑aði)}[44] \dicDirectTranslationCS{lichvařit} \dicExampleIS{okra á e‑u} \dicExampleCS{lichvařit s~(čím)}
\dicEntry[okrari] \dicTerm{okr··ar|i} \dicIPA{{\textopeno}{\textlengthmark}{\r{g}}{r}{a}{r}{\textsci}} \dicPos{m}[1] \dicFlx{(‑a, ‑ar)}[13] \dicDirectTranslationCS{lichvář(ka)}
\dicEntry[oksítanska] \dicTerm{oksítanska} \dicIPA{{\textopeno}{\r{g}}{s}{i}{\textsubring{d}}{a}{n}{s}{\r{g}}{a}} \dicPos{f}[1] \dicFlx{(oksítönsku)}[2] \dicLink{próvensalska}
\dicEntry[október] \dicTerm{október} \dicsymFrequent\  \dicIPA{{\textopeno}{x}{\textsubring{d}}{ou}{\textsubring{b}}{\textepsilon}{\textsubring{r}}} \dicPos{m}[15] \dicFlx{indecl}[2] \dicDirectTranslationCS{říjen} \dicExampleIS{einn dag í október} \dicExampleCS{jednoho říjnového dne}
\dicEntry[októbermánuður] \dicTerm{október··mán·|uður} \dicIPA{{\textopeno}{x}{\textsubring{d}}{ou}{\textsubring{b}}{\textepsilon}{r}{m}{au}{n}{\textscy}{ð}{\textscy}{\textsubring{r}}} \dicPos{m}[10] \dicFlx{(‑aðar, ‑uð\-ir)}[42] \dicDirectTranslationCS{(měsíc) říjen}
\dicEntry[okur] \dicTerm{okur} \dicIPA{{\textopeno}{\textlengthmark}{\r{g}}{\textscy}{\textsubring{r}}} \dicPos{n}[2] \dicFlx{(‑s)}[28] \dicDirectTranslationCS{lichva}
\dicEntry[okurleiga] \dicTerm{okur··leig|a} \dicIPA{{\textopeno}{\textlengthmark}{\r{g}}{\textscy}{r}{l}{ei}{\textbabygamma}{a}} \dicPos{f}[1] \dicFlx{(‑u, ‑ur)}[19] \dicDirectTranslationCS{lichvářský nájem}
\dicEntry[okurverð] \dicTerm{okur··verð} \dicIPA{{\textopeno}{\textlengthmark}{\r{g}}{\textscy}{r}{v}{\textepsilon}{r}{\texttheta}} \dicPos{n}[2] \dicFlx{(‑s)}[2] \dicDirectTranslationCS{přemrštěná\,/\addthin lichvářská cena}
\dicEntry[okurvextir] \dicTerm{okur··vextir} \dicIPA{{\textopeno}{\textlengthmark}{\r{g}}{\textscy}{r}{v}{\textepsilon}{x}{s}{\textsubring{d}}{\textsci}{\textsubring{r}}} \dicPos{m}[11] \dicFlx{pl}[1] \dicDirectTranslationCS{lichvářský úrok}
\dicEntry[olía] \dicTerm{olí|a} \dicsymFrequent\  \dicIPA{{\textopeno}{\textlengthmark}{l}{i}{j}{a}} \dicPos{f}[1] \dicFlx{(‑u, ‑ur)}[7] \textbf{1.} \dicDirectTranslationCS{olej (rostlinný, živočišný ap.)} \dicExampleIS{línolía} \dicExampleCS{lněný olej}  \textbf{2.} \dicDirectTranslationCS{benzín, nafta, ropa}  \textbf{3.} \dicSynonym*{olíulitur} \dicDirectTranslationCS{olejová barva};  \dicIdiom{olía}{ \dicPhraseIS{hella olíu í eld}} \dicLangCat{přen.} \dicDirectTranslationCS{přilévat olej do ohně}
\dicEntry[olíubor] \dicTerm{olíu··bor} \dicIPA{{\textopeno}{\textlengthmark}{l}{i}{j}{\textscy}{\textsubring{b}}{\textopeno}{\textsubring{r}}} \dicPos{m}[4] \dicFlx{(‑s, ‑ar)}[14] \dicDirectTranslationCS{ropný vrt}
\dicEntry[olíuborpallur] \dicTerm{olíu·bor··pall|ur} \dicIPA{{\textopeno}{\textlengthmark}{l}{i}{j}{\textscy}{\textsubring{b}}{\textopeno}{\textsubring{r}}{p\smash{\textsuperscript{h}}}{a}{\textsubring{d}}{l}{\textscy}{\textsubring{r}}} \dicPos{m}[6] \dicFlx{(‑s, ‑ar)}[5] \dicDirectTranslationCS{ropná plošina}
\dicEntry[olíuframleiðsla] \dicTerm{olíu··fram·leiðsl|a} \dicIPA{{\textopeno}{\textlengthmark}{l}{i}{j}{\textscy}{f}{r}{a}{m}{l}{ei}{ð}{s}{\textsubring{d}}{l}{a}} \dicPos{f}[1] \dicFlx{(‑u)}[5] \dicDirectTranslationCS{produkce\,/\addthin těžba ropy}
\dicEntry[olíugeymir] \dicTerm{olíu··geym|ir} \dicIPA{{\textopeno}{\textlengthmark}{l}{i}{j}{\textscy}{\r{\textObardotlessj}}{ei}{m}{\textsci}{\textsubring{r}}} \dicPos{m}[7] \dicFlx{(‑is, ‑ar)}[1] \dicDirectTranslationCS{ropná nádrž}
\dicEntry[olíuhreinsunarstöð] \dicTerm{olíu··hreinsunar·stöð} \dicIPA{{\textopeno}\-{\textlengthmark}\-{l}\-{i}\-{j}\-{\textscy}\-{\textsubring{r}}\-{ei}\-{n}\-{s}\-{\textscy}\-{n}\-{a}\-{\textsubring{r}}\-{s}\-{\textsubring{d}}\-{\oe}\-{\texttheta}\-} \dicPos{f}[6] \dicFlx{(‑var, ‑var)}[1] \dicDirectTranslationCS{rafinerie}
\dicEntry[olíukynding] \dicTerm{olíu··kynd·ing} \dicIPA{{\textopeno}{\textlengthmark}{l}{i}{j}{\textscy}{c\smash{\textsuperscript{h}}}{\textsci}{n}{\textsubring{d}}{i}{\ng}{\r{g}}} \dicPos{f}[4] \dicFlx{(‑ar, ‑ar)}[5] \dicDirectTranslationCS{naftové topení}
\dicEntry[olíuleiðsla] \dicTerm{olíu··leiðsl|a} \dicIPA{{\textopeno}{\textlengthmark}{l}{i}{j}{\textscy}{l}{ei}{ð}{s}{\textsubring{d}}{l}{a}} \dicPos{f}[1] \dicFlx{(‑u, ‑ur)}[19] \dicDirectTranslationCS{ropovod}
\dicEntry[olíumálverk] \dicTerm{olíu··mál·verk} \dicIPA{{\textopeno}{\textlengthmark}{l}{i}{j}{\textscy}{m}{au}{l}{v}{\textepsilon}{\textsubring{r}}{\r{g}}} \dicPos{n}[2] \dicFlx{(‑s, ‑)}[5] \dicDirectTranslationCS{olejomalba}
\dicEntry[olíupanna] \dicTerm{olíu··|panna} \dicIPA{{\textopeno}{\textlengthmark}{l}{i}{j}{\textscy}{p\smash{\textsuperscript{h}}}{a}{n}{a}} \dicPos{f}[1] \dicFlx{(‑pönnu, ‑pönnur)}[8] \dicFieldCat{techn.} \dicDirectTranslationCS{jímka}
\dicEntry[olíuskip] \dicTerm{olíu··skip} \dicIPA{{\textopeno}{\textlengthmark}{l}{i}{j}{\textscy}{s}{\r{\textObardotlessj}}{\textsci}{\textsubring{b}}} \dicPos{n}[2] \dicFlx{(‑s, ‑)}[5] \dicFieldCat{nám.} \dicDirectTranslationCS{tanker} \dicIndirectTranslationCS{(loď)}
\dicEntry[olli] \dicTerm{olli} \dicIPA{{\textopeno}{\textsubring{d}}{l}{\textsci}} \dicPos{v} \dicFlx{ind pf sg 1 pers} \textbf{1.} \dicLink{valda\smash{\textsuperscript{2}}}  \textbf{2.} \dicLink{valda\smash{\textsuperscript{3}}}
\dicEntry[ollið] \dicTerm{ollið} \dicIPA{{\textopeno}{\textsubring{d}}{l}{\textsci}{\texttheta}} \dicPos{v} \dicFlx{supin} \dicLink{vella\smash{\textsuperscript{2}}}
\dicEntry[ollum] \dicTerm{ollum} \dicIPA{{\textopeno}{\textsubring{d}}{l}{\textscy}{\textsubring{m}}} \dicPos{v} \dicFlx{ind pf pl 1 pers} \textbf{1.} \dicLink{valda\smash{\textsuperscript{2}}}  \textbf{2.} \dicLink{valda\smash{\textsuperscript{3}}}
\dicEntry[olnboga] \dicTerm{oln··bog|a} \dicIPA{{\textopeno}{l}{\textsubring{b}}{\textopeno}{\textbabygamma}{a}} \dicPos{v}[1] \dicFlx{(‑aði)}[1] \dicFlx{acc} \dicDirectTranslationCS{razit si cestu, drát se (davem ap.)} \dicExampleIS{olnboga sig áfram í mannþrönginni} \dicExampleCS{razit si cestu tlačenicí}
\dicEntry[olnbogabarn] \dicTerm{oln·boga··|barn} \dicIPA{{\textopeno}{l}{\textsubring{b}}{\textopeno}{\textbabygamma}{a}{\textsubring{b}}{a}{r}{\textsubring{d}}{\textsubring{n}}} \dicPos{n}[2] \dicFlx{(‑barns, ‑börn)}[8] \dicDirectTranslationCS{nechtěné dítě}
\dicEntry[olnbogarými] \dicTerm{oln·boga··rými} \dicIPA{{\textopeno}{l}{\textsubring{b}}{\textopeno}{\textbabygamma}{a}{r}{i}{m}{\textsci}} \dicPos{n}[2] \dicFlx{(‑s, ‑)}[14] \dicSynonym{svigrúm} \dicDirectTranslationCS{(časový\,/\addthin manévrovací) prostor, (časová) rezerva}
\dicEntry[olnbogi] \dicTerm{oln··bog|i} \dicIPA{{\textopeno}{l}{\textsubring{b}}{\textopeno i}{j}{\textsci}} \dicPos{m}[1] \dicFlx{(‑a, ‑ar)}[1] \dicFieldCat{anat.} \dicDirectTranslationCS{loket}
\dicEntry[oltið] \dicTerm{oltið} \dicIPA{{\textopeno}{\textsubring{l}}{\textsubring{d}}{\textsci}{\texttheta}} \dicPos{v} \dicFlx{supin} \dicLink{velta\smash{\textsuperscript{3}}}
\dicEntry[oní] \dicTerm{oní} \dicIPA{{\textopeno}{\textlengthmark}{n}{i}} \dicPos{prep\,/\addthin adv} \dicPhraseIS{ofan í} \dicLink{ofan}
\dicEntry[op] \dicTerm{op} \dicIPA{{\textopeno}{\textlengthmark}{\textsubring{b}}} \dicPos{n}[2] \dicFlx{(‑s, ‑)}[5] \dicSynonym{gap} \dicDirectTranslationCS{otvor, díra, štěrbina} \dicExampleIS{op á ílátinu} \dicExampleCS{otvor v~nádobě}
\dicEntry[opinber] \dicTerm{opin··ber} \dicsymFrequent\  \dicIPA{{\textopeno}{\textlengthmark}{\textsubring{b}}{\textsci}{n}{\textsubring{b}}{\textepsilon}{\textsubring{r}}} \dicPos{adj}[5] \dicFlx{(f ‑)}[8] \textbf{1.} \dicDirectTranslationCS{oficiální} \dicExampleIS{opinber heimsókn} \dicExampleCS{oficiální návštěva}  \textbf{2.} \dicDirectTranslationCS{veřejný} \dicExampleIS{opinber umræða} \dicExampleCS{veřejná diskuze}
\dicEntry[opinbera] \dicTerm{opin··ber|a} \dicIPA{{\textopeno}{\textlengthmark}{\textsubring{b}}{\textsci}{n}{\textsubring{b}}{\textepsilon}{r}{a}} \dicPos{v}[1] \dicFlx{(‑aði)}[31] \dicFlx{acc} \dicDirectTranslationCS{zveřejnit, uveřejnit} \dicExampleIS{opinbera trúlofun} \dicExampleCS{zveřejnit zasnoubení}
\dicEntry[opinberlega] \dicTerm{opin·ber··lega} \dicIPA{{\textopeno}{\textlengthmark}{\textsubring{b}}{\textsci}{n}{\textsubring{b}}{\textepsilon}{r}{l}{\textepsilon}{\textbabygamma}{a}} \dicPos{adv} \dicDirectTranslationCS{veřejně}
\dicEntry[opinberun] \dicTerm{opin··ber|un} \dicIPA{{\textopeno}{\textlengthmark}{\textsubring{b}}{\textsci}{n}{\textsubring{b}}{\textepsilon}{r}{\textscy}{\textsubring{n}}} \dicPos{f}[7] \dicFlx{(‑unar, ‑anir)}[8] \textbf{1.} \dicSynonym{birting} \dicDirectTranslationCS{zveřejnění, uveřejnění} \dicExampleIS{opinberun trúlofunar} \dicExampleCS{zveřejnění zasnoubení}  \textbf{2.} \dicFieldCat{náb.} \dicSynonym{vitrun} \dicDirectTranslationCS{zjevení}
\dicEntry[opineygður] \textls[15]{\dicTerm{opin··eygður}\dicTerm{, opineygur} \dicIPA{{\textopeno}\-{\textlengthmark}\-{\textsubring{b}}\-{\textsci}\-{n}\-{ei}\-{\textbabygamma}\-{ð}\-{\textscy}\-{\textsubring{r}}\-} \dicPos{adj}[2]\dicFlx{}[-1] \dicDirectTranslationCS{(jsoucí)  s~očima dokořán, (jsoucí) s~vyvalenýma očima}}
\dicEntry[opineygur] \dicTerm{opin··eygur} \dicIPA{{\textopeno}{\textlengthmark}{\textsubring{b}}{\textsci}{n}{ei}{\textbabygamma}{\textscy}{\textsubring{r}}} \dicPos{adj}[1]\dicFlx{}[-1] \dicLink{opineygður}
\dicEntry[opinmynntur] \dicTerm{opin··mynntur} \dicIPA{{\textopeno}{\textlengthmark}{\textsubring{b}}{\textsci}{n}{m}{\textsci}{\textsubring{n}}{\textsubring{d}}{\textscy}{\textsubring{r}}} \dicPos{adj}[1]\dicFlx{}[-10] \dicDirectTranslationCS{(jsoucí) s~otevřenou pusou, (jsoucí) s~otevřenými ústy}
\dicEntry[opinn] \dicTerm{opinn} \dicsymFrequent\  \dicIPA{{\textopeno}{\textlengthmark}{\textsubring{b}}{\textsci}{\textsubring{n}}} \dicPos{adj}[6]\dicFlx{}[-2] \textbf{1.} \dicSynonym*{opnaður} \dicDirectTranslationCS{otevřený} \dicExampleIS{opinn gluggi} \dicExampleCS{otevřené okno};  \dicPhraseIS{opið hús} \dicDirectTranslationCS{den otevřených dveří}  \textbf{2.} \dicDirectTranslationCS{otevřený (myšlenkám ap.)};  \dicPhraseIS{vera opinn fyrir e‑ð} \dicDirectTranslationCS{být otevřený (čemu)}
\dicEntry[opinskár] \dicTerm{opin··skár} \dicIPA{{\textopeno}{\textlengthmark}{\textsubring{b}}{\textsci}{n}{s}{\r{g}}{au}{\textsubring{r}}} \dicPos{adj}[4]\dicFlx{}[-1] \dicSynonym{hreinskilinn} \dicDirectTranslationCS{otevřený, upřímný} \dicExampleIS{opinská umræða} \dicExampleCS{otevřená diskuze}
\dicEntry[opna] \dicTerm{opn|a\smash{\textsuperscript{1}}} \dicIPA{{\textopeno}{h}{\textsubring{b}}{n}{a}} \dicPos{f}[1] \dicFlx{(‑u, ‑ur)}[7] \dicDirectTranslationCS{dvoustrana, dvojstránka (knihy ap.)}
\dicEntry[opna] \dicTerm{opn|a\smash{\textsuperscript{2}}} \dicsymFrequent\  \dicIPA{{\textopeno}{h}{\textsubring{b}}{n}{a}} \dicPos{v}[1] \dicFlx{(‑aði)}[17] \dicFlx{acc} \dicDirectTranslationCS{otevřít, otvírat, otevírat} \dicExampleIS{opna gluggann} \dicExampleCS{otevřít okno};  \dicIdiom{opnast}{ \dicPhraseIS{opnast}} \dicFlx{refl} \dicDirectTranslationCS{otevřít se, otevírat se} \dicExampleIS{Dyrnar opnast af sjálfu sér.} \dicExampleCS{Dveře se samy otevírají.}
\dicEntry[opnari] \dicTerm{opn··ar|i} \dicIPA{{\textopeno}{h}{\textsubring{b}}{n}{a}{r}{\textsci}} \dicPos{m}[1] \dicFlx{(‑a, ‑ar)}[13] \dicDirectTranslationCS{otvírák, otvírač} \dicExampleIS{dósaopnari} \dicExampleCS{otvírač na konzervy}
\dicEntry[opnun] \dicTerm{opn|un} \dicIPA{{\textopeno}{h}{\textsubring{b}}{n}{\textscy}{\textsubring{n}}} \dicPos{f}[7] \dicFlx{(‑unar)}[9] \textbf{1.} \dicSynonym*{það að opna} \dicDirectTranslationCS{otevření, otevírání}  \textbf{2.} \dicDirectTranslationCS{otevření, zahájení, vernisáž} \dicExampleIS{opnun sýningar} \dicExampleCS{otevření výstavy}
\dicEntry[opnunartími] \dicTerm{opnunar··tím|i} \dicIPA{{\textopeno}{h}{\textsubring{b}}{n}{\textscy}{n}{a}{\textsubring{r}}{t\smash{\textsuperscript{h}}}{i}{m}{\textsci}} \dicPos{m}[1] \dicFlx{(‑a)}[3] \dicDirectTranslationCS{otvírací\,/\addthin otevírací doba} \dicAntonym{lokunartími}
\dicEntry[orð] \dicTerm{orð} \dicsymFrequent\  \dicIPA{{\textopeno}{r}{\texttheta}} \dicPos{n}[2] \dicFlx{(‑s, ‑)}[5] \textbf{1.} \dicDirectTranslationCS{slovo};  \dicPhraseIS{samsett orð} \dicFieldCat{jaz.} \dicDirectTranslationCS{složenina, složené slovo}  \textbf{2.} \dicSynonym{loforð} \dicDirectTranslationCS{slovo, slib};  \dicPhraseIS{halda orð sín} \dicDirectTranslationCS{držet slovo};  \dicPhraseIS{ganga á bak orða sinna} \dicDirectTranslationCS{nedržet slovo}  \textbf{3.} \dicSynonym{ummæli} \dicDirectTranslationCS{poznámka, připomínka};  \dicPhraseIS{fá orð í eyra} \dicDirectTranslationCS{dostat vynadáno, být pokárán};  \dicPhraseIS{segja hálft orð við e‑n} \dicDirectTranslationCS{prohodit s~(kým) (jen) pár slov}  \textbf{4.} \dicSynonym{bón\smash{\textsuperscript{1}}} \dicDirectTranslationCS{prosba, žádost}  \textbf{5.} \dicSynonym{skilaboð} \dicDirectTranslationCS{vzkaz, zpráva} \dicExampleIS{senda e‑m orð} \dicExampleCS{poslat (komu) zprávu}  \textbf{6.} \dicSynonym{tal} \dicDirectTranslationCS{slovo, řeč (před publikem ap.)} \dicExampleIS{fá orðið} \dicExampleCS{dostat slovo}  \textbf{7.} \dicSynonym{orðstír} \dicDirectTranslationCS{reputace, pověst, renomé};  \dicPhraseIS{hafa gott orð á sér} \dicDirectTranslationCS{mít dobrou pověst};  \dicIdiom{orð}{ \dicPhraseIS{biðja um orðið}} \dicDirectTranslationCS{požádat o~slovo};  \dicPhraseIS{eiga síðasta orðið} \dicLangCat{přen.} \dicDirectTranslationCS{mít poslední slovo};  \dicPhraseIS{e‑ir skiptast á orðum} \dicFlx{refl} \dicLangCat{přen.} \dicDirectTranslationCS{(kdo) si vyměňuje pár slov}; { \dicPhraseIS{ekki er orðum að því eyðandi}} \dicLangCat{přen.} \dicDirectTranslationCS{to nestojí za řeč};  \dicPhraseIS{eyða ekki orðum að e‑u} \dicLangCat{přen.} \dicDirectTranslationCS{neztrácet čas (čím), nepovažovat za důležité se (čím) zabývat};  \dicPhraseIS{fara hörðum orðum um e‑n} \dicDirectTranslationCS{tvrdě (koho) kritizovat};  \dicPhraseIS{fá orð á sig fyrir e‑ð} \dicDirectTranslationCS{být pověstný (čím)};  \dicPhraseIS{gæta orða sinna} \dicDirectTranslationCS{vážit svá slova; dávat si pozor na to, co člověk říká};  \dicPhraseIS{hafa orð á e‑u} \dicDirectTranslationCS{zmínit se o~(čem)};  \dicPhraseIS{í einu orði sagt} \dicFlx{adv} \dicDirectTranslationCS{jedním slovem, krátce řečeno};  \dicPhraseIS{í orðsins fyllstu merkingu} \dicFlx{adv} \dicDirectTranslationCS{v~pravém významu slova};  \dicPhraseIS{koma ekki upp orði} \dicLangCat{přen.} \dicDirectTranslationCS{nenacházet slov};  \dicPhraseIS{koma orðum að e‑u} \dicDirectTranslationCS{najít způsob, jak (co) vyjádřit};  \dicPhraseIS{komast (þannig) að orði} \dicFlx{refl} \dicDirectTranslationCS{vyjádřit to (tak)};  \dicPhraseIS{láta ekki sitja við orðin tóm} \dicLangCat{přen.} \dicDirectTranslationCS{nezůstat u~prázdných slov};  \dicPhraseIS{leggja orð í belg um e‑ð} \dicDirectTranslationCS{poznamenat (co)};  \dicPhraseIS{meðal annarra orða} \dicFlx{adv} \dicDirectTranslationCS{mimo jiné, mimochodem};  \dicPhraseIS{með öðrum orðum} \dicFlx{adv} \dicDirectTranslationCS{jinými slovy, jinak řečeno};  \dicPhraseIS{orð í tíma töluð} \dicDirectTranslationCS{poznámka v~pravý čas};  \dicPhraseIS{orðum aukinn} \dicFlx{adj} \dicDirectTranslationCS{zveličený, přehnaný};  \dicPhraseIS{segja e‑ð berum orðum} \dicDirectTranslationCS{říct (co) na rovinu};  \dicPhraseIS{skipta orðum við e‑n} \dicDirectTranslationCS{vyměnit\,/\addthin prohodit s~(kým) pár slov};  \dicPhraseIS{taka aftur orð sín} \dicLangCat{přen.} \dicDirectTranslationCS{vzít\,/\addthin brát to zpět} \dicIndirectTranslationCS{(vzít zpět to, co člověk řekl)};  \dicPhraseIS{taka e‑n á orðinu} \dicLangCat{přen.} \dicDirectTranslationCS{brát (koho) za slovo};  \dicPhraseIS{taka (svona) til orða} \dicDirectTranslationCS{vyjádřit se (takto)};  \dicProverb\  \dicPhraseIS{Töluð orð verða ekki aftur tekin.} \dicLangCat{přís.} \dicDirectTranslationCS{Co jsi z~úst vypustil, ani párem koní nedostaneš zpět.}
\dicEntry[orða] \dicTerm{orð|a\smash{\textsuperscript{1}}} \dicIPA{{\textopeno}{r}{ð}{a}} \dicPos{f}[1] \dicFlx{(‑u, ‑ur)}[19] \dicDirectTranslationCS{řád, vyznamenání}
\dicEntry[orða] \dicTerm{orð|a\smash{\textsuperscript{2}}} \dicIPA{{\textopeno}{r}{ð}{a}} \dicPos{v}[1] \dicFlx{(‑aði)}[1] \textbf{1.} \dicSynonym*{koma orðum að} \dicDirectTranslationCS{vyjádřit, vyjadřovat, (z)formulovat} \dicExampleIS{orða hugsun sína skýrt} \dicExampleCS{vyjádřit jasně své myšlenky}  \textbf{2.} \dicSynonym{nefna} \dicDirectTranslationCS{zmínit (se)} \dicExampleIS{orða e‑ð við e‑n} \dicExampleCS{zmínit se (komu) o~(čem)}  \textbf{3.} \dicPhraseIS{orða e‑n við e‑ð} \dicDirectTranslationCS{spojovat (koho) s~(čím) (řidiče s~nehodou ap.)}
\dicEntry[orðabil] \dicTerm{orða··bil} \dicIPA{{\textopeno}{r}{ð}{a}{\textsubring{b}}{\textsci}{\textsubring{l}}} \dicPos{n}[2] \dicFlx{(‑s, ‑)}[5] \dicDirectTranslationCS{slovní mezera, mezera mezi slovy}
\dicEntry[orðabók] \dicTerm{orða··|bók} \dicIPA{{\textopeno}{r}{ð}{a}{\textsubring{b}}{ou}{\r{g}}} \dicPos{f}[8] \dicFlx{(‑bókar, ‑bækur)}[5] \dicDirectTranslationCS{slovník, lexikon} \dicExampleIS{semja orðabók} \dicExampleCS{psát slovník}
\dicEntry[orðabókarfræði] \dicTerm{orða·bókar··fræð|i} \dicIPA{{\textopeno}{r}{ð}{a}{\textsubring{b}}{ou}{\r{g}}{a}{\textsubring{r}}{f}{r}{a}{i}{ð}{\textsci}} \dicPos{f}[3] \dicFlx{(‑i)}[3] \dicFieldCat{jaz.} \dicDirectTranslationCS{lexikografie, slovníkářství}
\dicEntry[orðabókarhöfundur] \dicTerm{orða·bókar··höf·und|ur} \dicIPA{{\textopeno}\-{r}\-{ð}\-{a}\-{\textsubring{b}}\-{ou}\-{\r{g}}\-{a}\-{\textsubring{r}}\-{h}\-{\oe}\-{v}\-{\textscy}\-{n}\-{\textsubring{d}}\-{\textscy}\-{\textsubring{r}}\-} \dicPos{m}[6] \dicFlx{(‑ar, ‑ar)}[56] \dicFieldCat{jaz.} \dicDirectTranslationCS{lexikograf(ka), slovníkář(ka)}
\dicEntry[orðafar] \dicTerm{orða··far} \dicIPA{{\textopeno}{r}{ð}{a}{f}{a}{\textsubring{r}}} \dicPos{n}[2] \dicFlx{(‑s)}[2] \dicDirectTranslationCS{jazykový styl, styl vyjadřování}
\dicEntry[orðaforði] \dicTerm{orða··forð|i} \dicIPA{{\textopeno}{r}{ð}{a}{f}{\textopeno}{r}{ð}{\textsci}} \dicPos{m}[1] \dicFlx{(‑a)}[3] \dicFieldCat{jaz.} \dicDirectTranslationCS{slovní zásoba}
\dicEntry[orðagjálfur] \dicTerm{orða··gjálfur} \dicIPA{{\textopeno}{r}{ð}{a}{\r{\textObardotlessj}}{au}{l}{v}{\textscy}{\textsubring{r}}} \dicPos{n}[2] \dicFlx{(‑s)}[28] \dicSynonym{þvaður} \dicDirectTranslationCS{příval\,/\addthin proud slov, tiráda}
\dicEntry[orðalag] \dicTerm{orða··lag} \dicIPA{{\textopeno}{r}{ð}{a}{l}{a}{x}} \dicPos{n}[2] \dicFlx{(‑s)}[2] \dicDirectTranslationCS{vyjádření, formulace}
\dicEntry[orðaleikur] \dicTerm{orða··leik|ur} \dicIPA{{\textopeno}{r}{ð}{a}{l}{ei}{\r{g}}{\textscy}{\textsubring{r}}} \dicPos{m}[9] \dicFlx{(‑s, ‑ir)}[15] \dicDirectTranslationCS{slovní hra\,/\addthin hříčka, hra se slovy}
\dicEntry[orðalykill] \dicTerm{orða··lyk|ill} \dicIPA{{\textopeno}{r}{ð}{a}{l}{\textsci}{\r{\textObardotlessj}}{\textsci}{\textsubring{d}}{\textsubring{l}}} \dicPos{m}[6] \dicFlx{(‑ils, ‑lar)}[35] \dicSynonym{orðstöðulykill} \dicDirectTranslationCS{konkordance}
\dicEntry[orðamunur] \dicTerm{orða··mun|ur} \dicIPA{{\textopeno}{r}{ð}{a}{m}{\textscy}{n}{\textscy}{\textsubring{r}}} \dicPos{m}[10] \dicFlx{(‑ar)}[12] \dicDirectTranslationCS{slovní obměna\,/\addthin varianta}
\dicEntry[orðaröð] \dicTerm{orða··|röð} \dicIPA{{\textopeno}{r}{ð}{a}{r}{\oe}{\texttheta}} \dicPos{f}[7] \dicFlx{(‑raðar)}[19] \dicFieldCat{jaz.} \dicDirectTranslationCS{slovosled, pořádek slov} \dicExampleIS{grundvallarorðaröð} \dicExampleCS{základní slovosled}
\dicEntry[orðasafn] \dicTerm{orða··|safn} \dicIPA{{\textopeno}{r}{ð}{a}{s}{a}{\textsubring{b}}{\textsubring{n}}} \dicPos{n}[2] \dicFlx{(‑safns, ‑söfn)}[8] \dicDirectTranslationCS{glosář, seznam slov, slovník}
\dicEntry[orðasamband] \dicTerm{orða··sam·|band} \dicIPA{{\textopeno}{r}{ð}{a}{s}{a}{m}{\textsubring{b}}{a}{n}{\textsubring{d}}} \dicPos{n}[2] \dicFlx{(‑bands, ‑bönd)}[8] \dicFieldCat{jaz.} \dicDirectTranslationCS{úsloví, fráze, (ustálené) slovní spojení}
\dicEntry[orðaskil] \dicTerm{orða··skil} \dicIPA{{\textopeno}{r}{ð}{a}{s}{\r{\textObardotlessj}}{\textsci}{\textsubring{l}}} \dicPos{n}[2] \dicFlx{pl}[1] \dicSynonym*{orðagrein} \dicDirectTranslationCS{hranice slov};  \dicPhraseIS{heyra ekki orðaskil} \dicDirectTranslationCS{nerozeznávat jednotlivá slova}
\dicEntry[orðaskipti] \dicTerm{orða··skipti} \dicIPA{{\textopeno}{r}{ð}{a}{s}{\r{\textObardotlessj}}{\textsci}{f}{\textsubring{d}}{\textsci}} \dicPos{n}[2] \dicFlx{pl}[19] \dicDirectTranslationCS{výměna slov, rozhovor};  \dicPhraseIS{eiga orðaskipti við e‑n} \dicDirectTranslationCS{rozmlouvat s~(kým), vyměnit s~(kým) pár slov}
\dicEntry[orðatiltæki] \dicTerm{orða··til·tæki} \dicIPA{{\textopeno}{r}{ð}{a}{t\smash{\textsuperscript{h}}}{\textsci}{\textsubring{l}}{t\smash{\textsuperscript{h}}}{a}{i}{\r{\textObardotlessj}}{\textsci}} \dicPos{n}[2] \dicFlx{(‑s, ‑)}[16] \dicFieldCat{jaz.} \dicDirectTranslationCS{idiom, fráze, úsloví, (ustálené) slovní spojení, slovní obrat}
\dicEntry[orðaval] \dicTerm{orða··val} \dicIPA{{\textopeno}{r}{ð}{a}{v}{a}{\textsubring{l}}} \dicPos{n}[2] \dicFlx{(‑s)}[2] \dicDirectTranslationCS{volba\,/\addthin výběr slov}
\dicEntry[orðblinda] \dicTerm{orð··blind|a} \dicIPA{{\textopeno}{r}{ð}{\textsubring{b}}{l}{\textsci}{n}{\textsubring{d}}{a}} \dicPos{f}[1] \dicFlx{(‑u)}[5] \dicFieldCat{med.} \dicSynonym{lesblinda} \dicDirectTranslationCS{dyslexie}
\dicEntry[orðbragð] \dicTerm{orð··bragð} \dicIPA{{\textopeno}{r}{ð}{\textsubring{b}}{r}{a}{\textbabygamma}{\texttheta}} \dicPos{n}[2] \dicFlx{(‑s)}[2] \dicDirectTranslationCS{výrazivo, vyjadřování, (vyjadřovací) jazyk} \dicExampleIS{ljótt orðbragð} \dicExampleCS{ošklivé vyjadřování}
\dicEntry[orðfall] \dicTerm{orð··fall} \dicIPA{{\textopeno}{r}{ð}{f}{a}{\textsubring{d}}{\textsubring{l}}} \dicPos{n}[2] \dicFlx{(‑s)}[2] \dicPhraseIS{e‑m verður orðfall} \dicFlx{impers} \dicDirectTranslationCS{(komu) docházejí slova}
\dicEntry[orðfimi] \dicTerm{orð··fim|i} \dicIPA{{\textopeno}{r}{ð}{f}{\textsci}{m}{\textsci}} \dicPos{f}[3] \dicFlx{(‑i)}[3] \dicDirectTranslationCS{výmluvnost, výřečnost}
\dicEntry[orðfimur] \dicTerm{orð··fimur} \dicIPA{{\textopeno}{r}{ð}{f}{\textsci}{m}{\textscy}{\textsubring{r}}} \dicPos{adj}[1]\dicFlx{}[-1] \dicDirectTranslationCS{výřečný, výmluvný}
\dicEntry[orðflokkur] \dicTerm{orð··flokk|ur} \dicIPA{{\textopeno}{r}{ð}{f}{l}{\textopeno}{h}{\r{g}}{\textscy}{\textsubring{r}}} \dicPos{m}[6] \dicFlx{(‑s, ‑ar)}[8] \dicFieldCat{jaz.} \dicDirectTranslationCS{slovní druh}
\dicEntry[orðfæri] \dicTerm{orð··færi} \dicIPA{{\textopeno}{r}{ð}{f}{a}{i}{r}{\textsci}} \dicPos{n}[2] \dicFlx{(‑s)}[20] \dicSynonym{orðalag} \dicDirectTranslationCS{styl\,/\addthin způsob vyjadřování}
\dicEntry[orðheldinn] \dicTerm{orð··heldinn} \dicIPA{{\textopeno}{r}{\texttheta}{h}{\textepsilon}{l}{\textsubring{d}}{\textsci}{\textsubring{n}}} \dicPos{adj}[6]\dicFlx{}[-2] \dicDirectTranslationCS{držící slovo} \dicExampleIS{vera orðheldinn} \dicExampleCS{držet slovo}
\dicEntry[orðheldni] \dicTerm{orð··heldn|i} \dicIPA{{\textopeno}{r}{\texttheta}{h}{\textepsilon}{l}{\textsubring{d}}{n}{\textsci}} \dicPos{f}[3] \dicFlx{(‑i)}[3] \dicDirectTranslationCS{důvěryhodnost\,/\addthin spolehlivost (v~držení slova)}
\dicEntry[orðheppinn] \dicTerm{orð··heppinn} \dicIPA{{\textopeno}{r}{\texttheta}{h}{\textepsilon}{h}{\textsubring{b}}{\textsci}{\textsubring{n}}} \dicPos{adj}[6]\dicFlx{}[-2] \dicSynonym{orðfimur} \dicDirectTranslationCS{výřečný, výmluvný}
\dicEntry[orðhlutafræði] \dicTerm{orð·hluta··fræð|i} \dicIPA{{\textopeno}{r}{ð}{\textsubring{l}}{\textscy}{\textsubring{d}}{a}{f}{r}{a}{i}{ð}{\textsci}} \dicPos{f}[3] \dicFlx{(‑i)}[3] \dicFieldCat{jaz.} \dicDirectTranslationCS{morfologie, tvarosloví}
\dicEntry[orðhvatur] \dicTerm{orð··|hvatur} \dicIPA{{\textopeno}{r}{ð}{k\smash{\textsuperscript{h}}}{v}{a}{\textsubring{d}}{\textscy}{\textsubring{r}}} \dicPos{adj}[1] \dicFlx{(f ‑hvöt)}[2] \dicDirectTranslationCS{prostořeký, hubatý}
\dicEntry[orðið] \dicTerm{orðið} \dicIPA{{\textopeno}{r}{ð}{\textsci}{\texttheta}} \dicPos{v} \dicFlx{supin} \dicLink{verða}
\dicEntry[orðlagður] \dicTerm{orð··|lagður} \dicIPA{{\textopeno}{r}{ð}{l}{a}{\textbabygamma}{ð}{\textscy}{\textsubring{r}}} \dicPos{adj}[2] \dicFlx{(f ‑lögð)}[3] \dicSynonym{frægur} \dicDirectTranslationCS{proslulý, věhlasný}
\dicEntry[orðlaus] \dicTerm{orð··laus} \dicIPA{{\textopeno}{r}{ð}{l}{\oe i}{s}} \dicPos{adj}[5]\dicFlx{}[-1] \dicDirectTranslationCS{nemající slov, němý} \dicExampleIS{Ég er orðlaus.} \dicExampleCS{Nemám slov.}
\dicEntry[orðlengja] \dicTerm{orð··leng|ja} \dicIPA{{\textopeno}{r}{ð}{l}{ei}{\textltailn}{\r{\textObardotlessj}}{a}} \dicPos{v}[2] \dicFlx{(‑di, ‑t)}[137] \dicFlx{acc} \dicDirectTranslationCS{mluvit zdlouhavě, protahovat (vyprávění ap.)} \dicExampleIS{orðlengja (um) e‑ð} \dicExampleCS{protahovat (co)}
\dicEntry[orðliður] \dicTerm{orð··lið|ur} \dicIPA{{\textopeno}{r}{ð}{l}{\textsci}{ð}{\textscy}{\textsubring{r}}} \dicPos{m}[10] \dicFlx{(‑ar\,/\addthin ‑s, ‑ir)}[33] \dicFieldCat{jaz.} \dicDirectTranslationCS{člen kompozita\,/\addthin složeniny}
\dicEntry[orðljótur] \dicTerm{orð··ljótur} \dicIPA{{\textopeno}{r}{ð}{l}{j}{ou}{\textsubring{d}}{\textscy}{\textsubring{r}}} \dicPos{adj}[1]\dicFlx{}[-1] \dicDirectTranslationCS{hrubý, sprostý (ve vyjadřování)}
\dicEntry[orðmargur] \dicTerm{orð··|margur} \dicIPA{{\textopeno}{r}{ð}{m}{a}{r}{\r{g}}{\textscy}{\textsubring{r}}} \dicPos{adj}[11] \dicFlx{(f ‑mörg, comp ‑fleiri, sup ‑flestur)}[8] \textbf{1.} \dicSynonym*{sem inniheldur mörg orð} \dicDirectTranslationCS{rozvláčný (text ap.)}  \textbf{2.} \dicSynonym*{sem notar mörg orð} \dicDirectTranslationCS{mnohomluvný, upovídaný}
\dicEntry[orðmynd] \dicTerm{orð··mynd} \dicIPA{{\textopeno}{r}{ð}{m}{\textsci}{n}{\textsubring{d}}} \dicPos{f}[7] \dicFlx{(‑ar, ‑ir)}[1] \dicFieldCat{jaz.} \dicDirectTranslationCS{slovní tvar\,/\addthin forma}
\dicEntry[orðmyndun] \dicTerm{orð··mynd|un} \dicIPA{{\textopeno}{r}{ð}{m}{\textsci}{n}{\textsubring{d}}{\textscy}{\textsubring{n}}} \dicPos{f}[7] \dicFlx{(‑unar)}[9] \dicFieldCat{jaz.} \dicDirectTranslationCS{tvoření slov, slovotvorba}
\dicEntry[orðmyndunarfræði] \dicTerm{orð·myndunar··fræð|i} \dicIPA{{\textopeno}{r}{ð}{m}{\textsci}{n}{\textsubring{d}}{\textscy}{n}{a}{\textsubring{r}}{f}{r}{a}{i}{ð}{\textsci}} \dicPos{f}[3] \dicFlx{(‑i)}[3] \dicFieldCat{jaz.} \dicDirectTranslationCS{slovotvorba}
\dicEntry[orðrétt] \dicTerm{orð··rétt} \dicIPA{{\textopeno}{r}{ð}{r}{j}{\textepsilon}{h}{\textsubring{d}}} \dicPos{adv} \dicDirectTranslationCS{doslovně} \dicExampleIS{skrifa e‑ð orðrétt upp} \dicExampleCS{opsat (co) doslovně}
\dicEntry[orðréttur] \dicTerm{orð··réttur} \dicIPA{{\textopeno}{r}{ð}{r}{j}{\textepsilon}{h}{\textsubring{d}}{\textscy}{\textsubring{r}}} \dicPos{adj}[1]\dicFlx{}[-10] \dicDirectTranslationCS{doslovný} \dicExampleIS{orðrétt þýðing} \dicExampleCS{doslovný překlad}
\dicEntry[orðrómur] \dicTerm{orð··róm|ur} \dicIPA{{\textopeno}{r}{ð}{r}{ou}{m}{\textscy}{\textsubring{r}}} \dicPos{m}[6] \dicFlx{(‑s, ‑ar)}[24] \dicDirectTranslationCS{fáma, řeči (neověřené)};  \dicPhraseIS{liggja undir orðrómi} \dicLangCat{přen.} \dicDirectTranslationCS{přijít do řečí}
\dicEntry[orðsending] \dicTerm{orð··send·ing} \dicIPA{{\textopeno}{r}{ð}{s}{\textepsilon}{n}{\textsubring{d}}{i}{\ng}{\r{g}}} \dicPos{f}[4] \dicFlx{(‑ar, ‑ar)}[5] \dicSynonym{boðskapur} \dicDirectTranslationCS{zprá\-va, vzkaz, depeše, sdělení} \dicExampleIS{orðsending um rafmagnsbilun} \dicExampleCS{zpráva o~poruše dodávky elektřiny}
\dicEntry[orðsifjafræði] \dicTerm{orð·sifja··fræð|i} \dicIPA{{\textopeno}{r}{ð}{s}{\textsci}{v}{j}{a}{f}{r}{a}{i}{ð}{\textsci}} \dicPos{f}[3] \dicFlx{(‑i)}[3] \dicFieldCat{jaz.} \dicDirectTranslationCS{etymologie} \dicIndirectTranslationCS{(nauka)}
\dicEntry[orðsifjafræðingur] \dicTerm{orð·sifja·fræð··ing|ur} \dicIPA{{\textopeno}\-{r}\-{ð}\-{s}\-{\textsci}\-{v}\-{j}\-{a}\-{f}\-{r}\-{a}\-{i}\-{ð}\-{i}\-{\ng}\-{\r{g}}\-{\textscy}\-{\textsubring{r}}\-} \dicPos{m}[6] \dicFlx{(‑s, ‑ar)}[8] \dicDirectTranslationCS{etymolog, etymoložka}
\dicEntry[orðsifjar] \dicTerm{orð··sifjar} \dicIPA{{\textopeno}{r}{ð}{s}{\textsci}{v}{j}{a}{\textsubring{r}}} \dicPos{f}[12] \dicFlx{pl}[1] \dicFieldCat{jaz.} \dicDirectTranslationCS{etymologie (slova ap.)}
\dicEntry[orðskipan] \dicTerm{orð··skipan} \dicIPA{{\textopeno}{r}{ð}{s}{\r{\textObardotlessj}}{\textsci}{\textsubring{b}}{a}{\textsubring{n}}} \dicPos{f}[7] \dicFlx{(‑ar)}[3] \dicFieldCat{jaz.} \dicSynonym{orðalag} \dicDirectTranslationCS{formulace, vyjádření}
\dicEntry[orðskipting] \dicTerm{orð··skipt·ing} \dicIPA{{\textopeno}{r}{ð}{s}{\r{\textObardotlessj}}{\textsci}{f}{\textsubring{d}}{i}{\ng}{\r{g}}} \dicPos{f}[4] \dicFlx{(‑ar, ‑ar)}[5] \dicFieldCat{jaz.} \dicDirectTranslationCS{dělení slov} \dicIndirectTranslationCS{(mezi řádky v~textu)}
\dicEntry[orðskviður] \dicTerm{orðs··kvið|ur} \dicIPA{{\textopeno}{r}{ð}{s}{k\smash{\textsuperscript{h}}}{v}{\textsci}{ð}{\textscy}{\textsubring{r}}} \dicPos{m}[10] \dicFlx{(‑ar, ‑ir)}[19] \dicSynonym{málsháttur} \dicDirectTranslationCS{pořekadlo, aforismus}
\dicEntry[orðskýring] \dicTerm{orð··skýr·ing} \dicIPA{{\textopeno}{r}{ð}{s}{\r{\textObardotlessj}}{i}{r}{i}{\ng}{\r{g}}} \dicPos{f}[4] \dicFlx{(‑ar, ‑ar)}[5] \dicFieldCat{jaz.} \dicDirectTranslationCS{slovní význam\,/\addthin výklad}
\dicEntry[orðspor] \dicTerm{orð··spor} \dicIPA{{\textopeno}{r}{ð}{s}{\textsubring{b}}{\textopeno}{\textsubring{r}}} \dicPos{n}[2] \dicFlx{(‑s)}[2] \dicDirectTranslationCS{pověst, reputace}
\dicEntry[orðstír] \dicTerm{orðs··tír} \dicIPA{{\textopeno}{r}{ð}{s}{\textsubring{d}}{i}{\textsubring{r}}} \dicPos{m}[4] \dicFlx{(‑s)}[18] \dicSynonym{frægð} \dicDirectTranslationCS{proslulost, věhlas, renomé, (dobré) jméno} \dicExampleIS{Orðstír hennar barst alla leið til Ameríku.} \dicExampleCS{Její věhlas se donesl až do Ameriky.}
\dicEntry[orðstofn] \dicTerm{orð··stofn} \dicIPA{{\textopeno}{r}{ð}{s}{\textsubring{d}}{\textopeno}{\textsubring{b}}{\textsubring{n}}} \dicPos{m}[4] \dicFlx{(‑s, ‑ar)}[3] \dicFieldCat{jaz.} \dicDirectTranslationCS{(slovní) kmen}
\dicEntry[orðstöðulykill] \dicTerm{orð·stöðu··lyk|ill} \dicIPA{{\textopeno}{r}{ð}{s}{\textsubring{d}}{\oe}{ð}{\textscy}{l}{\textsci}{\r{\textObardotlessj}}{\textsci}{\textsubring{d}}{\textsubring{l}}} \dicPos{m}[6] \dicFlx{(‑ils, ‑lar)}[35] \dicDirectTranslationCS{konkordance} \dicExampleIS{orðstöðulykill Íslendingasagna} \dicExampleCS{konkordance islandských ság}
\dicEntry[orðtak] \dicTerm{orð··|tak} \dicIPA{{\textopeno}{r}{ð}{t\smash{\textsuperscript{h}}}{a}{\r{g}}} \dicPos{n}[2] \dicFlx{(‑taks, ‑tök)}[8] \dicFieldCat{jaz.} \dicSynonym{frasi} \dicDirectTranslationCS{fráze, slovní obrat, idiom}
\dicEntry[orðtaka] \dicTerm{orð··|taka} \dicIPA{{\textopeno}{r}{ð}{t\smash{\textsuperscript{h}}}{a}{\r{g}}{a}} \dicPos{v}[6] \dicFlx{(‑tek, ‑tók, ‑tókum, ‑tæki, ‑tekið)}[128] \dicDirectTranslationCS{sbírat\,/\addthin vybírat slova} \dicIndirectTranslationCS{(za účelem jejich vložení do slovníku ap.)}
\dicEntry[orðtek] \dicTerm{orð··tek} \dicIPA{{\textopeno}{r}{ð}{t\smash{\textsuperscript{h}}}{\textepsilon}{\r{g}}} \dicPos{v} \dicFlx{ind praes sg 1 pers} \dicLink{orðtaka}
\dicEntry[orðtekið] \dicTerm{orð··tekið} \dicIPA{{\textopeno}{r}{ð}{t\smash{\textsuperscript{h}}}{\textepsilon}{\r{\textObardotlessj}}{\textsci}{\texttheta}} \dicPos{v} \dicFlx{supin} \dicLink{orðtaka}
\dicEntry[orðtók] \dicTerm{orð··tók} \dicIPA{{\textopeno}{r}{ð}{t\smash{\textsuperscript{h}}}{ou}{\r{g}}} \dicPos{v} \dicFlx{ind pf sg 1 pers} \dicLink{orðtaka}
\dicEntry[orðtókum] \dicTerm{orð··tókum} \dicIPA{{\textopeno}{r}{ð}{t\smash{\textsuperscript{h}}}{ou}{\r{g}}{\textscy}{\textsubring{m}}} \dicPos{v} \dicFlx{ind pf pl 1 pers} \dicLink{orðtaka}
\dicEntry[orðtæki] \dicTerm{orð··tæki} \dicIPA{{\textopeno}{r}{ð}{t\smash{\textsuperscript{h}}}{a}{i}{\r{\textObardotlessj}}{\textsci}} \dicPos{v} \dicFlx{con pf sg 1 pers} \dicLink{orðtaka}
\dicEntry[orðvar] \dicTerm{orð··|var} \dicIPA{{\textopeno}{r}{ð}{v}{a}{\textsubring{r}}} \dicPos{adj}[5] \dicFlx{(f ‑vör)}[9] \dicDirectTranslationCS{opatrný\,/\addthin obezřetný (ve výběru slov)}
\dicEntry[orf] \dicTerm{orf} \dicIPA{{\textopeno}{r}{f}} \dicPos{n}[2] \dicFlx{(‑s, ‑)}[5] \dicDirectTranslationCS{kosiště, násada kosy}
\dicEntry[org] \dicTerm{org} \dicIPA{{\textopeno}{r}{\r{g}}} \dicPos{n}[2] \dicFlx{(‑s, ‑)}[5] \textbf{1.} \dicSynonym{óp} \dicDirectTranslationCS{řev, (za)řvaní, jekot, (za)ječení}  \textbf{2.} \dicSynonym{grátur\smash{\textsuperscript{1}}} \dicDirectTranslationCS{(silný) pláč, řev, řvaní}
\dicEntry[orga] \dicTerm{org|a} \dicIPA{{\textopeno}{r}{\r{g}}{a}} \dicPos{v}[1] \dicFlx{(‑aði)}[34] \textbf{1.} \dicSynonym{æpa} \dicDirectTranslationCS{řvát, zařvat, (za)ječet}  \textbf{2.} \dicSynonym{gráta} \dicDirectTranslationCS{řvát, (hlasitě) brečet}
\dicEntry[organisti] \dicTerm{organ··ist|i} \dicIPA{{\textopeno}{r}{\r{g}}{a}{n}{\textsci}{s}{\textsubring{d}}{\textsci}} \dicPos{m}[1] \dicFlx{(‑a, ‑ar)}[1] \dicDirectTranslationCS{varhaník, varhanice}
\dicEntry[orgel] \dicTerm{orgel} \dicIPA{{\textopeno}{r}{\r{\textObardotlessj}}{\textepsilon}{\textsubring{l}}} \dicPos{n}[2] \dicFlx{(‑s, ‑)}[5] \dicFieldCat{hud.} \dicDirectTranslationCS{varhany}
\dicEntry[orka] \dicTerm{ork|a\smash{\textsuperscript{1}}} \dicsymFrequent\  \dicIPA{{\textopeno}{\textsubring{r}}{\r{g}}{a}} \dicPos{f}[1] \dicFlx{(‑u)}[5] \textbf{1.} \dicSynonym{kraftur} \dicDirectTranslationCS{energie, síla}  \textbf{2.} \dicFieldCat{fyz.} \dicDirectTranslationCS{energie} \dicExampleIS{nýta orku} \dicExampleCS{využívat energii}
\dicEntry[orka] \dicTerm{ork|a\smash{\textsuperscript{2}}} \dicIPA{{\textopeno}{\textsubring{r}}{\r{g}}{a}} \dicPos{v}[1] \dicFlx{(‑aði)}[1] \textbf{1.} \dicSynonym{megna} \dicDirectTranslationCS{zvládnout, mít sílu} \dicExampleIS{orka að (gera e‑ð)} \dicExampleCS{zvládnout ((co) udělat)}  \textbf{2.} \dicSynonym{verka} \dicDirectTranslationCS{působit, mít vliv} \dicExampleIS{orka á e‑ð\,/\addthin e‑n} \dicExampleCS{působit na (co\,/\addthin koho)}  \textbf{3.} \dicPhraseIS{e‑að orkar tvímælis} \dicDirectTranslationCS{(co) je diskutabilní}
\dicEntry[Orkneyjar] \dicTerm{Orkn··eyjar} \dicIPA{{\textopeno}{\textsubring{r}}{\r{g}}{n}{ei}{j}{a}{\textsubring{r}}} \dicPos{f}[4] \dicFlx{pl}[15] \dicFieldCat{geog.} \dicDirectTranslationCS{Orkneje} \dicIndirectTranslationCS{(skotské souostroví)}
\dicEntry[orkugjafi] \dicTerm{orku··gjaf|i} \dicIPA{{\textopeno}{\textsubring{r}}{\r{g}}{\textscy}{\r{\textObardotlessj}}{a}{v}{\textsci}} \dicPos{m}[1] \dicFlx{(‑a, ‑ar)}[8] \dicDirectTranslationCS{zdroj energie}
\dicEntry[orkuleysi] \dicTerm{orku··leysi} \dicIPA{{\textopeno}{\textsubring{r}}{\r{g}}{\textscy}{l}{ei}{s}{\textsci}} \dicPos{n}[2] \dicFlx{(‑s)}[20] \dicDirectTranslationCS{nedostatek energie, netečnost}
\dicEntry[orkunýting] \dicTerm{orku··nýt·ing} \dicIPA{{\textopeno}{\textsubring{r}}{\r{g}}{\textscy}{n}{i}{\textsubring{d}}{i}{\ng}{\r{g}}} \dicPos{f}[4] \dicFlx{(‑ar)}[7] \dicDirectTranslationCS{využití\,/\addthin využívání energie}
\dicEntry[orkuver] \dicTerm{orku··ver} \dicIPA{{\textopeno}{\textsubring{r}}{\r{g}}{\textscy}{v}{\textepsilon}{\textsubring{r}}} \dicPos{n}[2] \dicFlx{(‑s, ‑)}[5] \dicDirectTranslationCS{elektrárna} \dicExampleIS{kjarnorkuver} \dicExampleCS{atomová elektrárna}
\dicEntry[orlof] \dicTerm{or··lof} \dicIPA{{\textopeno}{r}{l}{\textopeno}{f}} \dicPos{n}[2] \dicFlx{(‑s, ‑)}[5] \dicSynonym{frí} \dicDirectTranslationCS{prázdniny, dovolená} \dicExampleIS{fá orlof} \dicExampleCS{dostat dovolenou}
\dicEntry[ormur] \dicTerm{orm|ur} \dicIPA{{\textopeno}{r}{m}{\textscy}{\textsubring{r}}} \dicPos{m}[6] \dicFlx{(‑s, ‑ar)}[4] \textbf{1.} \dicFieldCat{zool.} \dicDirectTranslationCS{červ} \textit{(l.~{\textLA{Vermes}})}  \textbf{2.} \dicSynonym{slanga} \dicDirectTranslationCS{had}
\dicEntry[orna] \dicTerm{orn|a} \dicIPA{{\textopeno}{r}{\textsubring{d}}{n}{a}} \dicPos{v}[1] \dicFlx{(‑aði)}[1] \dicFlx{dat} \dicSynonym{hlýja\smash{\textsuperscript{2}}} \dicDirectTranslationCS{(za)hřát, ohřát} \dicExampleIS{orna e‑m} \dicExampleCS{zahřát (koho)};  \dicPhraseIS{orna sér} \dicDirectTranslationCS{zahřát se}
\dicEntry[orpið] \dicTerm{orpið} \dicIPA{{\textopeno}{\textsubring{r}}{\textsubring{b}}{\textsci}{\texttheta}} \dicPos{v} \dicFlx{supin} \dicLink{verpa\smash{\textsuperscript{2}}}
\dicEntry[orri] \dicTerm{orr|i} \dicIPA{{\textopeno}{r}{\textlengthmark}{\textsci}} \dicPos{m}[1] \dicFlx{(‑a, ‑ar)}[1] \dicFieldCat{zool.} \dicDirectTranslationCS{tetřev, tetřívek obecný} \textit{(l.~{\textLA{Tetrao tetrix}})}  \dicsymPhoto\ 
\dicFigure{ds_image_orri_0_1.jpg}{Orri}{Orri - Vnp, CC BY-SA 3.0}
\dicEntry[orrusta] \dicTerm{orr··ust|a} \dicIPA{{\textopeno}{r}{\textlengthmark}{\textscy}{s}{\textsubring{d}}{a}} \dicPos{f}[1] \dicFlx{(‑u, ‑ur)}[19] \dicLink{orusta}
\dicEntry[orsaka] \dicTerm{or··sak|a} \dicsymFrequent\  \dicIPA{{\textopeno}{\textsubring{r}}{s}{a}{\r{g}}{a}} \dicPos{v}[1] \dicFlx{(‑aði)}[13] \dicFlx{acc} \dicSynonym{framkalla} \dicDirectTranslationCS{způsobit, zapříčinit, přivodit} \dicExampleIS{Þetta getur orsakað slys.} \dicExampleCS{To může způsobit nehodu.};  \dicIdiom{orsakast}{ \dicPhraseIS{orsakast}} \dicFlx{refl} \dicDirectTranslationCS{být způsoben}
\dicEntry[orsakafræði] \dicTerm{or·saka··fræð|i} \dicIPA{{\textopeno}{\textsubring{r}}{s}{a}{\r{g}}{a}{f}{r}{a}{i}{ð}{\textsci}} \dicPos{f}[3] \dicFlx{(‑i)}[3] \dicDirectTranslationCS{etiologie}
\dicEntry[orsakalögmál] \dicTerm{or·saka··lög·mál} \dicIPA{{\textopeno}{\textsubring{r}}{s}{a}{\r{g}}{a}{l}{\oe}{\textbabygamma}{m}{au}{\textsubring{l}}} \dicPos{n}[2] \dicFlx{(‑s, ‑)}[5] \dicDirectTranslationCS{zákon příčinnosti\,/\addthin kauzality}
\dicEntry[orsakarsetning] \dicTerm{or·sakar··set·ning} \dicIPA{{\textopeno}{\textsubring{r}}{s}{a}{\r{g}}{a}{\textsubring{r}}{s}{\textepsilon}{h}{\textsubring{d}}{n}{i}{\ng}{\r{g}}} \dicPos{f}[4] \dicFlx{(‑ar, ‑ar)}[5] \dicFieldCat{jaz.} \dicDirectTranslationCS{vedlejší věta příslovečná příčinná}
\dicEntry[orsakarsögn] \dicTerm{or·sakar··|sögn} \dicIPA{{\textopeno}{\textsubring{r}}{s}{a}{\r{g}}{a}{\textsubring{r}}{s}{\oe}{\r{g}}{\textsubring{n}}} \dicPos{f}[7] \dicFlx{(‑sagnar, ‑sagnir)}[16] \dicFieldCat{jaz.} \dicDirectTranslationCS{kau\-za\-ti\-vum, kauzativní sloveso}
\dicEntry[orsakartenging] \dicTerm{or·sakar··teng·ing} \dicIPA{{\textopeno}{\textsubring{r}}{s}{a}{\r{g}}{a}{\textsubring{r}}{t\smash{\textsuperscript{h}}}{ei}{\textltailn}{\r{\textObardotlessj}}{i}{\ng}{\r{g}}} \dicPos{f}[4] \dicFlx{(‑ar, ‑ar)}[5] \dicFieldCat{jaz.} \dicDirectTranslationCS{příčinná spojka}
\dicEntry[orsakasamband] \dicTerm{or·saka··sam·band} \dicIPA{{\textopeno}{\textsubring{r}}{s}{a}{\r{g}}{a}{s}{a}{m}{\textsubring{b}}{a}{n}{\textsubring{d}}} \dicPos{n}[2] \dicFlx{(‑s)}[2] \dicDirectTranslationCS{příčinná souvislost}
\dicEntry[orsakatengsl] \dicTerm{or·saka··tengsl} \dicIPA{{\textopeno}{\textsubring{r}}{s}{a}{\r{g}}{a}{t\smash{\textsuperscript{h}}}{\textepsilon}{\ng}{s}{\textsubring{d}}{\textsubring{l}}} \dicPos{n}[2] \dicFlx{pl}[1] \dicDirectTranslationCS{příčinná spojitost}
\dicEntry[orsök] \dicTerm{or··|sök} \dicsymFrequent\  \dicIPA{{\textopeno}{\textsubring{r}}{s}{\oe}{\r{g}}} \dicPos{f}[7] \dicFlx{(‑sakar, ‑sakir)}[16] \dicDirectTranslationCS{příčina} \dicExampleIS{orsök og afleiðing} \dicExampleCS{příčina a~důsledek}
\dicEntry[ort] \dicTerm{ort} \dicIPA{{\textopeno}{\textsubring{r}}{\textsubring{d}}} \dicPos{v} \dicFlx{supin} \dicLink{yrkja\smash{\textsuperscript{2}}}
\dicEntry[orti] \dicTerm{orti} \dicIPA{{\textopeno}{\textsubring{r}}{\textsubring{d}}{\textsci}} \dicPos{v} \dicFlx{ind\,/\addthin con pf sg 1 pers} \dicLink{yrkja\smash{\textsuperscript{2}}}
\dicEntry[ortum] \dicTerm{ortum} \dicIPA{{\textopeno}{\textsubring{r}}{\textsubring{d}}{\textscy}{\textsubring{m}}} \dicPos{v} \dicFlx{ind pf pl 1 pers} \dicLink{yrkja\smash{\textsuperscript{2}}}
\dicEntry[orusta] \dicTerm{or··ust|a}\dicTerm{, orrusta} \dicIPA{{\textopeno}\-{\textlengthmark}\-{r}\-{\textscy}\-{s}\-{\textsubring{d}}\-{a}\-} \dicPos{f}[1] \dicFlx{(‑u, ‑ur)}[19] \dicDirectTranslationCS{bitva} \dicExampleIS{eftir harða og langa orustu} \dicExampleCS{po těžké a~dlouhé bitvě}
\dicEntry[orustuflugvél] \dicTerm{orustu··flug·vél} \dicIPA{{\textopeno}{\textlengthmark}{r}{\textscy}{s}{\textsubring{d}}{\textscy}{f}{l}{\textscy}{\textbabygamma}{v}{j}{\textepsilon}{\textsubring{l}}} \dicPos{f}[4] \dicFlx{(‑ar, ‑ar)}[1] \dicFieldCat{let.} \dicDirectTranslationCS{stíhačka, bitevní letoun}
\dicEntry[orustuskip] \dicTerm{orustu··skip} \dicIPA{{\textopeno}{\textlengthmark}{r}{\textscy}{s}{\textsubring{d}}{\textscy}{s}{\r{\textObardotlessj}}{\textsci}{\textsubring{b}}} \dicPos{n}[2] \dicFlx{(‑s, ‑)}[5] \dicFieldCat{voj.} \dicDirectTranslationCS{bitevní loď}
\dicEntry[oss] \dicTerm{oss} \dicIPA{{\textopeno}{s}{\textlengthmark}} \dicPos{pron}[17] \dicFlx{pers} \dicFlx{dat acc pl} \dicLangCat{kniž.} \dicDirectTranslationCS{nám, nás} \dicExampleIS{frá oss} \dicExampleCS{od nás} \dicLink{vér}
\dicEntry[ostakaka] \dicTerm{osta··|kaka} \dicIPA{{\textopeno}{s}{\textsubring{d}}{a}{k\smash{\textsuperscript{h}}}{a}{\r{g}}{a}} \dicPos{f}[1] \dicFlx{(‑köku, ‑kökur)}[20] \dicFieldCat{kulin.} \dicDirectTranslationCS{tvarohový koláč}
\dicEntry[ostaskeri] \dicTerm{osta··sker|i} \dicIPA{{\textopeno}{s}{\textsubring{d}}{a}{s}{\r{\textObardotlessj}}{\textepsilon}{r}{\textsci}} \dicPos{m}[1] \dicFlx{(‑a, ‑ar)}[1] \dicDirectTranslationCS{kráječ na sýr}
\dicEntry[ostborgari] \dicTerm{ost··borg·ar|i} \dicIPA{{\textopeno}{s}{\textsubring{d}}{\textsubring{b}}{\textopeno}{r}{\r{g}}{a}{r}{\textsci}} \dicPos{m}[1] \dicFlx{(‑a, ‑ar)}[13] \dicFieldCat{kulin.} \dicDirectTranslationCS{cheeseburger}
\dicEntry[ostra] \textls[15]{\dicTerm{ostr|a} \dicIPA{{\textopeno}{s}{\textsubring{d}}{r}{a}} \dicPos{f}[1] \dicFlx{(‑u, ‑ur)}[7] \dicFieldCat{zool.} \dicDirectTranslationCS{ústřice, ústřice jedlá}} \textit{(l.~{\textLA{Ostrea edulis}})}
\dicEntry[ostur] \dicTerm{ost|ur} \dicsymFrequent\  \dicIPA{{\textopeno}{s}{\textsubring{d}}{\textscy}{\textsubring{r}}} \dicPos{m}[6] \dicFlx{(‑s, ‑ar)}[4] \dicDirectTranslationCS{sýr} \dicExampleIS{þroskaður ostur} \dicExampleCS{zralý sýr}
\dicEntry[ota] \dicTerm{ot|a} \dicIPA{{\textopeno}{\textlengthmark}{\textsubring{d}}{a}} \dicPos{v}[1] \dicFlx{(‑aði)}[1] \dicFlx{dat} \dicDirectTranslationCS{(na)mířit, (za)cílit};  \dicPhraseIS{ota e‑u að e‑m} \dicDirectTranslationCS{mířit (čím) na (koho), ohrožovat (koho čím)} \dicExampleIS{ota hnífi að honum} \dicExampleCS{ohrožovat ho nožem};  \dicIdiom{ota}[fram]{ \dicPhraseIS{ota sér fram}} \dicLangCat{přen.} \dicDirectTranslationCS{cpát se dopředu, vnutit se, vnucovat se};  \dicIdiom{ota}{ \dicPhraseIS{ota sínum tota}} \dicLangCat{přen.} \dicDirectTranslationCS{hrabat si na svém písečku, prosazovat své zájmy}
\dicEntry[otur] \dicTerm{ot|ur} \dicIPA{{\textopeno}{\textlengthmark}{\textsubring{d}}{\textscy}{\textsubring{r}}} \dicPos{m}[5] \dicFlx{(‑urs, ‑rar)}[1] \dicFieldCat{zool.} \dicDirectTranslationCS{vydra, vydra říční} \textit{(l.~{\textLA{Lutra lutra}})}  \dicsymPhoto\ 
\dicFigure{114148.jpg}{Otur}{Otur - Hykel Michal, Biolib, Copyright/CC-BY-SA}
\dicEntry[oxíð] \dicTerm{oxíð} \dicIPA{{\textopeno}{\r{g}}{s}{i}{\texttheta}} \dicPos{n}[2] \dicFlx{(‑s, ‑)}[5] \dicFieldCat{chem.} \dicDirectTranslationCS{oxid, kysličník}
\dicEntry[oxun] \dicTerm{ox|un} \dicIPA{{\textopeno}{\r{g}}{s}{\textscy}{\textsubring{n}}} \dicPos{f}[7] \dicFlx{(‑unar)}[9] \dicFieldCat{chem.} \dicDirectTranslationCS{oxidace, okysličování}

\else
\makeatletter
\@for\next:=\alphabetF\do{\input{letters/\next.tex}}
\makeatother 
\fi
\restoregeometry
\pagestyle{plain}

\fi

% ================================================================ PHONETICS =

\ifinputphon

\cleardoublepage

\chapter{Seznam islandských fonémů}              \label{sec:phon_phonems}

\begin{table}[h]
\begin{tabular}{lllll} \toprule

\textbf{foném} & \textbf{vysvětlivka} & \textbf{příklad} & \textbf{fonetický přepis} & \textbf{překlad} \\
\midrule
{\textipa{[{a}]}} & otevřená přední nezaokrouhlená samohl. & rata & {\textipa{[{r}{a}{\textlengthmark}{\textsubring{d}}{a}]}} & trefit \\ 
{\textipa{[{\textepsilon}]}} & polootevřená přední nezaokrouhlená samohl. & meta & {\textipa{[{m}{\textepsilon}{\textlengthmark}{\textsubring{d}}{a}]}} & ocenit \\ 
{\textipa{[{e}]}} & polozavřená přední nezaokrouhlená samohl. & heiður & {\textipa{[{h}{ei}{\textlengthmark}{ð}{\textscy}{\textsubring{r}}]}} & čest \\
{\textipa{[{\textsci}]}} & téměř zavřená téměř přední nezaokrouhlená samohl. & vinna & {\textipa{[{v}{\textsci}{n}{\textlengthmark}{a}]}} & pracovat \\ 
{\textipa{[{i}]}} & zavřená přední nezaokrouhlená samohl. & fínn & {\textipa{[{f}{i}{\textsubring{d}}{\textsubring{n}}]}} & jemný \\ 
{\textipa{[{\textopeno}]}} & polootevřená zadní zaokrouhlená samohl. & lofa & {\textipa{[{l}{\textopeno}{\textlengthmark}{v}{a}]}} & slibovat \\ 
{\textipa{[{o}]}} & polozavřená zadní zaokrouhlená samohl. & bóndi & {\textipa{[{\textsubring{b}}{ou}{n}{\textsubring{d}}{\textsci}]}} & farmář \\ 
{\textipa{[{\textscy}]}} & téměř zavřená téměř přední zaokrouhlená samohl. & hundur & {\textipa{[{h}{\textscy}{n}{\textsubring{d}}{\textscy}{\textsubring{r}}]}} & pes \\ 
{\textipa{[{u}]}} & zavřená zadní zaokrouhlená samohl. & súkkulaði & {\textipa{[{s}{u}{h}{\r{g}}{\textscy}{l}{a}{ð}{\textsci}]}} & čokoláda \\ 
{\textipa{[{\oe}]}} & téměř otevřená přední nezaokrouhlená samohl. & lönd & {\textipa{[{l}{\oe}{n}{\textsubring{d}}]}} & země \\ 
{\textipa{[{m}]}} & bilabiální nazála & mýri & {\textipa{[{m}{i}{\textlengthmark}{r}{\textsci}]}} & mokřina \\ 
{\textipa{[{\textsubring{m}}]}} & neznělá bilabiální nazála & þyrmt & {\textipa{[{\texttheta}{\textsci}{\textsubring{r}}{\textsubring{m}}{\textsubring{d}}]}} & {\textit{$\shortrightarrow$ þyrma}} \\  
{\textipa{[{n}]}} & alveolární nazála & nemandi & {\textipa{[{n}{\textepsilon}{\textlengthmark}{m}{a}{n}{\textsubring{d}}{\textsci}]}} & žák \\ 
{\textipa{[{\textsubring{n}}]}} & neznělá alveolární nazála & mennt & {\textipa{[{m}{\textepsilon}{\textsubring{n}}{\textsubring{d}}]}} & vzdělání \\ 
{\textipa{[{\textltailn}]}} & palatální nazála & angi & {\textipa{[{au}{\textltailn}{\r{\textObardotlessj}}{\textsci}]}} & drobek \\ 
{\textipa{[{\r{\textltailn}}]}} & neznělá palatální nazála & banki & {\textipa{[{\textsubring{b}}{au}{\r{\textltailn}}{\r{\textObardotlessj}}{\textsci}]}} & banka \\
{\textipa{[{\ng}]}} & velární nazála & skyggndur & {\textipa{[{s}{\r{\textObardotlessj}}{\textsci}{\ng}{\textsubring{d}}{\textscy}{\textsubring{r}}]}} & {\textit{$\shortrightarrow$ skyggna}} \\ 
{\textipa{[{\r{\ng}}]}} & neznělá velární nazála & skyggnt & {\textipa{[{s}{\r{\textObardotlessj}}{\textsci}{\r{\ng}}{\textsubring{d}}]}} & {\textit{$\shortrightarrow$ skyggna}} \\ 
{\textipa{[{\textsubring{b}}]}} & neznělá bilabiální ploziva & bara & {\textipa{[{\textsubring{b}}{a}{\textlengthmark}{r}{a}]}} & jenom \\ 
{\textipa{[{p\textsuperscript{h}}]}} & neznělá bilabiální ploziva s přídechem & prestur & {\textipa{[{p\textsuperscript{h}}{r}{\textepsilon}{s}{\textsubring{d}}{\textscy}{\textsubring{r}}]}} & kněz \\ 
{\textipa{[{\textsubring{d}}]}} & neznělá alveolární ploziva & dæmi & {\textipa{[{\textsubring{d}}{a}{i}{\textlengthmark}{m}{\textsci}]}} & příklad \\ 
{\textipa{[{t\textsuperscript{h}}]}} & neznělá alveolární ploziva s přídechem & tómur & {\textipa{[{t\textsuperscript{h}}{ou}{\textlengthmark}{m}{\textscy}{\textsubring{r}}]}} & prázdný \\ 
{\textipa{[{\r{\textObardotlessj}}]}} & neznělá palatální ploziva & gefa & {\textipa{[{\r{\textObardotlessj}}{\textepsilon}{\textlengthmark}{v}{a}]}} & dát \\ 
{\textipa{[{c\textsuperscript{h}}]}} & neznělá palatální ploziva s přídechem & kenna & {\textipa{[{c\textsuperscript{h}}{\textepsilon}{n}{\textlengthmark}{a}]}} & učit \\
{\textipa{[{\r{g}}]}} & neznělá velární ploziva & gata & {\textipa{[{\r{g}}{a}{\textlengthmark}{\textsubring{d}}{a}]}} & ulice \\ 
{\textipa{[{k\textsuperscript{h}}]}} & neznělá velární ploziva s přídechem & koma & {\textipa{[{k\textsuperscript{h}}{\textopeno}{\textlengthmark}{m}{a}]}} & přijít \\
{\textipa{[{v}]}} & znělá labiodentální frikativa & hafa & {\textipa{[{h}{a}{\textlengthmark}{v}{a}]}} & mít \\ 
{\textipa{[{f}]}} & neznělá labiodentální frikativa & finna & {\textipa{[{f}{\textsci}{n}{\textlengthmark}{a}]}} & najít \\ 
{\textipa{[{ð}]}} & znělá dentální frikativa & borða & {\textipa{[{\textsubring{b}}{\textopeno}{r}{ð}{a}]}} & jíst \\ 
{\textipa{[{\texttheta}]}} & neznělá dentální frikativa & þola & {\textipa{[{\texttheta}{\textopeno}{\textlengthmark}{l}{a}]}} & snést \\
{\textipa{[{s}]}} & neznělá alveolární frikativa & systir & {\textipa{[{s}{\textsci}{s}{\textsubring{d}}{\textsci}{\textsubring{r}}]}} & sestra \\ 
{\textipa{[{\c{c}}]}} & neznělá palatální frikativa & hjóla & {\textipa{[{\c{c}}{ou}{\textlengthmark}{l}{a}]}} & jet na kole \\ 
{\textipa{[{\textbabygamma}]}} & znělá velární frikativa & sagði & {\textipa{[{s}{a}{\textbabygamma}{ð}{\textsci}]}} & řekl \\ 
{\textipa{[{x}]}} & neznělá velární frikativa & lag & {\textipa{[{l}{a}{\textlengthmark}{x}]}} & píseň \\
{\textipa{[{h}]}} & neznělá glotální frikativa & heima & {\textipa{[{h}{ei}{\textlengthmark}{m}{a}]}} & doma \\ 
{\textipa{[{l}]}} & alveolární laterální aproximanta & loforð & {\textipa{[{l}{\textopeno}{\textlengthmark}{v}{\textopeno}{r}{\texttheta}]}} & slib \\ 
{\textipa{[{\textsubring{l}}]}} & neznělá alveolární laterální aproximanta & hlutur & {\textipa{[{\textsubring{l}}{\textscy}{\textlengthmark}{\textsubring{d}}{\textscy}{\textsubring{r}}]}} & věc \\  
{\textipa{[{j}]}} & palatální aproximanta & jól & {\textipa{[{j}{ou}{\textlengthmark}{\textsubring{l}}]}} & Vánoce \\ 
{\textipa{[{r}]}} & alveolární vibranta & reynsla & {\textipa{[{r}{ei}{n}{s}{\textsubring{d}}{l}{a}]}} & zkušenost \\ 
{\textipa{[{\textsubring{r}}]}} & neznělá alveolární vibranta & hraun & {\textipa{[{\textsubring{r}}{\oe i}{\textlengthmark}{\textsubring{n}}]}} & láva \\ 
\bottomrule
\end{tabular}
\end{table}

\clearpage

%\section{Výslovnost islandských samohlásek}      \label{sec:phon_vowels}
%\LTXtable{\columnwidth}{phon/vowels}
%\clearpage

%\section{Výslovnost islandských souhlásek}       \label{sec:phon_consonants}
%\LTXtable{\columnwidth}{phon/consonants}
%\clearpage
\fi

% =============================================================== MORPHOLOGY =

\ifinputmorpho

\cleardoublepage

\chapter{Morfologický klíč}                      \label{sec:morpho}

\section{Podstatné jméno, rod mužský}            \label{sec:morpho_m}
{\small\LTXfw{morpho/m}}
\clearpage

\section{Podstatné jméno, rod ženský}            \label{sec:morpho_f}
{\small\LTXfw{morpho/f}}
\clearpage

\section{Podstatné jméno, rod střední}           \label{sec:morpho_n}
{\small\LTXfw{morpho/n}}
\clearpage

\section{Přídavné jméno}                         \label{sec:morpho_adj}
{\small\LTXfw{morpho/adj}}
\clearpage

\section{Zájmeno}                                \label{sec:morpho_pron}
{\small\LTXfw{morpho/pron}}
 \begin{center}
\begin{minipage}[t]{.45\textwidth}
\begin{tabular}{l>{\footnotesize\itshape}lll}
 {\textbf{\textit{třída}}} & {\textit{pád}}   & \textit{sg} & \textit{pl}  \\ 
 \hline
\multirow{3}{*}{{{\textbf{pron} \Large{\textbf{17}}}}}  &  nom & ég & við   \\*
 & acc &  mig  & okkur  \\*
 & dat & mér & okkur   \\*
 & gen & mín  & okkar  \\
\hline

\multirow{3}{*}{{{\textbf{pron} \Large{\textbf{17}}}}}  &  nom & hann & þeir   \\*
 & acc &  hann  & þá  \\*
 & dat & honum & þeim   \\*
 & gen & hans  & þeirra  \\
\hline

\multirow{3}{*}{{{\textbf{pron} \Large{\textbf{17}}}}}  &  nom & hún & þær   \\*
 & acc &  hana  & þær  \\*
 & dat & henni & þeim   \\*
 & gen & hennar  & þeirra  \\
\hline

\multirow{3}{*}{{{\textbf{pron} \Large{\textbf{17}}}}}  &  nom &  & vér   \\*
 & acc &    & oss  \\*
 & dat &  & oss   \\*
 & gen &   & vor  \\
\hline

\multirow{3}{*}{{{\textbf{pron} \Large{\textbf{17}}}}}  &  nom &  & þér   \\*
 & acc &    & yður  \\*
 & dat &  & yður   \\*
 & gen &   & yðar  \\
\hline

\multirow{3}{*}{{{\textbf{pron} \Large{\textbf{17}}}}}  &  nom & þú & þið   \\*
 & acc &  þig  & ykkur  \\*
 & dat & þér & ykkur   \\*
 & gen & þín  & ykkar  \\
\hline

\multirow{3}{*}{{{\textbf{pron} \Large{\textbf{17}}}}}  &  nom & það & þau   \\*
 & acc &  það  & þau  \\*
 & dat & því & þeim   \\*
 & gen & þess  & þeirra  \\
\hline
\end{tabular}
\end{minipage}\hfil
\begin{minipage}[t]{.45\textwidth}
\begin{tabular}{l>{\footnotesize\itshape}lll}
\hline
\multirow{3}{*}{{{\textbf{pron} \Large{\textbf{18}}}}}  &  nom &  &    \\*
 & acc &  sig  &   \\*
 & dat & sér &    \\*
 & gen & sín  &   \\
 \hline
\end{tabular}
\end{minipage}
\end{center}

\clearpage

\section{Číslovka}                               \label{sec:morpho_num}
 {\small\LTXfw{morpho/num}}
\clearpage

\begin{landscape}
\section{Sloveso}                                \label{sec:morpho_v}
{\small\LTXfw{morpho/v}}
\end{landscape}

\clearpage

\fi

% ============================================================== BACK MATTER =

\ifmakebackmatter

% ------------------------------------------------ Photographs authors index -

\cleardoublepage
%\pagestyle{empty}
%\pagestyle{indexstyle}
%\phantomsection
%\addcontentsline{toc}{chapter}
%  {Seznam autorů fotografií a ilustrací}        
\label{sec:photo}
\printindex[figures]

% ------------------------------------------------------------- Bibliography -
\cleardoublepage
\chapter{Bibliografie}
%\twocolumn
%\phantomsection
%\addcontentsline{toc}{chapter}{Bibliografie}
\begin{multicols}{2}
\nocite{*}\printbibliography[heading=none]
\end{multicols}

% ------------------------------------------------------------------ Apendix -
\cleardoublepage
\onecolumn
\ifPDF
\else
\chapter{Apendix: Islandská gramatika}
\clearpage
\covergeometry  
\makeatletter
\includepdf[
  pages={1,2,3,4,5,6,7,8,9,10,11,12,13,14,15,16,17,18,19,20,21,22,23,24,25,26,27,28,29,30,31,32,33,34,35,36,37,38,39,40,41,42,43,44,45,46,47,48,49,50,51,52,53,54,55,56,57,58,59,60,61,62,63,64,65,66,67,68,69,70,71,72,73},
  turn=false,
  scale=0.79,
  keepaspectratio,
  offset={\dimexpr(\paperwidth-\csname Gm@layoutwidth\endcsname)/2-\csname Gm@layouthoffset\endcsname\relax}
  {\dimexpr(\paperheight-\csname  Gm@layoutheight\endcsname)/2-\csname Gm@layoutvoffset\endcsname\relax}
   ]{gram.pdf}
\makeatother
\restoregeometry  
\fi

% ------------------------------------------------------------------ Barcode -
\pagenumbering{gobble}
\ifPDF
\else
\cleardoublepage\null\clearpage % clear to verso
\vspace*\fill
\EANisbn[SC3]
\fi
\fi

% --------------------------------------------------------------- Back cover -

\ifmakecovers
  \cleardoublepage\null\clearpage % inner side of the cover
\covergeometry  
  \makecoverwith\backcoverimages
\restoregeometry  
\fi

\end{document}

